\documentclass[10pt,a4paper]{article}

% ----------------- Paquetes -------------------%
\usepackage[T1]{fontenc}
\usepackage[utf8]{inputenc}
\usepackage[a4paper,margin=2.5cm]{geometry}
\usepackage{mathptmx}             
\usepackage{changepage} 
\usepackage{setspace}
\usepackage[spanish]{babel}
\usepackage[labelfont=bf,labelsep=space]{caption}
\usepackage{amssymb}
\usepackage{amsmath}
\usepackage{ebgaramond}
\usepackage{placeins}
\usepackage{etoolbox}
\usepackage{setspace}
\usepackage{parskip} 
\usepackage{tcolorbox}
\usepackage{needspace}
\usepackage{xparse}
\usepackage{ocgx2}
\usepackage{hyperref}
\usepackage{hypcap}


%%%%%%%%%%%%%%%%%%%%%%%%%%%%%%%%%%%%%%%%%%%%%%%%%%%%%%%%%%%%%%%%
% --- Fija primero tus valores globales "buenos" ---
\setstretch{1.2}                 % Interlineado global
\setlength{\parskip}{0.7\baselineskip}
\setlength{\parindent}{0pt}
\newlength{\GlobalParskip}
\setlength{\GlobalParskip}{\parskip}

\newcounter{eqnum}[subsection]
\renewcommand{\theeqnum}{\arabic{eqnum}}
\newcounter{alphaitem}
\newcounter{numitem}

%% ------------------- Recuadros----------------------%%
\newcommand{\puntosclave}[1]{%
  \begin{tcolorbox}[
    colframe=black,
    title={Puntos clave},
    before upper={\setlength{\parindent}{0.5cm}\setlength{\parskip}{0.5\baselineskip}}
  ]
  #1
  \end{tcolorbox}
}

\newcommand{\modificaciones}[1]{
  \begin{tcolorbox}[
    colback=white,
    colframe=red,
    title={Modificaciones a realizar},
    before upper={\setlength{\parindent}{0.5cm}\setlength{\parskip}{0.5\baselineskip}}
  ]
  #1
  \end{tcolorbox}
}

%% ------------------- Preguntas TEST----------------------%%
% Opciones: radio (single-choice)
\NewDocumentCommand{\OptRadio}{m m m}{%
  \noindent\CheckBox[radio,name=#1,value=#2,width=1.2ex,height=1.2ex]{}%
  \;\textbf{#2.}~#3\par
}

% Opciones: checkbox (multi-choice)
\NewDocumentCommand{\OptCheck}{m m}{%
  \noindent\CheckBox[name=#1,width=1.2ex,height=1.2ex,bordercolor={0 0 0}]{}%
  \;#2\par
}

% Pregunta single-choice (radio)
% Args: 1=id, 2=enunciado, 3=opciones, 4=solución+justificación, 5=imagen (o {}), 6=origen
\NewDocumentCommand{\PreguntaTestSingle}{m m m m m m}{%
  \par\medskip
  \noindent\textbf{#1.}~#2\par\smallskip

    % Imagen (si no es {})
    \IfBlankTF{#5}{}{%
      \begin{center}
        \includegraphics[width=0.75\linewidth]{#5}
      \end{center}
    }

    \begin{Form}
      #3
    \end{Form}

    \par\smallskip
    \switchocg{sol-#1}{\fbox{Mostrar / ocultar solución}}%
    \hfill{\footnotesize #6}\par\smallskip

    \begin{ocg}{Solución}{sol-#1}{0}
      \noindent #4
    \end{ocg}
  \par\medskip
}

% Pregunta multi-choice (checkboxes)
\NewDocumentCommand{\PreguntaTestMulti}{m m m m m m}{%
  \par\medskip
  \noindent\textbf{#1.}~#2\par\smallskip

    % Imagen (si no es {})
    \IfBlankTF{#5}{}{%
      \begin{center}
        \includegraphics[width=0.75\linewidth]{#5}
      \end{center}
    }
    \begin{Form}
      #3
    \end{Form}

    \par\smallskip
    \switchocg{sol-#1}{\fbox{Mostrar / ocultar solución}}%
    \hfill{\footnotesize #6}\par\smallskip

    \begin{ocg}{Solución}{sol-#1}{0}
      \noindent #4
    \end{ocg}
  \par\medskip
}

%% ------------------- Listas----------------------%%
    % --- Lista alfabética ---
    \newenvironment{la}
      {\begin{list}{\alph{alphaitem})}{%
          \usecounter{alphaitem}%
          \setlength{\leftmargin}{1cm}%
          \setlength{\parskip}{3pt}% 
          \setlength{\itemsep}{0pt}% ← espacio entre ítems
          \setlength{\parsep}{0pt}%  ← párrafos dentro de un ítem
      }}
      {\end{list}}
      
    % --- Lista numérica ---
    \newenvironment{lnum}
      {\begin{list}{\arabic{numitem}.}{%
          \usecounter{numitem}%
          \setlength{\leftmargin}{1cm}%
          \setlength{\parskip}{3pt}% 
          \setlength{\itemsep}{0pt}% ← espacio entre ítems
          \setlength{\parsep}{0pt}%  ← párrafos dentro de un ítem
      }}
      {\end{list}}
      
% ------------------- Imágenes----------------------%
    \graphicspath{{Img/}}% Rutas por defecto 
    % --- Imagen adaptativa ---
    % Registros de dimensión definidos una sola vez
    \newdimen\imgcap
    \newdimen\imgmin
    \newdimen\imgmax
    \newdimen\imgremain
    \newdimen\imgh
    \newdimen\imgneed
    
    \newcommand{\imagenfit}[3]{%
      \addvspace{10pt}% separación antes
      \par\begingroup
        % Parámetros de composición (asignaciones, no \newdimen)
        \imgcap=1\baselineskip            % alto estimado de título+fuente
        \imgmin=0.15\textheight
        \imgmax=0.15\textheight
        \imgremain=\pagegoal
        \ifdim\pagegoal=\maxdimen \imgremain=0pt\fi
        \advance\imgremain by -\pagetotal
        \advance\imgremain by -\imgcap
        % Decisión de altura destino
        \ifdim\imgremain<\imgmin
          \newpage
          \imgh=0.20\textheight
        \else
          \imgh=\imgremain
          \ifdim\imgh>\imgmax \imgh=\imgmax \fi
        \fi
        % Reservar espacio para evitar cortes del bloque completo
        \imgneed=\imgh \advance\imgneed by \imgcap
        \Needspace{\imgneed}
        % Cuerpo
        \noindent\begin{minipage}{\textwidth}
          \centering
          \captionof{figure}{\textemdash\ #2}\par\vspace{0.3em}% título arriba
          \includegraphics[width=\linewidth,height=\imgh,keepaspectratio]{\detokenize{#1}}\par
          \vspace{0.3em}\small\textit{Fuente: }#3% fuente abajo
        \end{minipage}%
      \endgroup
      \par\addvspace{10pt}%
    }

%------------ Bloque de RECUADRO ESPECIAL -------------%
    \newenvironment{specialblock}[1]{%
      \begin{quote} % crea sangría a izquierda y derecha
      \small % texto más pequeño
      \noindent\textbf{#1}\\[-0.3em] % título en negrita
      \rule{\linewidth}{0.4pt}\\[0.6em]}{
      \end{quote}}
      
%-------------------------- Portada -----------------------%
\newcommand{\oldtitle}[1]{\def\OldTitle{#1}}
\newcommand{\changerationale}[1]{\def\ChangeWhy{#1}}

\newcommand{\changebox}{%
\begin{tcolorbox}[title={Historial de cambios del tema}]
\textbf{Título anterior:} \OldTitle\par
\textbf{Motivo del cambio:} \ChangeWhy
\end{tcolorbox}
}

%-------------------------- Author Font -----------------------%
    \newcommand{\authorfont}[1]{{\fontfamily{EBGaramond-TLF}\selectfont #1}}

%-------------------------- Anexos -----------------------%
    \makeatletter
    \newcounter{anexo}
    \renewcommand{\theanexo}{\Roman{anexo}}
    \newcommand{\anexo}[1]{%
      \clearpage
      \phantomsection
      \refstepcounter{anexo}%
      \noindent\textbf{Anexo \theanexo.- }#1\par\medskip
      \addcontentsline{stc}{section}{Anexo \theanexo.- #1}%
    }
    \makeatother

    %-------------------------- Lista de figuras (personal) -----------------------%
\makeatletter
\newcommand{\MyListOfFigures}{%
  \clearpage
  \phantomsection
  {\centering\bfseries\Large Índice de figuras\par}%
  \medskip
  % Entrada en nuestro índice de contenido (stc)
  \addcontentsline{stc}{section}{Índice de figuras}%
  % Recuperación del listado (sin cabecera propia)
  \begingroup
    \parindent=0pt\parskip=2pt
    \@starttoc{lof}%
  \endgroup
}
\makeatother

%-------------------------- Bibliografía -----------------------%
\makeatletter
% Contador para las referencias
\newcounter{myref}
% Cabecera de la bibliografía, entra en el índice 'stc'
\newcommand{\MyBibliography}{%
  \clearpage
  \phantomsection
  {\centering\bfseries\Large Bibliografía y referencias\par}%
  \medskip
  \addcontentsline{stc}{section}{Bibliografía y referencias}%
  \setcounter{myref}{0}%
}
\newcommand{\MyBibItem}[4]{%
  \refstepcounter{myref}%
  \noindent[\themyref]\ #1.\ \emph{#2}. #3. #4\par
}
\makeatother

%--------- Establecer cuatro niveles de índice --------%
    \setcounter{tocdepth}{4}   
    \setcounter{secnumdepth}{4}% numera hasta \paragraph

%%%%%%%%%%%%%%%%%%%%%%%%%%%%%%%%

\makeatletter

    %--------- Interlineado Global --------%
    \edef\GlobalBaseLineStretch{\baselinestretch}
    
    %--------- Espaciado de seccinones --------%
    \renewcommand\section{\@startsection{section}
      {1}{\z@}
      {2.5ex \@plus 1ex \@minus .2ex} % espacio antes
      {1.5ex \@plus .2ex}             % espacio después
      {\normalfont\Large\bfseries}    % formato del título
    }
    \renewcommand\subsection{\@startsection{subsection}{2}{\z@}
      {3ex \@plus 0.8ex \@minus .2ex}
      {1.5ex \@plus .2ex}
      {\normalfont\large\bfseries}
    }
    \renewcommand\subsubsection{\@startsection{subsubsection}{3}{\z@}
      {3ex \@plus 0.5ex \@minus .2ex}
      {1ex \@plus .2ex}
      {\normalfont\normalsize\bfseries}
    }
    \renewcommand\paragraph{\@startsection{paragraph}{4}{\z@}%
    {-2ex \@plus -0.5ex \@minus -.2ex}% espacio antes
    {0.8ex \@plus .2ex}% espacio después
    {\normalfont\normalsize\itshape}}% ← cursiva en lugar de negrita

    %--------- Ecuaciones --------%   
    % Variables globales para almacenar títulos y notas
    \gdef\leq@current@section{}%
    \gdef\leq@current@subsection{}%
    \gdef\leq@last@printed@section{}%
    \gdef\leq@last@printed@subsection{}%
    
    % Comando para añadir nota (invisible en el cuerpo)
    \newcommand{\eqnote}[1]{%
      \addcontentsline{leq}{leqnote}{#1}%
    }
    
    % Comando para ecuaciones
    \newcommand{\eqblock}[2]{%
      % Verificar si necesitamos imprimir el título de sección
      \ifx\leq@current@section\leq@last@printed@section\else
        \addcontentsline{leq}{leqsection}{\leq@current@section}%
        \global\let\leq@last@printed@section\leq@current@section
        \gdef\leq@last@printed@subsection{}% Reset subsección al cambiar sección
      \fi
      % Verificar si necesitamos imprimir el título de subsección
      \ifx\leq@current@subsection\leq@last@printed@subsection\else
        \ifx\leq@current@subsection\@empty\else
          \addcontentsline{leq}{leqsubsection}{\leq@current@subsection}%
          \global\let\leq@last@printed@subsection\leq@current@subsection
        \fi
      \fi
      % Ecuación
      \refstepcounter{eqnum}%
      \begin{equation*}
        #1\tag{E.\theeqnum}
      \end{equation*}%
      \addcontentsline{leq}{leqt}{[E.\theeqnum] -- #2}%
      \addcontentsline{leq}{leqm}{#1}%
    }
    
    %%% ----- Índice de ecuaciones ----------%%%
    % Tamaño para []
    \newcommand{\leqIndexFont}{\normalsize}
    \newcommand{\leqTitleFont}{\tiny}
    \newcommand{\leqNoteFont}{\tiny}
    % Cabecera de sección 
    \newcommand*\l@leqsection[2]{%
      \addvspace{25pt}% Mayor separación antes de nueva sección
      \noindent\leqIndexFont
      \makebox[\linewidth]{\rule{.38\linewidth}{0.4pt}\ \ \textbf{  \MakeUppercase{#1}}\ \ \rule{.38\linewidth}{0.4pt}}%
      \par\vspace{-10pt}
    }
    % Cabecera de subsección 
    \newcommand*\l@leqsubsection[2]{%
      \par\addvspace{15pt}%
      \noindent\leqIndexFont #1
      \par\vspace{-7pt}
    }
    % Título de ecuación 
    \newcommand*\l@leqt[2]{%
    \par\addvspace{10pt}%
      \par\noindent\leqTitleFont #1\par
    }
    % Miniatura de ecuación
    \newcommand*\l@leqm[2]{%
      \vspace{0.3\baselineskip}%
      \par\scriptsize\noindent\hspace*{1.5em}\(#1\)\par%
    }
    % Nota de ecuación 
    \newcommand*\l@leqnote[2]{%
      \vspace{0.2\baselineskip}%
      \par\noindent\leqNoteFont\hspace*{1.5em}\textbf{Nota.--} #1\par\vspace{0.5\baselineskip}%
    }
    % Generación del índice
    \newcommand{\mylistofequations}{%
      \clearpage
      \phantomsection
      % Título visual de la lista (ajústalo a tu gusto):
      {\centering\bfseries\Large Índice de ecuaciones\par}
      % Entrada en nuestro índice 'stc':
      \addcontentsline{stc}{section}{Índice de ecuaciones}%
      % Llamada a tu lista real (usa el nombre exacto de tu macro):
      \@starttoc{leq}%
    }
    
    %--------- Captura de títulos de sección --------%
    \let\leq@original@section\section
    \let\leq@original@subsection\subsection
    % Flag para saber si estamos en una sección asterisco
    \newif\ifLeqStarSection
    \LeqStarSectionfalse
    
    % Redefinir \section CON CONTROL DE ASTERISCO
    \renewcommand{\section}{%
      \@ifstar{\leq@section@star}{\@dblarg\leq@section@nostar}%
    }
    \def\leq@section@star#1{%
      \global\LeqStarSectiontrue
      \gdef\leq@current@section{#1}%
      \gdef\leq@current@subsection{}%
      \leq@original@section*{#1}%
      \global\LeqStarSectionfalse
    }
    \def\leq@section@nostar[#1]#2{%
      \global\LeqStarSectionfalse
      \gdef\leq@current@section{#2}%
      \gdef\leq@current@subsection{}%
      \leq@original@section[#1]{#2}%
    }
    % Redefinir \subsection
    \renewcommand{\subsection}{%
      \@ifstar{\leq@subsection@star}{\@dblarg\leq@subsection@nostar}%
    }
    \def\leq@subsection@star#1{%
      \global\LeqStarSectiontrue
      \gdef\leq@current@subsection{#1}%
      \leq@original@subsection*{#1}%
      \global\LeqStarSectionfalse
    }
    \def\leq@subsection@nostar[#1]#2{%
      \global\LeqStarSectionfalse
      \gdef\leq@current@subsection{#2}%
      \leq@original@subsection[#1]{#2}%
    }

    % ----- Restauración de valores Globales de estilo ----------%
    % Flags
    \newif\ifInFootnote
    \InFootnotefalse
    \pretocmd{\@footnotetext}{\InFootnotetrue}{}{}
    \apptocmd{\@footnotetext}{\InFootnotefalse}{}{}
    % Profundidad de listas (soporta anidadas)
    \newcount\MyListDepth \MyListDepth=0
    \AtBeginEnvironment{la}{\advance\MyListDepth by 1}
    \AfterEndEnvironment{la}{\advance\MyListDepth by -1}
    \AtBeginEnvironment{lnum}{\advance\MyListDepth by 1}
    \AfterEndEnvironment{lnum}{\advance\MyListDepth by -1}
    \AtBeginEnvironment{itemize}{\advance\MyListDepth by 1}
    \AfterEndEnvironment{itemize}{\advance\MyListDepth by -1}
    \AtBeginEnvironment{enumerate}{\advance\MyListDepth by 1}
    \AfterEndEnvironment{enumerate}{\advance\MyListDepth by -1}
    % Restauración diferida: hazlo al INICIO del próximo párrafo real
    \newif\if@pendingRestore
    \@pendingRestorefalse
    \newcommand{\RequestRestore}{\global\@pendingRestoretrue}
    \everypar{%
      \if@pendingRestore
        \global\@pendingRestorefalse
        \ifInFootnote\else
          \ifnum\MyListDepth=0
            \setlength{\parskip}{\GlobalParskip}%
            \setstretch{\GlobalBaseLineStretch}%
          \fi
        \fi
      \fi
    }
    % Al final de bloques:
    \AfterEndEnvironment{equation}{\RequestRestore}
    \AfterEndEnvironment{equation*}{\RequestRestore}
    \AfterEndEnvironment{la}{\RequestRestore}
    \AfterEndEnvironment{lnum}{\RequestRestore}
    % Al final de macros:
    \apptocmd{\eqblock}{\RequestRestore}{}{}
    \apptocmd{\imagenfit}{\RequestRestore}{}{}
    
\makeatother

% ===================== ToC propio con 4 niveles (único bloque) =====================
\makeatletter

% --- Escritores: añadimos una línea al .stc en cada cabecera numerada ---
% Guardamos las versiones actuales para llamar al comportamiento normal
\let\MyTOC@orig@section\section
\let\MyTOC@orig@subsection\subsection
\let\MyTOC@orig@subsubsection\subsubsection
\let\MyTOC@orig@paragraph\paragraph

% SECTION
\renewcommand{\section}{%
  \@ifstar{\MyTOC@section@star}{\@dblarg\MyTOC@section@nostar}%
}
\def\MyTOC@section@star#1{%
  \MyTOC@orig@section*{#1}% sin entrada al .stc
}
\def\MyTOC@section@nostar[#1]#2{%
  \MyTOC@orig@section[{#1}]{#2}% comportamiento normal (escribe al .toc)
  % Escribimos también nuestra línea al .stc (texto con \numberline)
  \addcontentsline{stc}{section}{\protect\numberline{\thesection}#2}%
}

% SUBSECTION
\renewcommand{\subsection}{%
  \@ifstar{\MyTOC@subsection@star}{\@dblarg\MyTOC@subsection@nostar}%
}
\def\MyTOC@subsection@star#1{%
  \MyTOC@orig@subsection*{#1}%
}
\def\MyTOC@subsection@nostar[#1]#2{%
  \MyTOC@orig@subsection[{#1}]{#2}%
  \addcontentsline{stc}{subsection}{\protect\numberline{\thesubsection}#2}%
}

% SUBSUBSECTION
\renewcommand{\subsubsection}{%
  \@ifstar{\MyTOC@subsubsection@star}{\@dblarg\MyTOC@subsubsection@nostar}%
}
\def\MyTOC@subsubsection@star#1{%
  \MyTOC@orig@subsubsection*{#1}%
}
\def\MyTOC@subsubsection@nostar[#1]#2{%
  \MyTOC@orig@subsubsection[{#1}]{#2}%
  \addcontentsline{stc}{subsubsection}{\protect\numberline{\thesubsubsection}#2}%
}

% PARAGRAPH
\renewcommand{\paragraph}{%
  \@ifstar{\MyTOC@paragraph@star}{\@dblarg\MyTOC@paragraph@nostar}%
}
\def\MyTOC@paragraph@star#1{%
  \MyTOC@orig@paragraph*{#1}%
}
\def\MyTOC@paragraph@nostar[#1]#2{%
  \MyTOC@orig@paragraph[{#1}]{#2}%
  % Si no numeras los \paragraph, cambia \theparagraph por texto plano sin \numberline
  \addcontentsline{stc}{paragraph}{\protect\numberline{\theparagraph}#2}%
}

% --- Impresor del índice propio (lee el .stc) ---
\newcommand{\MyTOC}{%
  \par\bigskip
  {\centering\bfseries\Large Índice de contenido\par}%
  \medskip
  \begingroup
    \parindent=0pt\parskip=2pt
    \@starttoc{stc}%
  \endgroup
}

\makeatother
% ===================== fin ToC propio =====================


%%%%%%%%%%%%%%


% --- Datos del tema ---
\title{La Política Comercial (I)\\
\large Instrumentos y efectos. Barreras arancelarias y no arancelarias. Otros instrumentos tradicionales.}
\author{Marcelo Ortega}
\date{Abril 2026}
\newcommand{\TemaCode}{3B7}

\oldtitle{
    B.7. La política comercial (I). Instrumentos y efectos. Barreras arancelarias y no arancelarias. Otros instrumentos tradicionales.
}
\changerationale{
    Sin cambios.
}

\begin{document}


\maketitle
\changebox

\newpage

% --- Página de resumen y modificaciones ---
\puntosclave{ %% Escribir puntos clave Apuntes CECO
     
    }
\vspace{1cm}

\modificaciones{
    
    }

\newpage
\MyTOC
\newpage

\mylistofequations
\clearpage

\MyListOfFigures
\clearpage


\pagenumbering{arabic}
\setcounter{page}{1}


\section{Introducción}



\subsection{Contextualización}


\subsubsection{Relevancia}




\subsubsection{Problemática}



\subsection{Estructura de la exposición}



\clearpage
\section{Política Comercial: Marco conceptual}

Es bien sabido que la apertura comercial puede beneficiar a unos sectores de la sociedad y perjudicar a otros. Este hecho lleva a que las ganancias del comercio se repartan de manera desigual, de tal forma que existirán empresas, sindicatos o grupos de presión que intentarán que los poderes públicos intervengan para limitar las pérdidas (o maximizar las ganancias) originadas por el comercio internacional. De esta forma, la Política Comercial de un gobierno tendrá por objetivo influir sobre el volumen de comercio internacional. La política comercial puede instrumentarse de numerosas maneras mediante el uso de aranceles, cuotas, subvenciones, y multitud de otras barreras no arancelarias.\footnote{Como ya se ha resaltado, existen otros muchos instrumentos de política comercial que pueden aplicar los Estados, como son: las medidas antidumping, antisubvención y las salvaguardias. En principio no es necesario exponer ningún otro instrumento a parte de los aranceles, cuotas y subvenciones a la exportación. Sin embargo, el opositor debe estar familiarizado con estos otros instrumentos por si el tribunal le preguntara al respecto.
} Los efectos de los principales instrumentos de política comercial se pueden resumir de forma útil como sigue:

\begin{center}
\begin{tabular}{|p{0.15\linewidth}|p{0.15\linewidth}|p{0.15\linewidth}|p{0.15\linewidth}|p{0.15\linewidth}|}
\hline
 & \textbf{Arancel} & \textbf{Subsidio a la exportación} & \textbf{Cuota de importación} & \textbf{Restricción voluntaria a la exportación} \\
\hline
 Excedente del productor & Aumenta & Aumenta & Aumenta & Aumenta\\
\hline
Excedente del consumidor & Disminuye & Disminuye & Disminuye & Disminuye \\
\hline
Ingresos públicos & Aumentan & Disminuye (el Gasto Público aumenta) & Sin cambios (renta hacia tenedores de licencias) & Sin cambios (renta hacia los extranjeros) \\
\hline
Bienestar nacional total & Ambiguo (disminuye para país pequeño) & Disminuye & Ambiguo (disminuye para país pequeño) & Disminuye \\
\hline
\end{tabular}
\end{center}

Aunque existen multitud de clasificaciones para los instrumentos de Política Comercial, los presentaremos haciendo especial hincapié en su generalidad (¿cuántos países los utilizan?) y su desarrollo histórico (¿cuándo surgieron?). De esta forma, podemos formar dos grandes grupos de instrumentos: los convencionales y los no convencionales.

\subsection{Instrumentos convencionales}


\subsubsection{Aranceles}
Un arancel, la más simple de las políticas comerciales, es un impuesto a la importación de un bien. 

Los aranceles fijos son una cantidad preestablecida que se exige por cada unidad de bien importado (por ejemplo, tres dólares por barril de petróleo). Los aranceles ad valorem son impuestos exigidos como porcentaje del valor de los bienes importados (por ejemplo, el 25\% del valor de los camiones importados a los Estados Unidos; véase el recuadro más adelante en este capítulo). En ambos casos, el efecto del arancel es aumentar el coste de llevar los bienes a un país. Los aranceles son la forma más antigua de política comercial, y han sido utilizados tradicionalmente como una fuente de ingresos para los estados. 

Hasta la introducción del impuesto sobre la renta, por ejemplo, el gobierno de los Estados Unidos obtenía la mayor parte de sus ingresos en concepto de aranceles. Sin embargo, en general su verdadera finalidad ha sido no solo proporcionar ingresos, sino proteger sectores nacionales concretos. A principios del siglo XX, el Reino Unido utilizaba aranceles (las famosas Corn Laws) para proteger su agricultura de la competencia de las importaciones. A finales de ese siglo, tanto Alemania como los Estados Unidos protegían sus nuevos sectores industriales mediante la imposición de aranceles a la importación de bienes manufacturados. La importancia de los aranceles ha disminuido en los tiempos modernos, porque hoy día los estados generalmente prefieren proteger las industrias nacionales mediante diversas barreras no arancelarias, tales como cuotas de importación (limitaciones a la cantidad de importaciones) y restricciones a la exportación (sobre la cantidad de exportaciones, normalmente impuestas por el país exportador a solicitud del importador). No obstante, la comprensión de los efectos de un arancel es aún una base esencial para entender las demás políticas comerciales.\footnote{En el desarrollo de la teoría del comercio en los capítulos 3 a 8 adoptamos la perspectiva del equilibrio general. Es decir, éramos completamente conscientes de que lo que ocurre en una parte de la economía tiene repercusiones en otras. Sin embargo, en muchos casos (aunque no en todos) las políticas comerciales dirigidas a un sector pueden ser razonablemente bien entendidas sin entrar en los detalles de las repercusiones de esas políticas en el resto de la economía. Así pues, en su mayor parte la política comercial puede ser analizada en el marco del equilibrio parcial. Cuando los efectos sobre la economía en su conjunto sean cruciales volveremos a referimos al análisis del equilibrio general.
}


Desde el punto de vista de alguien que comercia con bienes, un arancel es como un coste de transporte. Si nuestro país impone un impuesto de dos dólares a cada tonelada de trigo importado, los comerciantes no transportarán el trigo a no ser que la diferencia de precios entre los dos mercados sea de al menos dos dólares. Por tanto, podemos definir el impacto esencial de un arancel con la siguiente secuencia de igualdades: 

\eqblock{
\begin{aligned}
    \text{Demanda Mundial}=\text{Oferta Mundial}\\[1em]
    \hspace{2em}\\[1em]
    \begin{aligned}
        &\not\exists t_h\quad \Longrightarrow\quad p_h^0=p_b^0\\[1em]
        &\exists t_h\quad \Longrightarrow\quad p_h^1-p_b^1=t_h
    \end{aligned}
\end{aligned}
}{}

La introducción de un arancel provoca una diferencia del precio en los dos mercados. El arancel eleva el precio en nuestro país hasta $p^h_1$ y reduce el precio en el extranjero hasta $p_b^1-t_h$. Por tanto, en nuestro país los productores ofrecen más a un precio más elevado, mientras que los consumidores demandan menos, por lo que se demandan menos importaciones. En el extranjero, el menor precio conduce a una reducción de la oferta y a un aumento de la demanda y, de ese modo, a una menor oferta de exportaciones. El incremento del precio en nuestro país es menor que la cuantía del arancel porque parte del arancel se refleja en la reducción del precio de las exportaciones del extranjero y, de ese modo, no se traslada a los consumidores nacionales.\footnote{Este es el resultado normal de un arancel y de cualquier política comercial que limita las importaciones. Sin embargo, la magnitud de este efecto sobre el precio de exportación es a menudo muy pequeña en la práctica. Cuando un país pequeño impone un arancel, su cuota del mercado mundial del bien que importa es generalmente de menor magnitud desde un principio, por lo que la reducción de sus importaciones tiene un efecto muy reducido sobre el precio mundial (de exportación del país extranjero).
Los efectos de un arancel en el caso de un país pequeño son más sencillas: como no puede afectar a los precios de exportación, un arancel aumenta el precio del bien importado en la cantidad total del arancel. Por tanto, la producción aumenta, mientras que el consumo se reduce. La consecuencia del arancel, pues, es que las importaciones disminuyen en el país que lo impone.
} En definitiva, dado el volumen de comercio existente tras el arancel, la demanda de importaciones será tal que iguale la demanda de importaciones en nuestro país al precio con arancel con la oferta de exportaciones del extranjero descontado el arancel: $p_1 -Pt = t$.

Un arancel sobre un bien importado aumenta el precio recibido por los productores nacionales de dicho bien. Este efecto es, a menudo, el principal objetivo del arancel (proteger a los productores nacionales frente a los bajos precios resultantes de la competencia de las importaciones). Al analizar la política comercial en la práctica es importante averiguar la protección que realmente proporciona un arancel u otra política comercial. La respuesta se expresa normalmente en porcentaje del precio que existiría con libre comercio.\footnote{Una cuota de importación sobre el azúcar podría, por ejemplo, aumentar el precio recibido por los productores de azúcar de los Estados Unidos en un 35\%. La medición de la protección parece inmediata en el caso de un arancel: si el arancel es un impuesto ad valorem proporcional al valor de las importaciones, el mismo tipo arancelario debería medir la cuantía de la protección; si el arancel es de cuantía fija, al dividir el arancel entre el precio, sin dicho arancel, se obtiene el equivalente ad valorem. Sin embargo, al intentar calcular la tasa de protección de esta forma tan sencilla aparecen dos problemas. En primer lugar, si el supuesto de país pequeño no es una buena aproximación, parte del efecto de un arancel será la reducción de los precios extranjeros de exportación sin incrementar los precios nacionales. El efecto de las políticas comerciales sobre los precios ·extranjeros de exportación es a veces importante. El segundo problema proviene del hecho de que los aranceles pueden tener efectos muy diferentes sobre las distintas fases de producción de un bien. Un ejemplo sencillo ilustra esta cuestión. Supongamos que un automóvil se vende en el mercado mundial por 8.000 dólares, y que los componentes de los que está hecho se comercializan por 6.000 dólares. Compararemos ahora dos países: uno que quiere desarrollar una industria de ensamblaje de automóviles y otro que ya tiene una industria de ensamblaje y aspira a desarrollar una industria de componentes. Para fomentar una industria nacional del automóvil, el primer país establece un arancel del 25\% a los automóviles importados, lo que permite a los ensambladores nacionales fijar un precio de 10.000 dólares en vez de 8.000. En este caso, sería falso decir que los ensambladores reciben solamente el 25\% de protección. Antes del arancel, la industria de ensamblaje nacional solo se desarrollaría si se pudiera ensamblar por 2.000 dólares (la diferencia entre el precio del automóvil completo, 8.000 dólares, y el coste de los componentes, 6.000 dólares) o menos; ahora funcionará incluso si cuesta 4.000 dólares (la diferencia entre el precio de 10.000 dólares y el coste de los componentes). Es decir, el tipo arancelario del 25\% protege a los ensambladores con una tasa de protección efectiva del 100\%. Supongamos ahora que el segundo país, para fomentar la producción nacional de componentes, impone un arancel del 10\% a la importación de componentes, con lo que el coste de los componentes para los ensambladores nacionales aumenta de 6.000 a 6. 600 dólares. Aunque no hay cambio en el arancel sobre los automóviles ensamblados, esta pol ítica hace menos rentable el ensamblaje nacional. Antes del arancel habría resultado ventajoso ensamblar un coche en el país si se pudiera hacer por 2.000 dólares (8. 000 dólares - 6.000 dólares); después del arancel el ensamblaje local tiene lugar solo si puede ser realizado por 1 .400 dólares (8.000 dólares - 6.600 dólares). El arancel sobre los componentes, aunque proporciona una protección positiva a los fabricantes de componentes, ofrece una protección negativa a los ensambladores a una tasa de -30\% (-600/2.000). Con un razonamiento análogo al de este ejemplo, los economistas han llegado a elaborar cálculos para medir el grado de protección efectiva proporcionado realmente por los aranceles y otras políticas comerciales a industrias concretas. 
} Un arancel incrementa el precio de un bien en el país importador y lo reduce en el país exportador. Debido a estos cambios de precios, los consumidores pierden en el país importador y ganan en el exportador. Los efectos netos de un arancel sobre el bienestar se resumen:

\imagenfit{CyB_T_Importaciones.png}{Costes y beneficios de un arancel para el país importador}{}

Los efectos negativos consisten en los dos triángulos b y d. El primer triángulo es una pérdida debida a la distorsión de la producción, resultante del hecho de que el arancel conduce a los productores nacionales a producir demasiado de este bien. El segundo triángulo es una pérdida debida a la distorsión del consumo, derivada del hecho de que un arancel hace que los consumidores consuman demasiado poco del bien . Frente a estas pérdidas se debe considerar la ganancia debida a la relación de intercambio medida por el rectángulo e, que resulta de la reducción del precio de la exportación extranjera causada por el arancel. En el importante caso de un país pequeño que no puede influir significativamente en los precios extranjeros, este último efecto desaparece; por tanto, los costes de un arancel exceden sin ambigüedad a sus beneficios:


\paragraph{Hipótesis STOLPER-SAMUELSON}


\paragraph{Paradoja de METZLER}



\subsubsection{Cuotas a la importación} 
Una cuota de importación es una restricción directa de la cantidad que se puede importar de algún bien. La restricción es impuesta normalmente mediante la concesión de licencias a algún grupo de individuos o empresas.\footnote{Por ejemplo, los Estados Unidos tienen una cuota de importación de queso extranjero. Las únicas empresas que pueden importar queso son algunas compañías comerciales, cada una de las cuales tiene derecho a importar un máximo de kilogramos de queso al año; el tamaño de la cuota de cada empresa está determinado por la cantidad de queso que importó en el pasado. En algunos casos importantes, especialmente el azúcar y la ropa, el derecho a vender en los Estados Unidos se concede directamente a los gobiernos de los países exportadores.
} La diferencia entre una cuota y un arancel es que, con una cuota, el estado no recibe ingresos. Cuando se utiliza una cuota en vez de un arancel para restringir las importaciones, la cantidad de dinero que habría aparecido como ingresos del estado con un arancel es recaudada por quienquiera que reciba las licencias de importación. 

Los poseedores de licencias pueden comprar productos importados y volver a venderlos a un precio más elevado en el mercado nacional. Por ello es importante evitar el error de creer que las cuotas de importación limitan las importaciones sin aumentar los precios nacionales. Lo cierto es que una cuota de importación siempre aumenta el precio nacional del bien importado. Cuando se limitan las importaciones, la consecuencia inmediata es que, al precio inicial, la demanda del bien excede a la oferta nacional más las importaciones. Esta circunstancia provoca un alza de precios hasta que se equilibra el mercado. Al final, una cuota de importación aumentará los precios nacionales en la misma cantidad que un arancel que límite las importaciones hasta el mismo nivel. 

En la valoración de los costes y los beneficios de una cuota de importación es crucial determinar quién obtiene las rentas. Cuando los derechos de vender en el mercado nacional son asignados a gobiernos de los países exportadores, como a menudo ocurre, la transferencia de rentas al exterior hace que el coste de una cuota sea sustancialmente más elevado que el del arancel equivalente.

\paragraph{Restricticciones voluntarias a la exportación}
Una variante de la cuota de importación es la restricción voluntaria de la exportación (RVE), que también se conoce como un acuerdo de restricción voluntaria (ARV). Una RVE es una cuota al comercio impuesta por parte del país exportador en vez del país importador.\footnote{El ejemplo más conocido es la limitación de las exportaciones de automóviles a los Estados Unidos, puesta en práctica por Japón a partir de 1981.

En 1979, el brusco incremento del precio del petróleo y la escasez temporal de gasolina provocaron un rápido desplazamiento del mercado de los Estados Unidos hacia automóviles más pequeños. Los fabricantes japoneses, cuyos costes se habían ido reduciendo en relación con sus competidores estadounidenses, se aprestaron a atender la nueva demanda. A medida que la participación de los japoneses en el mercado aumentaba y la producción estadounidense disminuía, importantes fuerzas políticas en los Estados Unidos exigían protección para la industria nacional. En lugar de actuar unilateralmente y arriesgarse a entablar una guerra comercial, el gobierno estadounidense pidió a las autoridades japonesas que limitaran sus exportaciones. Los japoneses, ante el temor a posibles medidas proteccionistas unilaterales de los Estados Unidos, acordaron lintitar sus ventas. El primer acuerdo, en 1 98 1 , limitó las exportaciones japoneses a los Estados Unidos a 1 ,68 millones de automóviles. Una revisión del acuerdo aumentó el total hasta 1 ,85 millones en 1 984. En 1 985 se permitió el cese del acuerdo.
} 

Desde un punto de vista económico, una restricción voluntaria a la exportación es exactamente como una cuota de importación en la que las licencias son asignadas a los gobiernos extranjeros y, por tanto, resulta muy costosa para el país importador. Una RVE es siempre más cara para el país importador que un arancel que limite las importaciones en la misma cantidad ya que, lo que habrían sido ingresos bajo el arancel, se convierten en rentas ganadas por los extranjeros en virtud de la RVE, por lo que la RVE produce claramente una pérdida para el país importador. 


Un estudio de los efectos de las tres principales restricciones voluntarias a la exportación de los Estados Unidos (textiles y confección, acero y automóviles) concluyó que las rentas ganadas por los extranjeros son aproximadamente dos tercios del coste para los consumidores derivadas de estas restricciones. En otras palabras, el grueso del coste representa una transferencia de renta más que una pérdida de eficiencia. Este cálculo también pone el acento en que, desde el punto de vista nacional, las RVE son mucho más costosas que los aranceles. Por tanto, la p referencia general de los gobiernos por las RVE sobre otras medidas de política comercial exige un detenido análisis. Algunos acuerdos voluntarios de exportación afectan a más de un país. El acuerdo multilateral más conocido es el Acuerdo Multifibras, que limitaba las exportaciones textiles de 22 países hasta principios de 2005. Estos acuerdos multilaterales de restricciones voluntarias se conocen también por otra abreviatura de tres letras, AOM, acuerdos de ordenación de mercado.

Las restricciones voluntarias a la exportación son impuestas, normalmente, por exigencia del importador, y aceptadas por el exportador para evitar otras restricciones comerciales. 

\begin{specialblock}{Restricciones Voluntarias a la Exportación -- Paneles solares chinos}
    Aunque las restricciones voluntarias a las exportaciones no están ya permitidas en virtud de las normas de la OMC, esta resolución se aplica únicamente a un acuerdo negociado por los gobiernos e impuesto a los exportadores. Recientemente, un conflicto comercial entre China y la Unión Europea sobre el rápido incremento de las exportaciones chinas de paneles solares fue resuelto por el acuerdp de los fabricarites chinos de limitar sus exportaciones a la UE de paneles solares con potencia inferior a 7-gigavatios al año, junto con una base mínima de précio para esas unidades. Los fabricantes europeos de paneles, solares se sintieron defraudados, dado que este acuerdo no incluía la imposición de tasas antiqumping del 47\% a todas las importaciones de paneles solares chinos (para contrarrestar la amenaza generada por las concesiones a los fabricantes chinos de paneles solares). Sin embargo, la imposición de tasas antidumping habría suscitado importantes represalias chinas, cuyas autoridades ya habían elaborado una lista de productos europeos, entre ellos bienes de lujo y vinos, a los que se sometería a elevados aranceles de importación. Las autoridades convencieron a los fabricantes chinos de que convinieran en limitar las exportaciones y en aceptar el precio mínimo, ya que ello permitiría mantener los precios más altos que se pagaban en la UE. Quienes resultaron perjudicados por el pacto fueron los consumidores europeos (y el medio ambiente), que pagan un precio sustancialmente superior por la energía solar.
\end{specialblock}

Como veremos en el capítulo 1 0, las ventajas políticas y legales de las RVE han hecho que sean los instrumentos preferidos de la política comercial en algunas ocasiones. 


\subsection{Instrumentos no convencionales}
Existen muchos otros modos en los que el estado puede influir en el comercio.

\subsubsection{Subsidios a la exportación}
Un subsidio a la exportación es un pago realizado a una empresa o individuo que vende un bien en el extranjero. Como un arancel, un subsidio a la exportación puede ser fijo (una cantidad fija por unidad) o ad va/orem (una proporción del valor exportado). Cuando el estado ofrece un subsidio a la exportación, los vendedores exportarán el bien hasta el punto en q ue los precios nacionales excedan a los extranjeros en la cantidad del subsidio.

\imagenfit{Efectois_Subsidios_X.png}{Efectos de un subsidio a la exportación}{}

Los efectos sobre los precios de un subsidio a la exportación son exactamente los opuestos a los de un arancel (Figura 9. 1 1 ). El precio en el país exportador aumenta de PM a Ps, pero, dado que el precio en el país importador se reduce de PM a P!, el incremento del precio es menor que el subsidio. En el país exportador, los consumidores resultan perjudicados, fos productores ganan y el estado pierde, porque debe gastar dinero en el subsidio. La pérdida de los consumidores es el área a + b; la ganancia de los productores es el área a +· b + e; el subsidio del estado es el área b + e + d + e + f + g. Por tanto, la pérdida neta de bienestar es la suma de las áreas b + d + e + f + g. De estas, b y d representan las pérdidas debidas a las distorsiones de la producción y del consumo del mismo tipo que las que produce un arancel. Además, y en contraste con el arancel, el subsidio a la exportación empeora la relación de intercambio del país al reducir el precio de las exportaciones en el ·mercado exterior de P M a P s, Esto genera pérdidas adicionales debidas a la relación de intercambio e + f + g, iguales a PM - PJ veces la cantidad exportada con el subsidio. Por tanto, un subsidio a la exportación conlleva, sin ambigüedad, un coste que supera sus beneficios.

\begin{specialblock}{La Política Agraria Común de la Unión Europea} 
    En 1 957, seis países de Europa occidental (Alemania, Francia, Italia, Bélgica, Países Bajos y Luxemburgo) constituyeron la Comunidad Económica Europea, que desde entonces ha crecido para incluir a la mayor parte de Europa. Ahora l lamada la Unión Europea (U E), sus dos mayores efectos se producen en relación con la política comercial. Primero, los miem bros de la UE eliminaron los ,aranceles entre sí, para crear una unión aduanera (que. se explica en el próximo capítulo). Segundo, la política agrícola de la tJE se ha convertido en un gigantesco programa ,de subsidios a la exportación . La Política Agrícola Común (PAC) de la UE comenzó no'como un subsidio a la exportación, sino como un:esfuerzo para garantizar precios elevados para los agricultores europeos, mediante la. compra. por parte de la UE de los productos agrícolas cuando los precios cayeran por debajo de determinados niveles de apoyo. Para prevenir 1 la atracciór;i de gr.andes cantidades de importaciones, esta política fue apoyada infoialmente mediante aranceles que compensaban la diferencia entre los precios agrícolas europeos.y los mundiales. Sin embargo, desde 1970, los precios de apoyo establecidos por.la DE. han resultado tan elevados que Europa, que con libre comercio sería un importador. neto de muchos pr,oducfos agrícolas, producía más de lo que los consumidores estaban dispμestos a comprar. La consecuencia fue que la propia UE se vio obligada a comprar" y almacenar enormes cantidades de alimentos. 
    
    A finales de 1985, las naciones europeas habían almacenado 780.000 toneladas de ternera, 1,2 millones de toneladas de mantequilla y 1 2 millones de toneladas de trigo. Para evitar el crecimiento ilimitado de estos excedentes acumulados, la UE pasó a una política de subsidio a las exportaciones para desprenderse del exceso de producción. La figura muestra cómo funciona la PAC:
    
    \imagenfit{PAC_Ejemplo.png}{Política Agrícola Común (PAC) europea}{}
    
    Es, por supuesto, exactamente igual que el subsidio a la exportación, con la salvedad de que Europa sería realménte un importador con libre comercio. El precio de apoyo se establece no solamente por encima del precio mundial que prevalecería sin dicho precio de apoyo, sino también por encima del precio que igualaría la demanda y la oferta sin importaciones. Para exportar el excedente resultante se paga un subsidio a la exportación que compensa la diferencia entre los precios europeos y los mundiales. Las exportaciones subsidiadas tienden a reducir el precio mundial, con lo que el subsidio necesario se incrementa. Un estudio reciente estimaba que el coste en términos de bienestar para los consumidores europeos excedía las ventajas para los productores agrícolas en casi 30.000 millones de dólares (21.500 millones de euros) en 2007.

    Pese a los considerables costes netos de la PAC para los consumidores y contribuyentes europeos, la fuerza política de los agricultores en la UE es tan grande que el programa ha encontrado pocos desafíos internos eficaces. Una importante fuente de presión contra la PAC ha procedido de los Estados Unidos y otros países exportadores de alimentos, que se quejan de que los subsidios europeos a la exportación han impulsado a la baja el precio de sus propias exportaciones. Las consecuencias presupuestarias de la PAC también han planteado dudas: en 2013 la PAC costó a los contribuyentes europeos casi 78.000 millones de dólares (58.000 millones de euros), y esa cifra no incluye los costes indirectos que pagan los consumidores de alimentos. Los subsidios públicos a los agricultores europeos representan aproximadamente el 22\% del valor de su producción agrícola, el doble que la cifra estadounidense del 8,6\%. (Los subsidios agrícolas estadounidenses se concentran en un conjunto limitado de cultivos). 
    
    Las recientes reformas de la política agrícola de Europa representan un esfuerzo para reducir la distorsión de los incentivos provocados por el apoyo a los precios, sin dejar de ofrecer ayudas a los agricultores.
\end{specialblock}


\subsubsection{Otros instrumentos}


\paragraph{Exigencias de contenido local}
Una exigencia de contenido local es una regulación que exige que una fracción específica de un producto final sea producida dentro del país. En algunos casos, esta fracción se especifica en unidades físicas, como la cuota de importación de petróleo en los Estados Unidos en los años sesenta. En otros casos, el requisito se establece en función del valor, al exigir que una cuota mínima del precio del bien derive de un valor añadido nacional. Las leyes de contenido local han sido aplicadas ampliamente por los países en desarrollo que intentan transformar su base manufacturera desde el ensamblaje hacia los bienes intermedios. En los Estados Unidos se propuso una ley de contenido local para los automóviles en 1 982, pero nunca fue aprobada. Desde el punto de vista de los productores nacionales de componentes, una regulación del contenido local proporciona protección en el mismo sentido en que lo hace una cuota. Sin embargo, desde la perspectiva de las empresas que deben comprar localmente, el efecto es un poco distinto. El contenido local no constituye un lím ite estricto de las importaciones. Al contrario, permite a las empresas importar más, con tal de que también compren más en el interior. Esto significa que el precio efectivo de los factores productivos de la em presa es un promedio de los precios de los inputs importados y los producidos en el país. Consideremos, por ejemplo, el anterior ejemplo del automóvil en el que el coste de los componentes importados es de 6.000 dólares. Supongamos que comprar los mismos componentes en el país costara 10.000 dólares, pero que se exigiera a las em presas de ensamblaje util izar un 50\% de componentes nacionales. Entonces, tendrían un coste medio de los componentes de 8.000 dólares (0,5 x 6.000 dólares + 0,5 x 1 0.000 dólares), que quedaría reílejado en el precio final del coche. La cuestión importante es q ue un req uisito de contenido local ni produce ingresos para el estado, ni rentas de cuotas de importación. En lugar de ello, la diferencia entre el precio de los bienes importados y de los nacionales se incluye, de hecho, en el precio final y se traslada a los consumidores.\footnote{Una innovación interesante en las regulaciones del contenido local ha sido permitir a las empresas satisfacer sus exigencias de ·contenido local mediante exportaciones de componentes nacionales en lugar de usarlos en la producción. Esto ha sido importante en varios casos. Por ejemplo, las empresas automovilísticas de los Estados Unidos que operan en México han preferidoexportar algunos componentes desde México a los Estados Unidos, aunque esos componentes pudieran ser producidos en los Estados Unidos a un menor coste, porque eso les permite usar un menor contenido mexicano en la producción de coches, en México, para el mercado mexicano.
}


\paragraph{Subsidios al crédito a la exportación}
Es como un subsidio a la exportación, pero adopta la forma de un préstamo subsidiado al comprador. Los Estados Unidos, como muchos países, tienen una institución estatal, el Banco de Exportación e Importación, que se dedica a proporcionar préstamos, siempre algo subsidiados, para ayudar a las exportaciones.

\paragraph{Compras estatales}
Las compras del estado o de las empresas fuertemente reguladas pueden ser dirigidas hacia bienes producidos en el país, aun cuando esos bienes sean más caros que los importados. El ejemplo clásico es la industria de telecomunicaciones europea. Las naciones de la Unión Europea tienen, en principio, libre comercio entre sí. Sin embargo, los principales compradores de equipos de telecomunicaciones son las compañías telefónicas y, en Europa, esas compañías han sido hasta hace poco propiedad del estado. Estas compañías telefónicas de propiedad estatal compran a los proveedores nacionales incluso cuando esos proveedores establecen precios más elevados que los de otros países. La consecuencia es que hay muy poco comercio en equipos de telecomunicaciones en Europa.

\paragraph{Barreras administrativas}
A veces un gobierno quiere restringir las importaciones sin hacerlo formalmente. Por suerte o por desgracia, es fácil complicar los procedimientos sanitarios, de seguridad y aduaneros normales para establecer sustanciales obstáculos al comercio. El ejemplo clásico es el decreto francés de 1982 por el que todos los reproductores de vídeo japoneses debían pasar a través del pequeño puesto aduanero de Poitiers, lo que en la práctica limitaba las importaciones.


\clearpage
\section{Política Comercial: Efectos}
\puntosclave{

}

\textcolor{red}{\textbf{OJO!}} Como puente natural entre la parte anterior de la exposición y esta, resulta muy conveniente presentar el ejemplo del azúcar en EEUU como forma de entender las dificultades que acarrea el análisis teórico de la Política Comercial. ¿Cuántos sectores se ven afectados? ¿A qué otorgamos prioridad? ¿Cómo cambia el Sector Exterior ante modificaciones de la Política Comercial?

\subsection{Marco de Equilibrio Parcial: efectos sobre la estructura del mercado}
Estudiemos en primer lugar que efectos tiene la Política Comercial sobre la estructura competitiva del mercado. Para ello, presentaremos brevemente el Modelo de MELITZ-OTTAVIANO, enmarcado dentro de la Nueva Teoría del Comercio Internacional. 

Por su extensión, nos vemos impedidos de mostrar todo el desarrollo analítico del Modelo,\footnote{Para ver todo el desarrollo: [\authorfont{Ver Anexo \ref{anx:melitz_aranceles}}]
} por lo que iremos directamente a la estructura funcional que nos de las intuiciones económicas que buscamos.

¿Cómo podemos visualizar los efectos de un arancel? Lo primero que caba recordar es que un arancel es un impuesto directo sobre el consumo de un determinado bien que se ha exportado del exterior. Por tanto, ¿quién internaliza este coste? Esencialmente dos agentes: (i) el consumidor nacional; y, (ii) el productor extranjero. Como valoraremos la incidencia sobre el consumidor nacional más adelante, centremos en el segundo agente. 

Considerando un arancel ad-valorem: ¿cómo lo experimenta? El productor extranjero lo verá como un coste adicional marginal sobre cada unidad vendida en nuestro país. Sea $h$ nuestro país, cualquier productor extranjero (considerando que no hay trato preferencial), verá incrementado su coste marginal de distribución en nuestro país en la siguiente cuantía:

\eqblock{
\begin{aligned}
    \underset{\text{\scriptsize barrera comercial: iceberg costs + arancel}}{\tau^h=(\tau_b^h+t_i)}\quad\Longrightarrow\quad \boxed{c_X^h=c\,\tau^h<c_D^h=c}
\end{aligned}
}{}

Supongamos, por ahora, que se cumple la condición: las barreras comerciales son tan elevadas como para que ningún productor extranjero pueda exportar a nuestro mercado nacional. Veamos como se estructura nuestro mercado cautivo.

\subsubsection{Demanda: Curva de demanda lineal}
Supongamos que la demanda del conjunto de variedades viene gobernada por una curva de demanda lineal:

\eqblock{
\begin{aligned}
    &p_i^h = \alpha - \gamma x_i^h - \eta X^h\quad\Longrightarrow\quad \boxed{x_i^h = \frac{L^h}{\gamma}(p_{\max}^h - p_i^h)}\\[2em]
    &\text{donde,}\quad p_{\max}=\alpha-\eta X^h=\frac{\gamma\alpha+\eta N^h\bar p}{\eta N^h+\gamma}\qquad 
    \left\{\begin{aligned}
        &N^h&&\equiv\text{Variedades en el mercado $h$}\\[0.5em]
        &\bar p^h&&\equiv \text{Precio medio de las variedades en el país $h$}\\[0.5em]
        &X^h&&\equiv\text{Gasto total en la cesta de variedades, $h$}
    \end{aligned}\right.\\[2em]
    &\Longrightarrow 
    \text{Bienestar}:
        \left\{\begin{aligned}
            &Corr(\cdot ,p_{\max})&&<0&&\equiv\text{}\\[0.5em]
            &Corr(\cdot, F)&&<0&&\equiv\text{}\\[0.5em]
            &Corr(\cdot ,N^h)&&>0&&\equiv\text{}\\[0.5em]
            &\underset{\text{\scriptsize (efecto dominado)}}{Corr(\cdot,\sigma_p^2)}&&>0&&\equiv\text{}
        \end{aligned}\right.
\end{aligned}
}{Demanda: Problema del consumidor y demanda residual (MELITZ y OTTAVIANO, 2008). Política Comercial: Instrumentos y efectos}

La demanda residual de cada variedad depende negativamente de su propio consumo y del consumo agregado. La curva de demanda inversa depende de la elasticidad de sustitución entre variedades y de la elasticidad de sustitución entre bien numerario ($\alpha$) y conjunto de consumo de variedades ($\eta$). Como el precio máximo depende a su vez del número de variedades disponibles en el mercado y de su precio medio, vemos que en el modelo la competencia es endógena y actúa reduciendo $p_{\max}$, a través del precio medio: cuanto menor es $\bar p$, menor es $p_{\max}$ y más dura es la competencia (mayor elasticidad).

\subsubsection{Productor}
Por lo que respecto al comportamiento de las empresas, dado que la demanda es lineal, el problema de la empresa es trivial. La condición de primer orden implica:

\eqblock{
\begin{aligned}
    &\underset{\text{\scriptsize (función de distribución PARETO: CDF)}}{\text{Tecnología productiva en $h$}}:\quad 
    \underset{\text{\scriptsize podemos suponer que $\varphi_M=0$ y por tanto, $c_M=\infty$. OJO: Graficaremos sobre $\varphi$ los beneficios, pero es equivente, por eso cabe mencionarlo}}{\left\{\begin{aligned}
        &\Pr[c\leq c_M]=C(c)\in[0,1],\,\,c\in[0,c_M]\\[0.5em] 
        &\Pr[\varphi\geq \varphi_M]=1-\Pr[\varphi_M\leq \varphi]:C\left(\frac{1}{\varphi}\right)\in[0,1],\,\,\varphi\in[\varphi_M,\infty)
    \end{aligned}\right.}\\[1em]
    &\underset{\text{\scriptsize (a nivel operativo $\not\Longrightarrow$ nivel final de empresas)}}{\text{En un entorno competitivo}}:\quad  \underset{\text{\scriptsize (infinitesimal) máximo para tener demanda}}{\boxed{c_D^h=p_i^h(c_D^h)\leq p^h_{\max}}}\\[1em]
    &\text{Cantidad producida}:\quad x_i^h(c) = \frac{L^h}{\gamma}(p_D^h(c) - c)\\[0.5em]
    &\text{Precio (con mark-up)}:\quad  p_D^h(c) = \frac{1}{2}(c_D^h + c)
\end{aligned}
}{Productor: Umbral de costes y relaciones estructurales (MELITZ y OTTAVIANO, 2008). Política Comercial: Instrumentos y efectos}

Por tanto, el precio depende del coste marginal que enfrente la empresa y del coste umbral del mercado $c_D$. A su vez, el umbral de corte queda determinado por la estructura competitiva del mercado, que lo iguala al precio máximo con demanda residual nula. Por tanto, operando con las ecuaciones podemos obersar que todo todo comportamiento de la empresa está efectivamente gobernado por la distancia que enfrenta con la frontera $c_D$: 

\eqblock{
\begin{aligned}
    &\mu_D^h(c) = \frac{1}{2}(c_D^h - c), \quad
    &x_D^h(c) = \frac{L}{2\gamma}(c_D^h - c), \qquad
    &\underset{\text{\scriptsize beneficio operativo}}{\boxed{\Pi_D^h(c) = \frac{L}{4\gamma}(c_D^h - c)^2}}.
\end{aligned}
}{Productor: Beneficio operativo (MELITZ y OTTAVIANO, 2008). Política Comercial: Instrumentos y efectos}

Por tanto, cuanto más eficiente es el proceso productivo de la empresa (más bajo $c$), más produce, más margen tiene y más gana.

\paragraph{Condición de libre entrada: Beneficio esperado}
Entonces, ahora que conocemos la estructura del mercado, debemos saber como acabe determinándose. Si bien es cierto que el modelo original diferencia entre número de potenciales entrantes y número de entrantes efectivos, podemos resumir dichos desarrollos en lo que conocemos como la condición de libre entrada, que se proyecta tanto a nivel operativo como a nivel absoluto. 

Como en un contexto de competencia monopolística donde existe un continuo de potenciales entrantes, el beneficio esperado de las firmas debe ser nulo y, en consecuencia, el umbral $c_D$ se determina por la condición de libre entrada (cubrir costes fijos):\footnote{\textbf{NOTA.-} Las relaciones estructurales presentadas pueden parecer confusas. Por eso lo mejor es entender que existen tres dimensiones que se interrelacionan endógenamente: (i) fijamos el umbral de costes mínimos ($c_D^h$) en base a la condición de libre entrada, el beneficio esperado de las firmas debe ser nulo y el $c_D$ resultante determinará la composición de firmas que permanecen en equilibrio; (ii) el umbral de costes fija el número de firmas activas (variedades) que el mercao es capaz de sostener dada la estructura de demanda ($N^h$) ya que, cuánto menor sea el umbral de costes capaces de cubrir la condición de libre entrada, mayor será el número de firmas que el mercado será capaz de sostener; y, (iii) definimos la tasa de supervivencia por cada cohorte de potenciales entrantes que, evidentemente, será menor cuánto menor sea el umbral de costes (relación positiva).

La intución económica que subyace es muy potente: un cuto ff más bajo no significa menos beneficios en el sentido relevante para la entrada, sino que significa selección más dura: es más difícil sobrevivir. Pero, al mismo tiempo, las firmas que sí sobreviven operan en un mercado en el que la escala y la rentabilidad de las más eficientes sostienen más competencia agregada. Por eso el modelo combina dos resultados que a primera vista parecen tensionados, pero no lo están: menor probabilidad de supervivencia y mayor número de firmas activas. \textbf{Uno es un resultado sobre la fracción de entrantes que pasa el filtro; el otro es un resultado sobre el tamaño del equilibrio una vez operan entrada, selección y competencia}.
}

\eqblock{
\begin{aligned}
    \underset{\text{\scriptsize ¿Quién? Selección $\Longrightarrow$ Determina $c_D$ de forma única}}{\int_0^{c_D^h} \Pi_D^h(c)\, \mathrm dC(c) = F}\quad\Longrightarrow\quad
    \left\{\begin{aligned}
        &\underset{\text{\scriptsize ¿cuánto aguante el mercado? Número de variedades efectivas}}{N^h = \frac{2\gamma}{\eta}\frac{\alpha - c_D^h}{c_D^h - \bar c}}\\[0.5em]
        &\underset{\text{\scriptsize ¿Cuántas empresas desean entrar? Efecto localización}}{N_E^h=\frac{N^h}{C(c_D^h)}}
    \end{aligned}\right.
\end{aligned}
}{Condiciones: Corte tecnológico (selección), estructura de mercado y potenciales entrantes (MELITZ y OTTAVIANO, 2008). Política Comercial: Instrumentos y efectos}

Como vemos, el modelo permite analizar de manera impecable las relaciones que existen entre la productividad del mercado, el atractivo del mercado, la supervivencia en el mercado y el número de empresas que es capaz de sostener el mercado. Todas estas relaciones pueden parecer muy confusas pero en realidad son muy sencillas: se plantea una secuencia vertical de ecuaciones que acaba resolviendo el sistema cerrado de ecuaciones simultaneas que gira en torno al umbral de eficiencia.

\eqblock{
\begin{aligned}
    &\left.\begin{aligned}
        &\left.\underset{\text{\scriptsize Condición de libre entrada: Selección}}{\int_0^{c_D^h} \Pi_D^h(c)\, \mathrm dC(c) = F}\hspace{2em}\right|\hspace{2em}\left.\underset{\text{\scriptsize Relación de equilibrio en Mcdo. Producto}}{N^h = \frac{2\gamma}{\eta}\frac{\alpha - c_D^h}{c_D^h - \bar c}}\hspace{2em}\right|\hspace{2em}\underset{\text{\scriptsize Supervivencia}}{N^h=C(c_D)N_E^h}
    \end{aligned}\right.\\[2em]
    &\text{O, de forma equivalente}:\\[2em]
    &\text{Libre entrada} \;\Rightarrow\; c_D^h\quad\Longrightarrow\quad
    \left\{\begin{aligned}
        &\text{Selección: } &&c \leq c_D^h\\
        &\text{Competencia: } &&N(c_D) \;\Rightarrow\; p_{\max}(N)\\
        &\text{Supervivencia: } &&C(c_D^h)
    \end{aligned}\right|\quad\text{con}\quad N^h = C(c_D^h)N_E^h.
\end{aligned}
}{Desarrollo analítico del problema: Tres ecuaciones principales (MELITZ y OTTAVIANO, 2008). Política Comercial: Instrumentos y efectos}

El umbral de costes satisface la condición de libre entrada (selección, fuertemente gravado por $F$). Este umbral está inversamente relacionado con el número de empresas que es capaz de sostener el mercado ($N^h$) y, de forma indirecta, con el $p_{\max}$ y $\bar p$ del mercado: aumenta el grado de competencia en el mercado (entre empresas más productivas). Además, incide directamente en la tasa de supervivencia en el mercado, que establece, implícitamente, una conexión entre el \textbf{atractivo} del mercado y la tasa de supervivencia efectiva (firmas operantes). Gráficamente, obtendríamos la relación de empresas, beneficios, productividad y condición de libre entrada como:


\subsubsection{Política Comercial: Efectos sobre los márgenes del mercado (microestructura)}
Entonces, ¿qué pasaría si redujesemos nuestras barreras comerciales? ¿qué impacto tendría en el bienestar y la estructura del mercado? Estudiar el impacto de un cambio en nuestra Política Comercial equivale a analizar los cambios que experimenta nuestro equilibrio cerrado y comperarlos en diferentes ámbitos de Política Económica. 

Desde una óptica nacional, si las barreras comerciales fuesen suficientemente pequeñas como para inducir la entrada de competidores exnjeros, la estructura económica del mercado puede descomponerse en dos tipos de empresas: empresas domésticas ($x_D^h$) y empresas extranjeras que operan en nuestro mercado ($x_X^h$). ¿Cómo podemos verlo? Manteniendo la demanda del consumidor constante (pues no cambia), las alteraciones que experimentará nuestro mercado realmente tendrán que ver con dos factores fundamentales: las diferencias en productividad y el número total de firmas operantes. Analíticamente:\footnote{La relación básica entre el umbral de costes marginales del exportador y el umbral de costes marginales de la empresa doméstica es $c_X^b=\frac{p_{\max}^h}{\tau^h}=\frac{c_D^h}{\tau^h}$.

Es decir, el umbral de costes para exportar desde $b$ hacia $h$ es el umbral doméstico del país importador dividido por la barrera comercial. Por tanto, cuanto mayor es $\tau^h$, menor es $c_X^b$, y por tanto más difícil es exportar: las barreras comerciales hacen más difícil que una empresa exportadora alcance el punto de beneficio nulo que una empresa doméstica. De ahí sale inmediatamente la condición de existencia de comercio. Para que haya exportación positiva desde $b$ hacia h, tiene que existir al menos un subconjunto de firmas con costes suficientemente bajos como para cumplir: $c\le c_X^b=\frac{c_D^h}{\tau_h}$.

Como el soporte de costes es $c\in[0,c_M]$, mientras $c_X^b>0$ siempre habrá alguna masa positiva de exportadores potenciales en el continuo. Así que el mero hecho de que $\tau^h$ sea alto no impide la existencia de comercio, lo que ocurre es que la masa exportadora se vuelve arbitrariamente pequeña. Formalmente, la masa de exportadores desde $b$ a $h$ será: $G(c_X^b)\,N_E^b$. Si $\tau^h$ sube, $c_X^b$ cae y también cae $G(c_X^b)$.

Ahora bien, el paper sí da una restricción de selección exportadora: el cutoff exportador debe ser más exigente que el cutoff doméstico.
}

\eqblock{
\begin{aligned}
    &\left.\begin{aligned}
        p_D^h(c)=\frac{1}{2}(c_D^h+c)\\[0.5em]
        p_X^b(c)=\frac{\tau^h}{2}(c_X^b+c)
    \end{aligned}\right\}\underset{\substack{\text{\scriptsize teniendo en cuenta}\\\text{\scriptsize que, $c_X^b=\frac{c_D^b}{\tau^b}$}}}{\Longrightarrow}
    \underset{\text{\scriptsize las barreas dificultan ser rentables en el mercado doméstico}}{\left\{\begin{aligned}
        &p_D^h(c)=\frac{1}{2}(p_{\max}^h+c)\\[0.5em]
        &p_X^b(c)=\frac{1}{2}(p_{\max}^h+\tau^hc)
    \end{aligned}\right.}\Longrightarrow\quad \boxed{c_X^b=\frac{c_D^h}{\tau^h}}\\[2em]
    &\Longrightarrow\underset{\text{\scriptsize vemos una idea clave: si las barreras son suficientemente elevadas, la masa de exportadores se vuelve arbitrariamente pequeña $C(c_X^b)\to0$}}{\boxed{\underbrace{N^h}_{\text{variedades totales en }h}=\underbrace{C(c_D^h)N_E^h}_{\text{productores domésticos en }l}+\underbrace{C(c_X^b)N_E^b}_{\text{exportadores de }b\text{ que venden en }h}}}.
\end{aligned}
}{Comercio exterior: Descomposición de la oferta de mercado (MELITZ y OTTAVIANO, 2008). Política Comercial: Instrumentos y efectos}

Por tanto, con esta especificación podemos estudiar fácilmente las implicaciones ecónomicas que tienen tanto las barreras comerciales como las diferencias intrínsecas entre los países. 

\paragraph{Liberalización bilateral: Tamaño de Mercado}
Supongamos que dos países establecen un acuerdo comercial para reducir los aranceles de forma bilateral. ¿Qué impacto tendrá  esta liberalización en la estructura de mercado en ambos países? Si ambos países tienen la misma tecnología (distribución de PARETO), las consecuencias serán muy diferentes en función del tamaño relativo de ambos mercados.

Si los países son muy parecidos, tienen tamaños iguales, podemos suponer que el impacto que generará dicho acuerdo será beneficioso para ambos. ¿Por qué? Por los fundamentos que rigen la estructura económico del mercado (y el modelo), cabe esperar que en mercados parecidos la dureza competitiva sea similar, al igual que el número de variedades ofertadas y, en consecuencia, el umbral mínimo de costes para operar (eficiencia productiva). 

\eqblock{
\begin{aligned}
    &L^h\sim L^b&&\Longrightarrow\quad c_D^h\sim c_D^b\quad \Longrightarrow\quad N^h\sim N^b\\[1em]
    &\tau^h=\tau^b\equiv \,\,\downarrow\tau&&\Longrightarrow\quad \underset{\text{\scriptsize (\textit{libertad} de acceso)}}{c_D^h(L^h,L^b)}\quad \text{\scriptsize(aumenta el tamaño efectivo de los mercados)}\\[1em]
    &\underset{\text{\scriptsize (estructura de mercado)}}{\Longrightarrow}&&
    \left\{\begin{aligned}
        &\underset{\text{\scriptsize (suponiendo países con tecnologías similares)}}{N^h=\overset{N_D^h}{C(c_D^h)N_E^h}+\overset{N_X^b}{C(c_X^b)N_E^b}}:\qquad c_X^b=\frac{c_D^h}{\tau}\\[0.5em]
        &\underset{}{\uparrow N^h}\quad \Longrightarrow\quad \downarrow \,p_{\max}\,(\text{\scriptsize presión competitiva})\\[0.5em]
        &\underset{}{\downarrow p_{\max}}\quad \Longrightarrow\quad \downarrow \,c_D^h\,(\text{\scriptsize cut off productivo})\\[0.5em]
        &\text{Pero, OJO}:\quad \Delta\{\text{Tamaño},\Pi,\mu\}:\underset{\text{\scriptsize (efectos)}}{\text{Directos}>\text{Indirectos}}
    \end{aligned}\right.\\[1em]
    &\underset{\text{\scriptsize (bienestar)}}{\Longrightarrow}&&
    \left\{\begin{aligned}
        &\underset{}{\downarrow c_D^h}\quad \Longrightarrow\quad \uparrow N^h\quad\Longrightarrow\quad\uparrow X^h\\[0.5em]
        &\underset{}{\downarrow p_{\max}}\quad \Longrightarrow\quad \downarrow \,\bar p\,\,\underset{\text{\scriptsize (dominado)}}{(+\downarrow\sigma_p^2)}\\[0.5em]
    \end{aligned}\right.
\end{aligned}
}{Liberalización bilateral: Efectos sobre dos países similares. Política Comercial: Instrumentos y efectos}


Como vemos, una liberalización bilateral entre dos países similares tiene beneficios potenciales tanto desde el punto de vista de la oferta como desde el punto de vista de la demanda. Las empresas se vuelven más grandes, aumentan sus márgenes y son más productivas a raíz del efecto \textbf{tamaño de mercado} e, indirectamente, se reduce el precio medio (y máximo) y aumenta la variedad. Estas dos últimas afectan significativamente en el bienestar del consumidor que mejora gracias a la disminución del precio medio (deflactor en términos reales del coste de la variedad) y por la disponibilidad de mayor número de variedades (efectos que dominan la menor dispersión de precios).

Ahora bien, si la liberalización bilateral se llevara a cabo entre países de diferente tamaño las implicaciones serían muy diferentes. ¿Por qué? Analíticamente, un proceso semejante conduciría a lo siguiente:

\eqblock{
\begin{aligned}
    &L^h> L^b&&\Longrightarrow\quad c_D^h> c_D^b\quad \Longrightarrow\quad N^h>N^b\\[1em]
    &\tau^h=\tau^b\equiv \,\,\downarrow\tau&&\Longrightarrow\quad \underset{\text{\scriptsize (\textit{libertad} de acceso)}}{c_D^h(L^h,L^b)}\quad \text{\scriptsize(aumenta el tamaño efectivo de los mercados)}\\[1em]
    &\underset{\text{\scriptsize (estructura de mercado)}}{\Longrightarrow}&&
    \left\{\begin{aligned}
        &\downarrow \tau\quad\Rightarrow\quad  \uparrow c_X\quad\Rightarrow\quad \uparrow N_X \quad (\text{en ambos países})\\[0.5em]
        &\Longrightarrow\underset{\text{\scriptsize (afectará igual que antes a las variables base en ambos países)}}{\{N,p_{\max}, c_D^h\}}\\[0.5em]
        &\text{pero: } \underset{\substack{\text{\scriptsize relocalización de entrada:}\\\text{\scriptsize  una mayor proporción de nuevas empresas elige instalarse allí}}}{\quad N_E^h \uparrow\quad >\quad N_E^b}\\[0.5em]
        &\underset{\text{\scriptsize se experimenta cuando una país liberaliza unilateralmente, pero más extremo: NO > SÍ, para los dos}}{\text{¿Por qué?}:\underset{h>b}{\int_0^{c_D}\Pi(c)\mathrm dC(c)}+\underset{h=s}{\int_0^{c_X}\Pi(c)\mathrm dC(c)}=F}\\[1em]
        &\underset{\text{\scriptsize relativamente más atractivio}}{\frac{\mathrm d(N_E^h-N_E^b)}{\mathrm d\rho}\gg0}\qquad\wedge\qquad \frac{\mathrm d(N_D^h - N_D^b)}{\mathrm d\rho}\gg0\quad \wedge\quad \frac{\mathrm d(c_D^h-c_D^b)}{\mathrm d\rho}\ll0
    \end{aligned}\right.\\[1em]
    &\underset{\text{\scriptsize (bienestar)}}{\Longrightarrow}&&
    \left\{\begin{aligned}
        &\underset{}{\downarrow c_D}\quad \Longrightarrow\quad \uparrow N\quad\Longrightarrow\quad\uparrow X\\[0.5em]
        &\underset{}{\downarrow p_{\max}}\quad \Longrightarrow\quad \downarrow \,\bar p\,\,\underset{\text{\scriptsize (dominado)}}{(+\downarrow\sigma_p^2)}\\[0.5em]
    \end{aligned}\right.
\end{aligned}
}{Liberalización bilateral: Efectos sobre dos países similares. Política Comercial: Instrumentos y efectos}

En definitiva, la liberalización bilateral aumenta la competencia en ambos mercados, dando lugar a cambios proporcionales en los umbrales de costes (y, por tanto, a aumentos proporcionales de la productividad agregada) en ambos países. Los efectos de dicha liberalización son cualitativamente idénticos a los que se experimentan con la apertura: la variedad de productos aumenta como resultado de la mayor competencia, lo que también induce una disminución de los markups y de los precios. De nuevo, el bienestar aumenta como consecuencia de la combinación de una mayor productividad, menores markups y una mayor variedad de productos. Si bien los efectos (generales) no dependen del tamaño relativo de los países (siempre que el bien diferenciado se produzca en ambos), las diferencias en tamaño entre países sí inducen cambios importantes en el patrón relativo de entrada tras la liberalización. 

\textbf{NOTA.-} Cabe destacar que este mismo efecto sesgado es precisamente el que se experimentaría en caso de liberalización unilateral de las barreras comerciales. De la misma forma, cuando se lleva a cabo una liberalización preferencial se experimenta un sesgo hacia los países que liberalizan en un contexto multi-país.

\paragraph{Liberalización comercial: Ventaja comparativa}
Por otro lado, los efectos de la Política Comercial pueden analizarse desde el punto de vista clásico de la ventaja comparativa. Desde este punto de vista, los diferentes rendimientos entre países no solo vendrán dados por la combinación del efecto tamaño de mercado y el efecto de las barreras arancelarias, sino también por los efectos inducidos por la ventaja tecnológica. Proyectando directamente la estructura base de la intensidad competitiva (que marca las relaciones de rendimiento de las empresas):

\eqblock{
\begin{aligned}
    &\boxed{c_D^h=\left[\frac{\gamma \Gamma^h}{L^h}+\frac{1}{1+2\rho}\right]^{\frac{1}{k^h+2}}}\quad \text{donde,}
    \left\{\begin{aligned}
        &\Gamma^h(k^h,F^h)\equiv \text{Resume los parámetros tecnológicos}\\[1em]
        &k^h\equiv\text{Distrubución tecnológica PARETO (colas gruesas)}\\[1em]
    \end{aligned}\right.\\[2em]
    &\text{Por tanto}:\qquad \underset{\text{\scriptsize incidencia de los CF}}{\frac{\mathrm dc_D^h}{\mathrm dF^h}>0}\quad\wedge\quad\underset{\substack{\text{\scriptsize Incidencia distribución}\\\text{\scriptsize Colas gruesas}\Rightarrow\text{\scriptsize Menos costes}}}{\frac{\mathrm dc_D^h}{\mathrm dk^h}>0}\qquad\iff\qquad\frac{\mathrm d\Gamma^h}{\mathrm dk^h}<0\quad\wedge\quad\frac{\mathrm d\Gamma^h}{\mathrm dF^h}<0
\end{aligned}
}{}

Como vemos, introducir diferencias tecnológicas implica alterar dos parámetros: o bien el coste de entrada $F^h$, o bien la distribución de costes  $k$ (por ejemplo, un menor $c_M$, que implica mejores tecnologías potenciales). En ambos casos, esto se resume en el parámetro tecnológico agregado que aparece en la condición de libre entrada $\Gamma^h$, de modo que una mejora tecnológica (menor $c_M^h$ o menor $F^h$) implica un menor $\Gamma^h$. En este caso, el país con mejor tecnología (menor $\Gamma^h<\Gamma^b$) tendrá un menor $c_D^h$ en equilibrio, lo que tiene implicaciones inmediatas: solo sobreviven empresas más productivas, mayor productividad media, menores precios y markups (aumentando ventas y beneficios) y mayor competitividad internacional.\footnote{Para poder ver gráficamente la incidencia de la tecnología en los umbrales de corte tecnológico: \href{https://www.geogebra.org/calculator/dvtnwnpg}{GeoGebra.org (Elaboración propia)}
}

Por tanto, el país tecnológicamente avanzado (con menor $c_D^h<c_D^b$) tendrá una mayor masa de empresas capaces de exportar, porque su distribución de costes está más concentrada en niveles bajos. Además, bajo esta misma heterogeneidad podemos estudiar los efectos anteriores (efectos tamaño de mercado, efectos de la liberalización) ya que esta ventaja se proyecta amplificando los efectos anteriores: más entrada, más, competencia y selección más dura ($\downarrow c_D^h$). Es decir, la ventaja comparativa se retroalimenta endógenamente lo que conduce a la generación de \textit{hubs} comerciales, de manera parecida pero amplificado a la generada por los efectos del trato preferencial: los países con mejor tecnología o mejor acceso se convierten en hubs, mientras que los demás quedan en desventaja.


\subsection{Marco de Equilibrio General: efectos sobre el bienestar agregado}
Pero, ¿se ajusta esto a la realidad observada? La mayoría de mercados desarrollados mantienen una Política Comercial favorable a países muchos más pequeños. Además, las barreras comerciales predominan esencialmente entre países en vías de desarrollo que se integran progresivamente en el Comercio Internacional.

Por tanto, aunque el marco anterior parece muy útil para describir (positivo) como cambia la microestructura sectorial ante cambios en la Política Comercial (productividad, integración comercial, las decisiones de localización, etc.); no parece proporcionar un marco normativo muy adecuado para la Política Comercial que cada país debe (y parece) emprender. 

Por ello pasamos a analizar la Política Comercial desde una perspectiva general en el que podamos enfatizar la heterogeneidad real sobre dos fundamentos clave: (i) la elasticidad comercial (margen intensivo y extensivo); y, (ii) el impacto sobre el bienestar vía salarios. Como se puede intuir, distanciarnos de un marco microeconómico y centrarnos en una perspectiva agregada tiene sus ventajas e inconvenientes. En particular, el siguiente enfoque normativo parece ajustarse mucho más a la observación empírica y nos permite evaluar el impacto de la Política Comercial sobre dos indicadores clave en términos de Política Económica: (i) competitividad exterior; y, (ii) 

Brevemente, las ecuaciones principales del modelo se pueden resumir en las siguientes (seguimos la formulación de ARA, 2025).\footnote{Se puede ver el desarrollo completo del modelo traducido en el [\authorfont{Ver Anexo \ref{anx:modelo_ara}}]
}

\subsubsection{Demanda: CES-isoelástica}
Desde el punto de vista de la demanda, consideramos una Función de Utilidad con elasticidad de sustitución constante (CES-isoelástica) y sin la presencia de numerario. De esta forma, la utilidad del consumidor viene exclusivamente proporcionado por el consumo de productos diferenciados y, en particular, por su amor por la variedad.

\eqblock{
\begin{aligned}
    &U_h=\left(\sum_{n=h,b}\int_{i\in I_n}\left(x_i^{nh}\right)^{\rho}\mathrm di\right)^{\frac{1}{\rho}}\quad\overset{\substack{\text{\scriptsize en un marco}\\\text{\scriptsize dos países}}}{\Longrightarrow}\quad 
    U_i=\left(\underset{\text{\scriptsize variedades domésticas}}{\int_{i\in I_h}\left(x_i^{hh}\right)^\rho\mathrm di}+\underset{\text{\scriptsize variedades importadas}}{\int_{i\in I_b}\left(x_i^{bh}\right)^\rho\mathrm di}\right)^{\frac{1}{\rho}}\\[2em]
    &\boxed{x_i^{bh}= R_D \mathbf p_D^{\sigma-1} \left(p_i^{bh}\right)^{-\sigma}}\qquad\text{donde},
    \left\{\begin{aligned}
        &R_D=\mathbf p_DX_D&&\underset{\text{\scriptsize $R_D$ es crítica en el modelo: deja de ser exógena (o proporcional a $L_D^h$) y pasa a depender de $w^h$ y $\tau^{bh}$}}{\equiv\text{Consumo agregado del mercado doméstico, $h$}}\\[1em]
        &\mathbf p_D&&\equiv\text{Índice de precios (deflactor de bienestar), en $h$}
    \end{aligned}\right.
\end{aligned}
}{Demanda: Curva de demanda de cada variedad y problema dual del consumidor (Modelo ARA, 2025). Instrumentos y efectos: Marco general y evaluación sobre el bienestar.}

\subsubsection{Productor: }
Desde el punto de vista del productor, su comportamiento no difiere mucho del planteado por MELITZ (2003). Supondremos que existe un número infinito de posibles entrantes, que todos ellos enfrentan unos costes fijos de instalación y unos costes marginales dependientes de su productividad ($\varphi$, con CDF acotada pero indefinida, para poder generalizar), costes de transporte ($\theta$) y costes arancelarios ($\tau$). Además, todos los costes de producción se expresan en unidades de salario doméstico. Por tanto, maximizando la función objetivo del productor obtendremos la segunda ecuación fundamental, la regla de precios con markup constante, que introduce salarios directamente en el comportamiento de la empresa:

\eqblock{
\begin{aligned}
    &\left.\begin{aligned}
        &\varphi\in[\varphi_{\min},\varphi_{\max}]:C(\varphi)=\Pr[\varphi\leq \bar\varphi]\\[1em]
        &l^{hh}(\varphi)=F^{hh}+\frac{1}{\varphi}x^{hh}:\quad \theta^{hh}=1\\[1em]
        &l^{bh}(\varphi)=F^{bh}+\frac{\theta^{bh}}{\varphi}x^{bh}:\quad \theta^{bh}>1\\[1em]
    \end{aligned}\right\}\quad\overset{CT(\varphi)=w^b\,l^{bh}(\varphi)}{\Longrightarrow}\quad{\Pi^{bh}(\varphi)=\left(\frac{p^{bh}(\varphi)}{\tau^{bh}}-\frac{\theta^{bh}w^b}{\varphi}\right)x^{bh}(\varphi)-w^bF^{bh}}\\[2em]
    &\underset{\text{\scriptsize \textbf{Nuevo canal}: $w^b$ afecta directamente a la \textbf{competitividad} (empresas)}}{\boxed{p^{bh}(\varphi) = \underset{\mu}{\frac{\sigma}{\sigma-1}}\cdot\underset{c}{\theta^{bh}\frac{w^b}{\varphi}}\cdot\underset{\tau}{\tau^{bh}} }}
\end{aligned}
}{Demanda: Curva de demanda de cada variedad y problema dual del consumidor (Modelo ARA, 2025). Instrumentos y efectos: Marco general y evaluación sobre el bienestar.}



\paragraph{Condición de libre entrada}
Además del problema e maximización de la función objetivo del productor, que nos proporciona la regla óptima de fijación de precios, debido a la competencia monopolística sabemos, igual que en el supuesto anterior, que la condición de libre entrada nos obliga a condicionar el equilibrio en variedades (i.e. empresas efectivas en el mercado) como función del beneficio esperado, que a su vez depende de la productividad que observe la empresa al entrar. Por tanto, la condición de corte tecnológico se puede sintetizar en:

\eqblock{
\begin{aligned}
    &\underset{\substack{\\[0.5em]\\ \text{\scriptsize i.e. } r^{hb}(\bar{\varphi}^{hb})=w^hF^{hb}}}{\Pi^{hb}(\varphi)=0}\iff \underset{\text{\textbf{Competitividad} de las empresas domésticas, $h$, respecto $b$}}{\overset{\text{\scriptsize \textbf{Selección}: doméstica-exportación. Obtenemos umbral doméstico con: $b=h$}}{\boxed{B_D^b\left(\tau^{hb}\right)^{-\sigma}\left(\theta^{hb}w^h\right)^{1-\sigma}(\bar{\varphi}^{hb})^{\sigma-1} = \sigma w^{h} F^{hb}}}}\\[2em]
    &\text{donde},\quad B_D^b = R_D^b \left(\underbrace{\frac{\sigma-1}{\sigma}}_{\rho}\mathbf p_D^b\right)^{\sigma-1}\equiv\text{Índice de demanda de mercado (exterior, $b$)}
\end{aligned}
}{}

Esta ecuación nos sirve de referencia por dos razones fundamentales:
\begin{lnum}
    \item \textbf{Corte tecnológico}. En primer lugar, nos proporcionará el umbral de eficiencia mínimo que necesitará observar una empresa para poder operar. Este será diferente en función de si lo evaluamos para el mercado doméstico, ajustando los parámetros necesarios ($\theta^{hh}=\tau^{hh}=1$), o para el mercado exterior: $\bar{\varphi}^{hh}<\bar{\varphi}^{hb}$. Además, la selección (domésticas-exportadoras) prevalece siempre y cuando los costes de exportación sean suficientemente elevados y el tamaño de mercado ($L_D$) no sea excesivamente diferente;\footnote{Esto último se mantiene siempre y cuando la demanda relativa de mercado, $\frac{B_D^b}{B_D^h}$ (que son proporcionales al tamaño de mercado medido como gasto relativo $\frac{B_D^b}{B_D^h}$), no sea muy grande ni muy pequeña.
    } y,
    \item \textbf{Competitividad}. Además, nos permitirá evaluar la competitividad exterior de nuestra economía como función de nuestros salarios. Es decir, como los costes del productor se evaluan siempre en unidades nacionales, establecerá una relación unívoca entre el umbral tecnológico que determina la selección hacia la exportación y los salarios domésticos.
\end{lnum}

Por tanto, esta ecuación de corte tecnológico endogeneiza también los cambios experimentados en salarios ($w^b$, $w^h$) y en gasto agregado, vía $B_D^h(R_D)$. Sería una condición equivalente a la condición de selección en el modelo de MELITZ y OTTAVIANO (recordemos con la ecuación de selección pretendemos obtener el umbral de costes máximos para operar rentablemente, $c_D$). Se desprenden dos implicaciones: (i) nos permitirá relacionar el nivel de salarios domésticos ($w^h$) y el umbral de productividad para la exportación ($\bar{\varphi^{hb}}$) a través de la \textbf{Curva de Competitividad}; y, (ii) nos permite obtener la condición de libre entrada a partir de la suma de los beneficios esperados en la exportación y la producción doméstica. Analíticamente:

\eqblock{
\begin{aligned}
    &\text{Curva de Competitividad}:\boxed{w^h=\mathcal C^{hb}\left(\bar{\varphi}^{hb}\right)^{(\sigma-1)/\sigma}\qquad\text{donde,}\quad \mathcal C^{hb}\equiv \left(\frac{B_D^b\left(\theta^{hb}\right)^{(1-\sigma)}}{\left(\tau^{hb}\right)^{\sigma}\sigma F^{hb}}\right)}\\[2em]
    &\text{Condición de Libre Entrada}:\\[1em]
    &\sum_{n=h,b} w_bF^{hn} \, J_D(\bar{\varphi}^{hn}) = w_bF_E^h\quad\Longrightarrow\quad \underset{\text{\scriptsize $\mathbb E[\Pi^h]$: mercado doméstico}}{\left[F^{hh}J_D(\bar{\varphi}^{hh})\right]}+\underset{\text{\scriptsize $\mathbb E[\Pi^h]$: mercado extranjero}}{\left[F^{hb}J_D(\bar{\varphi}^{hb})\right]}=\underset{\text{\scriptsize Coste fijo soportado (ex ante)}}{F_E^h}\\[2em]
\end{aligned}
}{}

Por un lado, la curva de competitividad nos permite relacionar el salario del país $h$ como función del cutoff tecnológico de sus exportadores hacia $b$, dado el resto de variables relevantes. Esta expresión nos muestra mucho más de lo que parece:
\begin{lnum}
    \item Es creciente en $\bar{\varphi}_{hb}$. Eso quiere decir que, para sostener salarios más altos en $h$, el umbral exportador debe ser más exigente. La lógica es muy clara: si producir en $h$ es más caro porque $w_h$ sube, entonces solo las firmas suficientemente productivas pueden seguir exportando al mercado $b$. En otras palabras, mayores salarios domésticos erosionan la competitividad exterior y empujan hacia arriba el cutoff exportador;
    \item Depende positivamente de $B^b_D$: demanda del mercado de destino. Un mercado exterior más grande o más atractivo permite sostener salarios más altos en $h$ para un mismo cutoff, $\bar{\varphi}^{hb}$. Por tanto, si el mercado $b$ ofrece más demanda, las firmas de $h$ pueden soportar costes salariales más altos sin perder acceso exportador. Recordar, precisamente que $B_D$ depende positivamente del nivel de gasto agregado ($R_D$) y del indice de precios ($\mathbf p_D$);
    \item Depende negativamente de $\tau^{hb}$ y de $\theta^{hb}$, el coste iceberg. Si los mercados están menos integrados o tienen mayores barreras arancelarias, la pendiente de la curva de competitividad $\mathcal C^{hb}$ cae. Por tanto, se exige una mayor productividad para sostener el mismo salario. Así es como la política comercial o los costes comerciales endurecen la condición de acceso al mercado exterior. Eso también está en la selección exportadora del modelo: el cutoff de exportación aumenta cuando exportar se vuelve más costoso;
    \item Cuarto, la curva depende negativamente de $F^{hb}$, el coste fijo de exportar desde $h$ a $b$. Si los costes suben, sostener un mismo cutoff con un salario dado se vuelve más difícil, porque el margen variable debe cubrir un coste fijo mayor. Por eso, para un salario dado, solo sobreviven exportadores aún más productivos.
\end{lnum}

Por otro lado, la condición de libre entrada exige que el beneficio esperado operando en el mercado doméstico ($w_hF^{hh}J^h(\bar{\varphi}^{hh})$) más el beneficio esperado vendiendo al extranjero ($w_hF^{hb}J^h(\bar{\varphi}^{hb})$) sea igual al coste fijo que debe soportar cualquier empresa potencial antes de materializar su productividad (es decir, de experimentar el proceso de selección). ¿Qué significa esto? Esta condición tiene la intención de establecer el umbral mínimo necesario para que la producción de mercado se lleva a cabo en nuestra economía doméstica. Es decir, si el beneficio esperado no es suficiente para cubrir los costes fijos de entrada, nadie experimentará proceso de selección (salida del mercado) porque nadie entrará al mercado. Por tanto, se trata de una ecuación análoga a la expuesta antes (MELITZ y OTTAVIANO) pero, en este caso, se endogeneiza también el efecto de los salario, ya que el umbral tecnológico ($\bar{\varphi}^{hn}$) depende de éstos ($w^b$).

The FE condition determines the market demands Bi;Bj by adjusting the price indices Pi; Pj in equilibrium, so that potential entrants make zero expected profits.

\subsubsection{Equilibrio General: Balanza Comercial}
Con toda esta especificación podemos presentar la ecuación de cierre del modelo como la condición de Equilibrio General. ¿De dónde saladrá esta condición? Como vemos, hemos presentado el equilibrio en la entrada  (beneficio esperado), en el proceso de selección (umbral productivo que es capaz de sostener la estructura de la demanda) dado una regla óptima de fijación de precios y, en última instancia (aunque en primer lugar), el equilibrio en la demanda de cada variedad (vaciamiento de mercados) como función de su precio, del consumo agregado y del índice de precios. Ahora bien, todo el consumo que realiza el consumidor se sostiene por su ingreso en forma de salario que informa todas las ecuaciones del modelo presentadas hasta ahora. ¿De dónde sale este salario? Esta es precisamente la ecuación de cierre del modelo.

El equilibrio en el mercado de trabajo que atiende a la condición de vaciamiento en el mercado, fija el nivel de salarios que sostiene el equilibrio general. La intuición económica es muy potente: se utiliza la condición de vaciamiento en el mercado de trabajo como proxy para la condición de equilibrio en la Balanza Comercial. Por tanto, será el equilibrio exterior el que condicionará el resto de condiciones marginales, como la productividad, la competitividad, el númro de variedades, etc.

\eqblock{
\begin{aligned}
    &L_D^h = \frac{R_D^h-T_D^h}{w^h}\,;\qquad\text{donde,}\,\,
    \underset{\text{\scriptsize ¿Por qué? Porque el gasto doméstico de bienes importados trae incluido el arancel: $R_D^h=\sum_{n=h,b}\tau^{nh}R_D^{nh}$}}{\left\{\begin{aligned}
        &T_D^h=(\tau^{bh}-1)R^{bh}&&\equiv\text{Recaudación agregada}\\[1em]
        &R^{bh}_D&&\equiv\text{Consumo agregado en bienes importados, $b\to h$}
    \end{aligned}\right.}\\[1em]
    &\underset{\text{\scriptsize equivalentemente}}{\Longrightarrow}\quad \boxed{R_D^h=w_iL_D^h+T_D^h}
\end{aligned}
}{}

El salario doméstico ($w^h$) viene determinado por la igualdad entre consumo agregado del mercado de variedades ($R_D^h$) y el ingreso agregado del trabajo ($w^hL_D^h$) más el ingreso agregado por recaudación arancelaria ($T_D^h$). De esta forma, no hay un excedente neto más allá del ingreso por trabajo y recaudación arancelaria ya que la condición de libre entrada deja beneficios nulos en el equilibrio. La condición de vaciamiento en el mercado de trabajo es equivalente al equilibrio de la Balanza Comercial ya que, aunque el equilibrio exterior se simplifica a la igualdad entre gasto doméstico de bienes importados y gasto extranjero de bienes domésticos ($R_D^{bh}=R_D^{hb}$), las dos ecuaciones conducen a la misma igualdad: $R_D^h=w^hL_D^h+T_D^h$. Téngase en cuenta que el ingreso agregado del mercado de trabajo incluye en los ingresos obtenidos por firmas domésticas y los ingresos obtenidos por empresas exportadoras, descontando los aranceles del socio comercial.

\paragraph{Equilibrio: Balanza Comercial (descomposición por cuotas de gasto)}
Teniendo todo esto presente podemos llevar a cabo una serie de simplificaciones que aproximan este equilibrio a algunos de los estadísticos más familiares en el análisis empírico del comercio internacional (y, en consecuencia, esenciales en la evaluación de la Política Comercial):

\eqblock{
\begin{aligned}
    &\left.\begin{aligned}
        &\text{Sea},\,\underset{\text{\scriptsize cuota de importaciones en el gasto agregado (\textbf{incluyendo aranceles})}}{\lambda^{bh}=\frac{\tau^{bh}R_D^{bh}}{\sum_n\tau^{nh}R_D^{nh}}}\\[2em]
        &\text{Sea},\,\underset{\text{\scriptsize cuota de importaciones en el gasto agregado (\textbf{sin aranceles})}}{\tilde{\lambda}^{bh}=\frac{R_D^{bh}}{\sum_nR_D^{nh}}}\\[2em]
    \end{aligned}\right\}&&\underset{\substack{\text{\scriptsize consumos agregados}\\\text{\scriptsize netos de aranceles}}}{\Longrightarrow}\quad R_D^{nh}=\tilde{\lambda}^{nh}w^hL_D^h\\[1em]
    &\left(\substack{\text{\scriptsize definición anterior}}\right)&&\Longrightarrow\quad
    \boxed{\underset{\text{\scriptsize Gasto agregado}}{w^hL_D^h}=\underset{\text{\scriptsize Cuota de gasto doméstico}}{\tilde{\lambda}^{hh}w^hL_D^h}+\underset{\text{\scriptsize Cuota de gasto en bienes importados}}{\tilde{\lambda}^{bh}w^hL_D^h}}\\[2em]
    &\left(\substack{\text{En equilibrio $R^{bh}=R^{hb}$}\\+\\\text{\scriptsize condición LMC}}\right)&&\Longrightarrow\quad
    \underset{\text{\scriptsize De esta forma podemos reescribir el equilibri ocomercial como cuotas de gasto en los países}}{\boxed{\underset{\text{\scriptsize Gasto agregado}}{w^hL_D^h}=\underset{\text{\scriptsize Cuota de gasto doméstico}}{\tilde{\lambda}^{hh}w^hL_D^h}+\underset{\text{\scriptsize Importados por $b$, $h\to b$}}{\tilde{\lambda}^{hb}w^bL_D^b}}}
\end{aligned}
}{}

Por tanto, hay dos formas equivalentes de escribir la restricción agregada. La primera, como igualación del consumo agregado al ingreso generado por aranceles y rentas laborales. La segunda mediante el desarrollo del gasto agregado del país por cuotas de gasto que, por la condición de vaciamiento del mercado laboral (y en consecuencia, de la Balanza Comercial), permite representar el equilibrio descomponiendo el gasto doméstico y el equilibrio comercial bilateral (en \textbf{valor}, $h\to b=b\to h$ : $\tilde{\lambda}^{bh} w^h L^h = \tilde{\lambda}^{hb} w^b L^b$). Por tanto, si empezamos con una identidad de gasto, acabamos con una condición general descompuesta en cuotas gasto y equilibrio comercial. La interpretación económica es muy potente: la renta agregada no solo se gasta en bienes domésticos e importaciones, sino que en equilibrio está estructuralmente relacionada al ingreso/renta del resto del mundo. Es decir, \textbf{lo que un país importa está directamente financiado por lo que es capaz de exportar}. Esto es precisamente lo que introduce el nuevo canal de ajuste: la demanda de importaciones afecta a los salarios extranjeros, y los salarios extranjeros determinan la demanda de nuestras exportaciones.



\subsubsection{Política Comercial: Efectos sobre el bienestar (macroestructura)}
Como es de esperar, esto tendrá unas consecuencias muy profundas a la hora de evaluar el impacto de la Política Comercial, puesto que las interdependencias económicas entre economías incidirán a través del canal renta, algo que no se podía evaluar bajo el marco microeconómico anterior. Entonces, ¿qué impacto tiene la Política Comercial?

Veamos algunos ejemplos. 

\paragraph{Liberalización: Selección doméstica débil}


si se elimina ese bien exterior, los salarios pasan a ser endógenos y el efecto del tamaño de mercado sobre la selección puede invertirse: un país grande puede sufrir una selección doméstica más débil porque los salarios más altos desplazan entrada hacia el exterior y permiten sobrevivir a firmas menos eficientes. Él mismo dice expresamente que su efecto de tamaño de mercado contradice el de Melitz y Ottaviano, y que la razón clave es la presencia del outside good en éstos.

\paragraph{Liberalización: Efecto Mercado doméstico en salarios}






\clearpage
\section{Política Comercial: Economía Política}
\puntosclave{
    Pese a que son pocos los países que practican el libre comercio, muchos economistas aún apoya el libre comercio como una política deseable. Esta defensa se basa en tres líneas argumentales: primero, un argumento formal para las ganancias de eficiencia del libre comercio que es simplemente el análisis coste-beneficio de la ausencia de política comercial; segundo, muchos economistas creen que el libre comercio produce ganancias adicionales que van más allá de este análisis formal; finalmente, dadas las dificultades de trasladar los complejos análisis económicos a las políticas reales, incluso quienes no ven el libre comercio como la mejor política imaginable, lo contemplan como una regla práctica útil.
    
    Hay motivos, intelectualmente aceptables, para desviarse del libre comercio. Un argumento que es claramente válido, en principio, sostiene que los países pueden mejorar su relación de intercambio mediante aranceles óptimos e impuestos a la exportación. Sin embargo, este argumento no es demasiado importante en la práctica. Los países pequeños no pueden tener mucha influencia sobre sus precios de importación o exportación; por tanto, no utilizan aranceles u otras políticas para mejorar su relación de intercambio. Por su parte, los países grandes pueden modificar su relación de intercambio, pero con la imposición de aranceles incurren en el riesgo de perturbar acuerdos comerciales y provocar represalias.
    
    Otro motivo de desviación del libre comercio se basa en los fallos del mercado nacional. Si algún mercado nacional, como el mercado de trabajo, no funciona correctamente, la desviación del libre comercio puede, a veces, ayudar a reducir las consecuencias de este mal funcionamiento. La teoría del segundo óptimo sostiene que, si un mercado no funciona correctamente y el estado no interviene, se aleja del óptimo. Un arancel puede aumentar el bienestar si en la producción de un bien existe un beneficio marginal social adicional al excedente del productor.
    
    A pesar de que los fallos del mercado son frecuentes, el argumento del fallo de mercado no debería aplicarse con demasiada ligereza. En primer lugar, es un argumento para políticas nacionales, más que para políticas comerciales; los aranceles son siempre un método peor, un «segundo óptimo», para compensar el fallo del mercado, que siempre es mejor atajar en su raíz. Además, el fallo del mercado es dificil de analizar suficientemente como para estar seguro de la recomendación política adecuada. 
    
    En la práctica, la política comercial está dominada por consideraciones sobre la distribución de la renta. No existe una forma sencilla de hacer modelos sobre los aspectos políticos de la política comercial, pero se han propuesto varias ideas útiles. Los expertos en ciencias políticas argumentan a menudo que las políticas se determinan por competencia entre partidos políticos que intentan atraer el máximo número de votantes. En el caso más sencillo, esta práctica conduce a la adopción de políticas que sirven a los intereses del votante mediano. Sin embargo, este enfoque, aunque útil para reflexionar sobre muchas cuestiones, parece proporcionar predicciones poco realistas sobre las políticas comerciales, que típicamente favorecen el interés de grupos pequeños y concentrados, por encima del interés público en general. Los economistas y los expertos en ciencias políticas lo explican generalmente apelando al problema de la acción colectiva. Dado que los individuos no tienen muchos incentivos para actuar políticamente en nombre de los grupos a los que pertenecen, aquellos grupos que están bien organizados (normalmente grupos pequeños que se juegan mucho) pueden conseguir a menudo que se aprueben las políticas que promueven sus intereses a expensas de la mayoría. 
    
    Si la política comercial se realizara a partir de consideraciones puramente nacionales, sería dificil conseguir el progreso hacia la liberalización comercial. Sin embargo, los países i ndustriales han alcanzado reducciones arancelarias sustanciales a través de un proceso de negociación internacional. La negociación i nternacional apoya la causa de la reducción arancelaria por dos caminos: ayuda a ampliar la base de la liberalización comercial al otorgar a los exportadores un apoyo directo, y ayuda a los estados a evitar las guerras comerciales mutuamente desventajosas, que podrían ser provocadas por la falta de coordinación de las políticas internacionales.
    
    A pesar de algunos progresos logrados en los años treinta hacia la liberalización comercial mediante acuerdos bilaterales, desde la Segunda Guerra Mundial la coordinación internacional ha tenido lugar principalmente mediante acuerdos multilaterales bajo el auspicio del Acuerdo General de Aranceles y Comercio. El GATT, que comprende una burocracia y unconjunto de reglas de conducta, es la institución central del sistema comercial internacional. El acuerdo mundial más reciente del GATT creó también una nueva entidad, la Organización Mundial del Comercio (OMC), para supervisar y poner en práctica el acuerdo.
    
    Además de las diversas reducciones de aranceles que han tenido lugar a través de la negociación multilateral, algunos grupos de países han negociado acuerdos comerciales preferentes que reducen los aranceles entre sí, pero no con respecto al resto del mundo. En virtud del GATT se permiten dos tipos de acuerdos preferentes: las uniones aduaneras, en las que los miembros del acuerdo establecen aranceles exteriores comunes, y las áreas de libre comercio, en las que los países no impqnen aranceles a los productos de los demás partícipes en el acuerdo pero mantienen sus propios aranceles frente al mundo exterior. Ambos tipos de acuerdo tienen efectos ambiguos sobre el bienestar económico. Si la integración conduce a reemplazar la producción nacional de alto coste por importaciones procedentes de otro país de la unión (el caso de la creación de comercio) el país gana. Si, por el contrario, la integración lleva a sustituir importaciones baratas de fuera del área integrada por otras más caras del interior (el caso de la desviación de comercio), el país pierde.
}

Hemos visto el marco conceptual en el que nos movemos al hablar de Intrumentos de Política Comercial y hemos proporcionado algunas de las aportaciones teóricas más sólidad para el análisis de impacto de dichas políticas, tanto en los microfundamentos del mercado como en el bienestar agregado de la economía. Ahora bien, hay algo que no debe pasar desapercibido: desde cualquier óptica, la Política Comercial activista que pone impedimentos al libre comercio menoscaba el Bienestar general, tanto desde la perspectiva doméstica como desde la perspectiva internacional. Y, sin embargo, pocos países se han aproximado completamente al libre comercio.\footnote{Tal vez la ciudad de Hong Kong, que legalmente pertenece a China pero mantiene una política económica independiente, sea la única economía moderna si n aranceles o cuotas de importación. No obstante, desde los tiempos de Adam Smith, los economistas han defendido el libre comercio como un ideal por el que la política comercial debe luchar. Las razones de esta defensa no son tan simples como la idea en sí. En un primer nivel, los modelos teóricos sugieren que el libre comercio evitará las pérdidas de eficiencia asociadas a la protección. Muchos economistas creen que el li bre comercio produce ganancias adicionales además de la eliminación de distorsiones a la producción y al consumo. Finalmente, incluso entre l os economistas que creen que el libre comercio no es una política perfecta, muchos creen que, en general, esta opción es mejor que ninguna otra política que pueda aplicar un gobierno.
}

Por eso resulta imperativo aterrizar la abrumadora cantidad de aproximaciones teóricas, tanto positivas como normativas, a la realidad en la que vivimos. Esto es precisamente los que se conoce como Economía Política de la Política Comercial. 

\textbf{¿Qué queremos decir?} así que revisaremos brevemente algunas de sus aportaciones. 

\subsection{Argumentos liberales: En favor del libre comercio}
La eficiencia como justificación del libre comercio es, simplemente, la otra cara de la moneda del análisis coste-beneficio de un arancel. 

La referencia básica del caso del país pequeño que no puede influir sobre los precios de exportación extranjeros. Un arancel genera una pérdida neta en la economía, medida por el área de los dos triángulos; se produce por la distorsión de los incentivos económicos de productores y consumidores. Analogamente, un cambio hacia el libre comercio elimina estas distorsiones y aumenta el bienestar nacional.

En el mundo moderno, por razones que explicaremos más adelante en este capítulo, los tipos impositivos de los aranceles suelen ser reducidos y las cuotas a la importación, bastante poco frecuentes. Por ello, las estimaciones de los costes totales provocados por las distorsiones generadas por los aranceles y las cuotas a la importación tienden a ser reducidas. Para el conjunto del mundo, según estas estimaciones, los costes de la protección son inferiores al 1\% del PIB. Las ganancias del libre comercio son algo menores para las economías avanzadas, corno los Estados Unidos y Europa, y ligeramente superiores para los «países en desarrollo» más pobres.

Entre los economistas se ha extendido la creencia de que este tipo de cálculos, a pesar de que afirman que en algunos casos se pueden lograr sustanciales·ganancias gracias al libre comercio, no cuentan toda la historia. En los países pequeños, en general; y en los países en desarrollo, en particular, muchos economistas consideran que hay importantes ganancias asociadas al libre comercio que no se pueden contabilizar en un análisis convencional coste-beneficio. Los tres argumentos esbozados en el apartado anterior representan probablemente el punto de vista general de la mayor parte de los economistas especializados en economía internacional, al menos en los Estados Unidos:
\begin{lnum}
    \item Los costes convencionales de desviarse del libre comercio son elevados;
    \item Hay otros beneficios del libre comercio además del coste de las políticas proteccionistas; y,
    \item Cualquier intento de conseguir sofisticadas desviaciones del libre comercio será subvertido por el proceso político.
\end{lnum}

\subsubsection{Ganancias dinámicas: Economías deescala}
Un tipo de ganancias adicionales comprende las economías de escala.\footnote{Las ganancias adicionales del libre comercio que se analizan aquí se definen, a veces, como ganancias dinámicas, porque la mayor innovación y competencia puede tardar más en producirse que la eliminación de las distorsiones en la producción y el consumo.
} 

Los mercados protegidos limitan las ganancias de las economías de escala externas al inhibir la concentración de industrias; cuando las economías de escala son internas, no solamente fragmentan la producción internacional sino que, al reducir la competencia y aumentar los beneficios, también atraen a demasiadas empresas a la industria protegida. Con la proliferación de empresas en mercados nacionales de un tamaño reducido, la escala de producción de cada empresa se hace ineficiente.\footnote{Un buen ejemplo de cómo genera la protección una escala ineficiente es el caso de la industria del automóvil argentina, que surgió debido a la existencia de restricciones a la importación. Una fábrica de ensamblaje de escala eficiente debería fabricar de 80.000 a 200.000 automóviles al año; a pesar de eso, en 1964, la industria argentina, que producía solamente 166.000 coches, estaba formada por no menos de 13 empresas. Algunos economistas consideran que la necesidad de detener un número de entradas excesivas en una industria y las resultantes escalas ineficientes de producción son justificaciones del libre comercio que desbordan el cálculo estándar del análisis coste-beneficio.
} 

\paragraph{Ganancias dinámicas: Curva de aprendizaje}
Otro argumento en favor del libre comercio es que, al proporcionar a los empresarios un incentivo para buscar nuevas vías para exportar o competir con las importaciones, el libre comercio ofrece más oportunidades para el aprendizaje y la innovación que un sistema de comercio administrado, en el que el gobierno dicta en gran parte el patrón de importaciones y exportaciones. 

El capítulo 1 1 analiza las experiencias de los países menos desarrollados, que descubrieron oportunidades de exportación inesperadas cuando cambiaron sus sistemas de cuótas de importación y aranceles por políticas comerciales más abiertas. 

Otra forma relacionada de ganancias derivadas del libre comercio deriva de la tendencia a que las empresas más productivas exporten, mientras que las menos productivas se quedan en el mercado nacional. Esto sugiere que la transición al libre comercio hará que el conjunto de la economía sea más eficiente al pasar a tener una combinación de industrias con empresas con mayor productividad. Estos argumentos adicionales en favor del libre comercio son difíciles de cuantificar, aunque algunos economistas lo han intentado.\footnote{Por lo general, los modelos que intentan tener en cuenta las economías de escala y la competencia imperfecta ofrecen cifras superiores a la estimación de las pérdidas derivadas de la Política Comercial activista. Sin embargo, no hay consenso en torno a la mayor cuantía real de las ganancias del libre comercio. Si las ganancias adicionales son tan importantes como consideran algunos economistas, los costes de distorsionar el comercio con aranceles, cuotas, subsidios a la exportación, etcétera, serán por tanto muy superiores a las medidas solamente por el análisis costebeneficio  convencional.
}

\subsubsection{Ineficiencia productiva: Búsqueda de rentas}
Cuando las importaciones se limitan con una cuota en vez de un arancel, el coste se ve amplificado en ocasiones por un proceso conocido como búsqueda de rentas. Para aplicar una cuota a las importaciones, el gobierno tiene que emitir licencias de importación, con lo que quienquiera que las reciba obtendrá unas rentas económicas por tenerlas. En algunos casos, los individuos y las empresas incurren en costes sustanciales (con lo que se desperdician, en la práctica, algunos de los recursos productivos de la economía) en un esfuerzo por conseguir licencias de importación. 

Un conocido ejemplo ocurrió en la India en las décadas de 1950 y 1960. En aquella época, las empresas indias recibían el derecho de comprar bienes importados en función de su capacidad instalada. Se creó así un incentivo para invertir en exceso: por ejemplo, una acería podría construir más hornos de los que consideraba que necesitaría porque le otorgaría un mayor número de licencias de importación.\footnote{Un ejemplo más reciente e inhabitual de la búsqueda de rentas está relacionado con las importaciones estadounidenses de atún enlatado. El atún está protegido por una cuota arancelaria: se puede importar una pequeña cantidad de atún (el 4,8\% del consumo estadounidense) a una tasa arancelaria reducida, el 6\%, aunque las importaciones adicionales tienen que pagar una tasa arancelaria del 12,5\%. Por alguna razón, no hay licencias de importación; cada año, el derecho de importar atún a un tipo arancelario reducido se asigna por riguroso orden de solicitud. El resultado es una costosa carrera para Hevar el atún a los Estados Unidos lo antes posible. He aquí cómo describe el proceso de búsqueda de rentas la Comisión del Comercio Internacional de los Estados Unidos: Los importadores intentan obtener la mayor proporción de la cuota arancelaria a su alcance mediante la acumulación de atún enlatado en almacenes en las aduanas a finales de diciembre para sacarlo de los almacenes al comienzo del año natural. El dinero que gastan los importadores en almacenar ingentes cantidades de atún en diciembre representa un coste adicional de la protección para la economía estadounidense.
}

\subsubsection{Proceso Político}
Un argumento político en favor del libre comercio refleja el hecho de que un compromiso político con el libre comercio puede ser una buena idea en la práctica, incluso a pesar de que puedan existir, en principio, mejores políticas. Los economistas consideran a menudo que las políticas comerciales, en la práctica, están dominadas por intereses políticos especiales más que por una comparación de los costes y beneficios nacionales. 

A veces, los economistas pueden demostrar que, en teoría, una determinada selección de aranceles y subsidios de exportación puede aumentar el bienestar nacional pero, en realidad, cualquier organismo estatal que intente conseguir un sofisticado programa de intervención comercial puede caer presa de grupos de interés y verse convertido en un aparato para redistribuir la renta en favor de sectores con influencia política. Si este argumento es correcto, puede ser preferible defender el libre comercio sin excepciones, incluso a pesar de que en el terreno puramente económico el libre comercio puede no ser siempre la mejor política concebible.

\subsection{Argumentos inervencionistas: Política Comercial Activista}
No obstante, cabe detacar también que no hay nada, hoy por hoy, que descarte de una vez por todas la deseabilidad de la intervención gubernamental en el comercio. De hecho, la Economía Internacional muestra queexisten tanto ganancias derivadas del libre comercio, como ganancias derivadas de la intervención gubernamental mediante el recurso a la Política Comercial.

Ahora que conocemos algunos de los argumentos en pro de la liberalización comercial, revisamos brevemente algunos otros en favor de la intervención.

\subsubsection{Política Comercial óptima}
Cuando hemos los dos marcos analíticos más general sobre los cuales analizar la incidencia de la Política Comercial, hemos visto cómo la Política Comercial influye de manera determinante en un amplio conjunto de medidas de desempeño industrial (productividad, tamaño, precios, markups). Hemos observado también que la reducción de barreras comerciales tiene fuertes efectos sobre el bienestar agregado: la apertura (cierre) comercial provoca cambios que mejoran (empeoran) el bienestar de los consumidores en su conjuneto, pero también genera perdedores (ganadores) domésticos en el mercado deproducto que no deben ser minimizados (maximizados). 

Además, en nuestro análisis inicial hemos visto que tendrá efectos muy diferentes en función de: (i) el tamaño de mercado los mercados implicados (ganadores y perdedores); (ii) las disparidades tecnológicas entre países (ventaja comparativa); y, (iii) la distorsión generada sobre las decisiones de localización/entrada (desplazando producción entre países). 

En definitiva, hemos visto que incide a través de diversos canales que podríamos sintetizar en tres:
\begin{lnum}
    \item \item \textbf{Canal de bienestar}. Un aumento de barreras comerciales reduce las importaciones, lo que implica menor acceso a variedades extranjeras y, por tanto, una pérdida de eficiencia estática (menor bienestar del consumidor);
    \item \textbf{Términos de intercambio}. Los aranceles permiten al país grande ejercer poder de mercado, forzando a los exportadores extranjeros a reducir sus precios netos para seguir accediendo al mercado doméstico. Esto genera una transferencia de renta desde el exterior hacia el país que impone el arancel; y,
    \item \textbf{Canal renta (Balanza Comercial)}. La política comercial afecta al equilibrio externo y, por tanto, a la renta nacional a través de: (i) cambios en salarios (endógenos en ARA); y, (ii) recaudación arancelaria. En particular, los aranceles reducen importaciones, alteran la balanza comercial y generan ingresos públicos que se redistribuyen a los consumidores, afectando directamente al bienestar. Además, el ajuste de la balanza comercial induce cambios en los salarios que afectan costes, entrada/salida de empresas y selección.
\end{lnum}

Por tanto, si la Política Comercial incide a través de diferentes canales de maneras diferentes, ¿podemos minimizar las pérdidas maximizando las ganancias? La respuesta a esto es precisamente lo que se conoce como Política Comercial óptima. En este contexto, la Política Comercial óptima (o arancel óptimo) será exactamente el punto en el que estos dos efectos se anulan:\footnote{La elasticidad de equilibrio, evaluada en el nivel de comercio que resulta sin el arancel (o en su entorno): la elasticidad de la función de demanda de importaciones que genera el modelo en equilibrio general.
}

\eqblock{
\begin{aligned}
    \boxed{t^* = \frac{1}{\varepsilon_M}}\qquad \text{donde,}
    \left\{\begin{aligned}
        &\text{TOT}=\frac{p_X^h}{p_M^h},&&p_M^h=\frac{p_{\max}^b}{\tau}\quad\Longrightarrow\quad\substack{\text{Mejora nuestra relación}\\\text{ de intercambio}}\\[1em]
        &\downarrow M^h\Longrightarrow\,\,\downarrow N^h,&&\text{Empeora el Bienestar agregado}
    \end{aligned}\right.
\end{aligned}
}{Política Comercial óptima: Marco MELITZ-OTTAVIANO. Política Comercial: Instrumentos y efectos}

Donde $\varepsilon_M$ es la elasticidad de la demanda de importaciones. Es decir, cuanto menos elástica ($\varepsilon_M\to0$) sea la demanda externa, mayor poder de mercado tiene el país y mayor es el arancel óptimo. ¿Por qué? Porque un país pequeño puede sustituir facilmente entre muchos proveedores extranjeros, no tiene influencia sobre el precio mundial y los exportadores extranjeros (con alta productividad) no tienen que modificar sus precios para introducirse en el mercado. En consecuencia, países pequeños tendrán una elasticidad mucho mayor que la de mercados grandes. En esencia, la Política Comercial óptima \textbf{será aquella que mejore el bienestar a través de la relación de intercambio minimizando los efectos por menor variedad y menor competencia}. 


\subsubsection{Externalidades: Tecnología}
Por último, uno de los argumentos más utilizados en favor de la Política Comercial activista (en los países avanzados) es el de la generación de externalidades positivas al desarrollarse una determinada industria.

Si las empresas en una industria generan conocimientos que otras empresas también pueden utilizar sin pagar, la industria genera, en la práctica, una producción adicional (el beneficio marginal social del conocimiento) que no se refleja en los incentivos de las empresas. Cuando se puede demostrar que esas externalidades (beneficios que se acumulan en el exterior de las empresas) pueden ser importantes, hay una buena razón para subsidiar a la industria. 

En los países avanzados el argumento adquiere un carácter especial porque en estas naciones existen importantes industrias de alta tecnología en las que la generación de conocimientos es, en muchos casos, el aspecto central de la empresa. En las industrias de alta tecnología, las empresas dedican gran parte de sus recursos a mejorar la tecnología, mediante gastos en investigación y desarrollo, o con su disposición a padecer pérdidas iniciales en nuevos productos y procesos para ganar experiencia. Estas actividades tienen lugar en casi todas las industrias, por supuesto, por lo que no existe una separación nítidamente definida entre la alta tecnología y el resto de la economía. Sin embargo, se aprecian claras diferencias de grado y tiene sentido hablar de un sector de alta tecnología en el que la inversión en conocimientos es el aspecto clave del negocio.

La razón de ser de la política comercial activista es que, mientras las empresas pueden apropiarse de algunos beneficios de su propia inversión en conocimientos (de otro modo no invertirían en esa actividad), normalmente no pueden adueñarse de ellos por completo. Algunos de los beneficios se obtienen en otras empresas que pueden emular las ideas y las técnicas de las compañías punteras. En electrónica, por ejemplo, es frecuente que las empresas lleven a cabo una ingeniería inversa de los diseños de sus rivales, un proceso en virtud del cual toman los productos de la competencia para comprender cómo funcionan y cómo están producidos. Dado que las leyes de patentes proporcionan sol o una débil protección a los innovadores, puede presumirse de forma razonable que en un ámbito de laissez-faire, las empresas de alta tecnología no reciben los incentivos que deberían para proseguir en su tarea de innovación.

\subsubsection{Trilema de RODRIK}
Por último, pero no menos importante, es necesario mencionar las aportaciones del economista de Harvard RODRICK (2009). El autor, que a lo largo de su carrera se ha especializado en la Economía Internacional no solo desde un punto de vista analítico sino también descriptivo (más coherente con la Economía Política) ha sido un ferviente crítico de las tendencias globalizadoras que han ido surgiendo desde la Ronda de Uruguay (catalogándolas en muchas ocasiones como: anti-librecomercio).

Una de sus aportaciones más citadas es el conocido Trilema de RODRICK, que plantea una coyuntura tripartita frente al contexto de globalización, según la cual las sociedades (o estados) no pueden optar al mismo tiempo por: (i) sostener altos niveles de globalización económica; (ii) soberanía nacional; y, (iii) democracia.

\imagenfit{TrilemaRodrick.png}{Trilema de RODRICK}{\href{https://www.realinstitutoelcano.org/comentarios/la-voladura-del-trilema-de-rodrik/#:~:text=En\%20su\%20libro\%20La\%20paradoja\%20de\%20la,soberanía\%20nacional\%20y\%20la\%20democracia\%2C\%20sino\%20sólo}{\textit{La voladura del trilema de RODRICK}. Real Instituto El Cano (2018)}}

Según Rodrik, al intentar elevar los niveles de dos de cualquiera de los tres factores mencionados, el tercero restante tiende a resentirse y, por tanto, verse afectado negativamente; situación que crea un trilema. Como solución, Rodrik plantea que los Estados deben, necesariamente, renunciar a una de las alternativas

































\clearpage
\section{Conclusión} 

\textcolor{red}{\textbf{OJO!}} Puede ser interesante concluir ligando con el Tema siguiente: En conclusión, hemos estudiado los Instrumentos tradicionales de la Política Comercial. Quizás puedan categorizarse como la forma más brusca de intervenir y hacer bascular la balanza a nuestro favor. También hemos visto como sus efectos son díficilmente cuantificables pues la interrelación entre multitud de mercados y hace práctiamente imposible su análisis aislado. 

Por esta misma razón (y, evidentemente, por las restricciones impuestas por la OMC), la Política Comercial actual se caracteriza por su sutileza y dirección. Esta nueva forma de hacer Política Comercial dio origen, en un primer momento, a la Política Comercial estratégica (aunque sería cuestionable no tratarla como un instrumento comercial convencional) y, más en particular, a la Política de Promoción Exterior.




\subsection{Recapitulación}


\subsection{Extensiones y relación con otras partes del Temario}



\subsection{Opinión}

\subsubsection{Idea final (Salida o cierra)}

\clearpage
\pagenumbering{gobble}
\MyBibliography



\MyBibItem{DAWID y GATTI}{Chapter 2 - Agent-Based Macroeconomics}{2018}{https://doi.org/10.1016/bs.hescom.2018.02.006}


% ANEXOS
\clearpage
\anexo{\textbf{Preguntas Test}}
%\input{\TemaCode/questions.tex} 



\anexo{Modelo MELITZ-OTTAVIANO: Comercio intrasectorial y aranceles}\label{anx:melitz_aranceles}
Consideremos una economía con $L$ consumidores, cada uno de ellos suministrando una unidad de trabajo. 

$$\text{\textbf{2.1. (Economía cerrada)} Precios y demanda}$$
Las preferencias se definen sobre un continuo de variedades diferenciadas, indexadas por $i \in I$, y sobre un bien homogéneo que se toma como numerario. Todos los consumidores comparten la misma función de utilidad, dada por:

$$U = m + \alpha \int_{i \in I} x_i^d \, \mathrm di \;-\; \frac{\gamma}{2} \int_{i \in I} (x_i^d)^2 \, \mathrm di \;-\; \frac{\eta}{2} \left( \int_{i \in I} x_i^d \, \mathrm di \right)^2,\quad \text{(1)}$$

Donde $m$ y $x_i^d$ representan, respectivamente, los niveles de consumo individual del bien numerario y de cada variedad $i$. Los parámetros de demanda $\alpha$, $\eta$ y $\gamma$ son todos positivos. Los parámetros $\alpha$ y $\eta$ determinan el patrón de sustitución entre las variedades diferenciadas y el bien numerario: aumentos en $\alpha$ y disminuciones en $\eta$ desplazan la demanda hacia las variedades diferenciadas en relación con el numerario. El parámetro $\gamma$ mide el grado de diferenciación entre variedades. En el límite en que $\gamma = 0$, los bienes son \textbf{sustitutivos perfectos}.\footnote{Por tanto, los consumidores solo valoran su consumo total agregado de variedades, $X = \int_{i \in I} x_i^d \, \mathrm di$, por lo que la estructura competitiva del mercado es cercana a la competencia perfecta. A medida que $\gamma$ aumenta, también lo hace el grado de diferenciación, ya que los consumidores otorgan mayor peso a la distribución del consumo entre variedades.
} 

Como las utilidades marginales de todos los bienes están acotadas, un consumidor puede no demandar alguna variedad en particular. Suponiendo que los consumidores siempre demandan una cantidad positiva del bien numerario, es decir, $m_0 > 0$. La demanda inversa de cada variedad consumida $i$ viene entonces dada por:

$$p_i = \alpha - \gamma x_i^d -\eta X\quad \text{(2)}$$,

Sea $I^* \subseteq I$ el subconjunto de variedades efectivamente consumidas (aquellas para las que $x_i > 0$), la expresión anterior puede invertirse para obtener el sistema de demanda lineal de mercado para estas variedades:

$$\begin{aligned}
    &X_i^d \equiv L x_i^d\Longrightarrow \underset{\text{\scriptsize (utilizar en exposición oral)}}{\boxed{X_i^d=\frac{L}{\gamma}\left[\alpha-\eta X-p_i\right]}}\\[1em]
    &X=\int x_i\mathrm di\underset{\substack{\text{\scriptsize sustituyendo $X$}\\\text{\scriptsize (endógeno)}}}{\Longrightarrow}\quad X_i^d=\frac{\alpha L}{\eta N + \gamma} - \frac{L}{\gamma} p_i + \frac{\eta N}{\eta N + \gamma} \frac{L}{\gamma} \bar{p}, \quad \forall i \in I^*\qquad \text{(3)}
\end{aligned}$$,

Donde $N$ es la medida de variedades consumidas en $I^*\in I$ y $\bar {p} = \frac{1}{N} \int_{i \in I^*} p_i \, \mathrm di$ es su precio medio. El conjunto $I^*$ es el mayor subconjunto de $I$ tal que se cumple:

$$p_i \leq \frac{1}{\eta N + \gamma} (\gamma \alpha + \eta N \bar{p}) \equiv p_{\max}, \qquad \text{(4)}$$.

Donde el término del lado derecho, $p_{\max}$, representa el precio al cual la demanda de una variedad se reduce a cero (\textbf{precio máximo compatible con demanda positiva para una variedad dada}). 

OJO! La ecuación de demanda lineal para cada variedad implica que $p_{\max} \leq \alpha$. Por tanto, a diferencia del caso de demanda C.E.S., la elasticidad-precio de la demanda no viene determinada únicamente por el grado de diferenciación del producto $\gamma$. Dado este último, una reducción del precio medio $\bar{p}$ o un aumento en el número de variedades competidoras $N$ (de variedades efectivamente consumidas) provocan una disminución de $p_{\max}$ y un aumento de la elasticidad-precio de la demanda $\varepsilon_i$ para cualquier $p_i$ dado. Por tanto, esto se interpreta como un \textbf{entorno competitivo más duro} y, por tanto, dado dicho entorno competitivo (esto es, dados $\Omega$ y $\mathbf{p}$), la elasticidad $\varepsilon_i$ aumenta monótonamente con el precio $p_i$ a lo largo de la curva de demanda.\footnote{La elasticidad precio de la demanda viene dado por: $\varepsilon_i = \left| \frac{\partial q_i}{\partial p_i} \cdot \frac{p_i}{q_i} \right| = \left[\left(\frac{p_{\max}}{p_i}\right) - 1\right]^{-1}$.
} 

De la misma forma, a diferencia del caso CES-isoelástica donde el índice de precios se consideraba la función de utilidad indirecta (y por tanto un \textit{deflactor} de bienestar), bajo esta especificación el bienestar puede evaluarse mediante la función de utilidad indirecta dada por:\footnote{Para garantizar niveles de demanda positivos del bien numerario, se supone que $m_0 > \int_{i \in I^*} p_i x_i^c \, \mathrm di = \mathbf{p} X^c - \frac{\Omega \sigma_p^2}{\gamma}$.
}

$$U = m_0 + \frac{1}{2}\left(\eta +\frac{\gamma}{N}\right)^{-1}(\alpha - \bar{p})^2 + \frac{1}{2} \frac{N}{\gamma} \sigma_p^2,\qquad\text{donde},
\left\{\begin{aligned}
    &\sigma_p^2 = \frac{1}{N} \int_{i \in I^*} (p_i - \bar{p})^2 \, \mathrm di&&\equiv \text{Varianza precios}\\[1em]
    &m_0&&\equiv\text{Renta del consumidor}
\end{aligned}\right.$$

Por tanto, el bienestar aumenta:
\begin{la}
    \item \textbf{Disminución del precio medio}. Cuando disminuye el precio medio $\mathbf{p}$, el consumidor podrá consumir más del bien \textit{variedad} (como conjunto) y, por tanto, aumentará su bienestar;
    \item \textbf{Varianza de los precios}. También aumenta cuando se incrementa la varianza de los precios $\sigma_p^2$ (manteniendo constante $\mathbf{p}$), ya que los consumidores pueden reasignar su gasto hacia las variedades más baratas, así como hacia el bien numerario; y,
    \item \textbf{Amor por la variedad}. Finalmente, el sistema de demanda presenta amor por la variedad: manteniendo constante la distribución de precios (es decir, manteniendo constantes tanto $\mathbf{p}$ como $\sigma_p^2$), el bienestar aumenta con el número de variedades $\Omega$ (efectivamente consumida).
\end{la}

$$\text{\textbf{2.2} Producción y Decisiones empresariales}$$
En cuanto a la producción y el comportamiento de las empresas, el trabajo es el único factor productivo y se ofrece inelásticamente en un mercado competitivo. El bien numerario se produce bajo rendimientos constantes a escala y con coste unitario, y su mercado es también competitivo. Estos supuestos implican un \textbf{salario unitario}, $w=1$. 

Por otro lado, la entrada en el sector de bienes diferenciados es costosa, ya que cada empresa incurre en un coste de desarrollo del producto y de puesta en marcha de la producción, $F$. Una vez establecida, la producción presenta rendimientos constantes a escala con un coste marginal $c$, que coincide con el requerimiento unitario de trabajo.\footnote{En este marco, con utilidad marginal acotada, las empresas con costes suficientemente elevados no sobreviven incluso en ausencia de dichos costes fijos.
} 

Este proceso se modeliza como una extracción de una distribución común y conocida $C(c)$, con soporte en el intervalo $[0, c_M]$. Dado que el coste de entrada es hundido, $F$, las empresas que pueden opera cubren su coste marginal y maximizan sus beneficios utilizando la función de demanda residual. Al hacerlo, dado el continuo de competidores, cada empresa toma como dados el nivel medio de precios $\mathbf{p}$ y el número de empresas $\Omega$ (efectivamente productoras). De esta forma, el precio $p(c)$ y el nivel de producción $x_i^s(c)$ que maximizan los beneficios de una empresa con coste $c$ deben entonces satisfacer:

$x_i^s(c) = \frac{L}{\gamma}\,[p_i(c) - c]$.

Si notamos como $c_D$ el coste de la empresa que es indiferente entre permanecer o salir de la industria, que obtiene beneficios nulos: su precio se reduce hasta igualar su coste marginal, $p(c_D) = c = p_{\max}$, y su nivel de demanda se reduce a cero, $x_i^d(c_D)$.\footnote{
El precio que maximiza beneficios, $p_i(c)$, podría situarse por encima del precio máximo, $p_{\max}$, en cuyo caso la empresa abandona el mercado.

Se supone que $c_M$ es lo suficientemente elevado como para situarse por encima de $c_D$, de modo que algunas empresas con niveles de coste entre ambos valores abandonan la industria: $c_D<c_M$
} Todas las empresas con $c < c_D$ obtienen beneficios positivos y permanecen en el mercado. Por tanto, el umbral de costes $c_D$ resumirá los efectos del precio medio y del número de empresas sobre las \textbf{variables de desempeño de todas las firmas}, que pueden expresarse como:\footnote{
Sea $r_j(c) = p_i(c) x_i^s(c)$ el ingreso, $\Pi_j(c) = r_j(c) - x_i^s(c)c$ el beneficio operativo y $\mu_j(c) = p_i(c) - c$ el markup (en términos absolutos) de una empresa con costes marginales $c$. Todas estas magnitudes pueden escribirse únicamente como funciones de $c$ y de $c_D$
}

$$\begin{aligned}
    &p_i(c) = \frac{1}{2}(c_D + c),\\[1em]
    &\mu_j(c) = \frac{1}{2}(c_D - c),\qquad \text{(8)}\\[1em]
    &x_i^s(c) = \frac{L}{2\gamma}(c_D - c),\\[1em]
    &r_j(c) = \frac{L}{4\gamma}\left[\left(c_D\right)^2 - c^2\right],\\[1em]
    &\Pi_j(c) = \frac{L}{4\gamma}\left(c_D - c\right)^2\qquad \text{(11)}.
\end{aligned}$$

Como es de esperar, las empresas más eficientes, tienen menores costes marginales, fijan precios más bajos y obtienen mayores ingresos y beneficios que aquellas con niveles de eficiencia más bajos. No obstante, las empresas más eficientes no trasladan completamente su ventaja de costes a los consumidores en forma de precios más bajos: también fijan markups más elevados, tanto en términos absolutos como relativos, que las empresas con mayores costes.

$$\text{\textbf{2.3.} Libre entrada y equilibrio (condición de libre entrada)}$$
Antes de la entrada, el beneficio esperado de una empresa es $\mathbb E_t[\Pi_j (c)]=\int_0^{c_D} \Pi_j(c)\, \mathrm dC(c) - F$.\footnote{Si este beneficio fuese negativo, ninguna empresa entraría en la industria.
} Por la condición de libre entrada y la existencia de un número infinito de potenciales competidores: el beneficio esperado se reduce a cero. Ahora, de estas dos condiciones (incertidumbre, ex ante, y realización, ex post) obtendremos: (i) el número de empresas que permanecen ($N$); y, (ii) el número de empresas que entran al mercado ($N_E$). En este contexto $N<N_E$, porque algunas empresas, al observar su coste decidirán salir del mercado pues no son lo suficientemente productivas ($c\in(c_D,c_M]$) 

Empezamos por definir la condición de equilibrio de libre entrada en base al beneficio esperado:

$$\begin{aligned}
    &\mathbb E_t[\Pi_j(c)]=\int_0^{c_D} \Pi_j(c)\, \mathrm dC(c) = \frac{L}{4\gamma} \int_0^{c_D} \left(c_D - c\right)^2 \, \mathrm dC(c) = F\\[1em]
    &\Longrightarrow\underset{\text{\scriptsize (exposición oral)}}{\boxed{\frac{L}{4\gamma}\int_0^{c_D}(c_D-c)^2\mathrm dC(c)=F}}
\end{aligned}$$,

Esta condición determina el umbral de eficiencia (o costes) $c_D$ efectivo del mercado. Es decir, el nivel mínimo de productividad capaz de cumplir la condición de beneficio nulo. Por tanto, este umbral determina el número de empresas supervivientes ya que, el umbral de beneficio nulo $c_D = p_i(c_D)$ debe coincidir, por construcción, con el precio máximo compatible con demanda positiva:

$$c_D = p_i(c_D) = \frac{1}{\eta N + \gamma}(\gamma \alpha + \eta N \bar{p})\leq p_{\max}$$.

Operando, obtenemos que el nùmero de empresas que efectivamente cumplan la condición de equilibrio con libre entrada (beneficio nulo), será:

$$\begin{aligned}
    &\underset{\substack{\\[0.5em]\text{\scriptsize Umbral de beneficio nulo}\\\text{\scriptsize (condición de libre entrada)}}}{N = \frac{2\gamma}{\eta} \cdot \frac{\alpha - c_D}{c_D-\bar{c}}},\qquad \underset{\text{\scriptsize número de empresas entrantes (ex ante)}}{N_E = \frac{N}{C(c_D)}}\\[1em]
    &\text{Tal que}:\qquad N\leq N_E
\end{aligned}$$

Donde $\bar{c} =\frac{\int_0^{c_D}c\, \mathrm dC(c)}{C(c_D)}$ es el coste medio de las empresas que sobreviven.\footnote{Está ya verificado que $\bar p=\frac{c_D+\bar c}{2}$.
} De esta forma vemos que, dada una tecnología de producción caracterizada por la función de distribución acumulada (CDF) $C(c)\in[0,1]$, la productividad media será mayor (es decir, $\bar{c}$ será menor) cuando: (i) los costes hundidos sean más bajos; (ii) cuando las variedades sean más sustitutivas (menor $\gamma$); y, (iii) en mercados de mayor tamaño (mayor número de consumidores $L$). 

Esto es así por todas las relaciones implícitas entre las tres ecuaciones anterior, de forma sintetizada:\footnote{El último paso sale de combinar algebraicamente la ecuación del precio de corte con la expresión del precio medio. No sale directamente de la condición de libre entrada, sino de usar el hecho de que el umbral $c_D$ debe coincidir con el precio máximo compatible con demanda positiva: primero obtenemos la condición de libre entrada, que determina $c_D$, y después dice que ese mismo $c_D$ igualamos al threshold de precio para demanda positiva, de donde sale la condición de beneficio nulo. Algebraicamente: $c_D=\frac{\gamma \alpha+\eta N\bar p}{\eta N+\gamma}$, desarrollamos ambos lados (multiplicamos $\eta N+\gamma$): $c_D(\eta N+\gamma)=\gamma\alpha+\eta N\bar p$.

Después, expandimos la expresión $\eta N c_D+\gamma c_D=\gamma\alpha+\eta N\bar p$ y agrupamos los términos con $N$ y depejamos: $\eta N(c_D-\bar p)=\gamma(\alpha-c_D)\quad\to\quad N=\frac{\gamma}{\eta}\frac{\alpha-c_D}{c_D-\bar p}$. Hasta aquí la fórmula intermedia correcta. 

Ahora usamos el resultado $\bar p=\frac{c_D+\bar c}{2}$, donde $\bar c$ es el coste medio de las firmas supervivientes. Entonces $c_D-\bar p=c_D-\frac{c_D+\bar c}{2}=\frac{c_D-\bar c}{2}$. Sustituyendo esto en la expresión anterior: $N=\frac{\gamma}{\eta}\frac{\alpha-c_D}{(c_D-\bar c)/2}=\frac{2\gamma}{\eta}\frac{\alpha-c_D}{c_D-\bar c}.$

La idea económica es que el número de firmas activas $N$ no se obtiene de integrar beneficios, sino de imponer que la firma marginal tiene precio $p(c_D)=c_D$ y que ese precio coincide con el precio de cierre (threshold) del mercado. La libre entrada fija el umbral $c_D$; la condición de demanda en el umbral traduce ese $c_D$ en un número de firmas $N$.
}

$$\begin{aligned}
    &\text(1)\,\,\int_0^{c_D} \Pi_j(c)\, \mathrm dC(c) - F=0\quad\Longrightarrow\quad\text{(2)}\,\, \int_0^{c_D} \Pi_j(c)\, \mathrm dC(c) = F\\[1em]
    &\hspace{10em}\Downarrow \text{\scriptsize Utilizando: $\Pi_j(c)=\frac{L}{4\gamma}(c_D-c)^2$}\\[1em]
    &\text{(3)}\,\,\frac{L}{4\gamma}\int_0^{c_D}(c_D-c)^2\, \mathrm dC(c) = F\quad\underset{\substack{\text{\scriptsize (Utilizando: )}\\\text{\scriptsize $p_i(c_D)=c_D$ y $p_i\leq p_{\max}$}}}{\Longrightarrow}\quad\text{(4)}\,\,N=\frac{2\gamma}{\eta}\frac{\alpha-c_D}{c_D-\bar c}
\end{aligned}$$

En todos los casos: (i) la tasa de salida de empresas es mayor que la probabilidad de supervivencia previa a la entrada, $C(c_D)$, que es menor; (ii) los parámetros de demanda $\alpha$ y $\eta$, que determinan el nivel agregado de demanda de las variedades diferenciadas, no afectan a la selección de empresas ni a la productividad del sector, sino únicamente al número de empresas en equilibrio; y, (iii) la competencia es más dura en mercados más grandes, ya que compiten más empresas y los precios medios $\bar{p} = (c_D + \bar{c})/2$ son más bajos. 

Una empresa con coste $c$ responde a este entorno más competitivo fijando un menor markup (en comparación con el que fijaría en un mercado más pequeño, como se desprende de la ecuación del markup).

$$\text{\textbf{2.4.} Parametrización de la Tecnología}$$
Todos los resultados derivados hasta ahora se mantienen para cualquier distribución de costes $C(c)$. En la modelización a la MELITZ se suele suponer una Función de distribución de PARETO. 

\begin{specialblock}{Distribución de PARETO}
    ¿Qué significa que la productividad siga una distribución de Pareto? Elegir una distribución de PARETO significa asumir que hay muchas empresas poco productivas y pocas muy productivas, y que la cola derecha (las más productivas) decrece de forma lenta: se trata de una distribución de \textbf{colas gruesas}. Su razón de ser es esencialmente el ajuste empírico: los datos de empresas muestran precisamente esa estructura, con unas pocas firmas muy eficientes dominando el mercado. 

    Formalmente, si defines los costes marginales como función de la productividad creamos dos relaciones: $c = 1/\varphi$: cuánta mayor productividad, menores costes marginales, y viceversa. De esta forma, una distribución de PARETO implica que:
    
    $$\begin{aligned}
        &\text{Teniendo en cuenta que}: c=\frac{1}{\varphi}\quad\iff\quad \varphi=\frac{1}{c}\\[1em]
        &\underset{\text{\scriptsize la cola gruesa queda hacia la derecha: la cota inferior ($\varphi_M$) se situa a la iquierda}}{C(\varphi)=\Pr(\varphi_M\leq \varphi)=1-\Pr(\varphi\leq\varphi_M)=1-\left(\frac{\varphi_M}{\varphi}\right)^k}\\
        &\underset{\text{\scriptsize llevado a costes marginales}}{\Longrightarrow}\quad \underset{\text{\scriptsize la cola gruesa queda hacia la izquierda: la cota inferior ($c_M$) se situa a la derecha}}{C(c)=\left(\frac{c}{c_M}\right)^k}
    \end{aligned}$$
    
    El parámetro $k\geq1$ mide la dispersión: cuanto menor es $k$, más dispersión (colas más gruesas y, por tanto, mayor peso en firmas muy productivas); cuanto mayor es $k$, colas menos gruesas y, por tanto, más concentradas están las firmas en niveles altos de coste (menos heterogeneidad).\footnote{Del Gatto, Mion y Ottaviano (2006) estiman la distribución de la productividad total de los factores utilizando datos a nivel de empresa para un panel de 11 países de la UE y 18 sectores manufactureros. Encuentran que la distribución de Pareto proporciona un ajuste muy adecuado para la productividad empresarial en distintos sectores y países. El valor medio estimado de k es cercano a 2. Combes et al. (2007) amplían el modelo y consideran distribuciones generales de costes/productividad, mostrando que los principales resultados de estática comparada no dependen de esta elección de parametrización.
    } En el límite, cuando $k\to\infty$, la distribución se vuelve degenerada en $c_M$. Cualquier truncamiento de la distribución de costes por arriba mantiene la misma forma funcional y el mismo parámetro $k$. 
    
    El soporte de ambas distribuciones viene exógenamente determinado: $c \in [0, c_M]$. Siendo $c_M$ el peor coste posible (la peor tecnología). Y, de forma equivalente, $\varphi\in[\varphi_M,\infty)$. Para poder ver gráficamente las funciones de distribución de ambas variables, se puede consultar: \href{https://www.geogebra.org/calculator/bf8nttpe}{GeoGebra: CDF Costes -- En el cuadrante superior | CDF Productividad -- En el cuadrante inferior}
\end{specialblock}


Por tanto, si suponemos una distrinución de PARETO para $C(c)$, la distribución de productividad de las empresas supervivientes también será de PARETO con parámetro $k$, y la distribución truncada de costes vendrá dada por:

$$C_D(c) = \left( \frac{c}{c_D} \right)^k, \quad c \in [0, c_D]$$.

Bajo esta parametrización, el nivel de coste umbral $c_D$ determinado por la condición de libre entrada es:

$$c_D = \left[ \frac{2(k + 1)(k + 2)\gamma}{L} (c_M)^k F \right]^{\frac{1}{k+2}}$$,

Donde suponemos que $c_M > \sqrt{\frac{2(k + 1)(k + 2)\gamma F}{L}}$ para garantizar que $c_D < c_M$. De esta forma, el número de empresas supervivientes viene dado por:

$$N= \frac{2(k + 1)}{\eta} \cdot \frac{\alpha - c_D}{c_D}$$.

Esta parametrización también permite obtener expresiones simples para las medias de todas las variables de desempeño empresarial:\footnote{Todas las derivaciones se basan en el supuesto de que los consumidores demandan positivamente el bien numerario. Los consumidores obtienen toda su renta del trabajo: no existen beneficios empresariales redistribuidos, ya que los beneficios de la industria (netos de costes de entrada) son nulos. Por tanto, es necesario garantizar que cada consumidor gasta menos que su renta unitaria en variedades diferenciadas. El gasto por consumidor en variedades es $N\bar{r}/L = (\alpha - c_D)c_D (k + 1)/[\gamma (k + 2)]$. Una condición suficiente para que esto sea menor que 1 es $\alpha < \frac{2\gamma (k + 2)}{k + 1}$.
}

$$\begin{aligned}
    &\text{Siguiendo la siguiente fórmula}: \bar{z} = \frac{\int_0^{c_D} z(c)\, \mathrm dC(c)}{C(c_D)},\quad\text{donde $\bar z\equiv$ cualquier variable (medio)}\\[1em]
    &\text{P.e. costes medios}:\quad \bar c=\frac{\int_0^{c_D}c\,\mathrm dC(c)}{C(c_D)}\quad\underset{\substack{
    \text{\scriptsize donde,}\\
    \text{\scriptsize$\left\{\begin{aligned}
        dC(c)=c(c)\,\mathrm dc=\frac{k}{c_M^k}c^{k-1}\mathrm dc\\
        C(c_D)=\left(\frac{c_D}{c_M}\right)^k=\frac{c_D^k}{c_M^k}.
    \end{aligned}\right.$}}}{\Longrightarrow}\quad\bar{c} = \frac{k}{k + 1} c_D\\[1em]
    &p_i(c) = \frac{1}{2}(c_D + c)\quad \to\quad \bar{p} = \frac{2k + 1}{2k + 2} c_D,,\\[1em]
    &\mu_j(c) = \frac{1}{2}(c_D - c),\quad\to\quad \bar{\mu} = \frac{1}{2}\frac{1}{k + 1} c_D\\[1em]
    &x_i^s(c) = \frac{L}{2\gamma}(c_D - c)\quad \to\quad\bar{x}^s = \frac{L}{2\gamma} \frac{1}{k + 1} c_D = \frac{(k + 2)(c_M)^k}{(c_D)^{k+1}} F,\\[1em]
    &r_j(c) = \frac{L}{4\gamma}\left[\left(c_D\right)^2 - c^2\right]\quad\to\quad \bar{r} = \frac{L}{2\gamma} \frac{1}{k + 2} (c_D)^2 = \frac{(k + 1)(c_M)^k}{(c_D)^k}F,\\[1em]
    &\Pi_j(c) = \frac{L}{4\gamma}(c_D - c)^2\quad\to\quad \bar{\Pi} = F\frac{(c_M)^k}{(c_D)^k}.
\end{aligned}$$

Al igual que ocurre con el coste medio, $\bar{c}$, la media de cualquier variable de desempeño $z(c)$ viene dada aplicando el desarrollo presentado. 

Aunque el coste medio se calcula como la media no ponderada de los costes, proporciona un índice representativo de un conjunto más amplio de medidas inversas de productividad. La productividad media de las empresas, $1/c$, ya sea no ponderada, ponderada por ingresos $r(c)$ o por producción $x(c)$, es proporcional a $1/\bar{c}$ (y, por tanto, a $1/c_D$).\footnote{Además, las varianzas de todas las variables de desempeño pueden expresarse como funciones simples de la varianza de los costes.
} 

Dado que el nivel de corte $c_D$ resume completamente la distribución de precios y de todas las demás variables, también determina de forma única el bienestar, a partir de (5):

$$U = m_0 + \frac{1}{2} \left(\eta+\frac{\gamma}{N}\right)^{-1} (\alpha - \bar{p})^2 + \frac{1}{2} \frac{N}{\gamma} \sigma_p^2\qquad \Longrightarrow\quad U = 1 + \frac{1}{2\eta} (\alpha - c_D)\left[\alpha - \frac{k + 1}{k + 2} c_D\right]$$.

Por tanto, el \textbf{bienestar aumenta cuando disminuye el umbral $c_D$}, ya que esto induce un aumento en las variedades disponibles, $N$, y una reducción del precio medio $\bar{p}$ (estos efectos dominan el impacto negativo de la menor varianza de precios).\footnote{Esta medida de bienestar refleja la \textbf{reducción en el consumo del bien numerario necesaria para cubrir los costes de entrada mediante el uso de trabajo.}
}

Por tanto, el bienestar aumenta:
\begin{la}
    \item \textbf{Disminuye el umbral de productividad mínimo ($c_D$)}. Ya que induce un aumento en las variedades disponibles y una reducción del precio medio. Lo que también provoca que el consumidor pueda consumir más del bien \textit{variedad} (como conjunto); y,
    \item \textbf{Tamaño de mercado}. Como se ha señalado anteriormente, mercados de mayor tamaño implican una selección más intensa (menor $c_D$), lo que se traduce en mayor productividad media (menor $\bar{c}$) y menores precios medios.
\end{la}

Estos efectos dominan el impacto negativo de la menor varianza de precios ya que, recordemos que mayor varianza mejoraba el bienestar por reasignación de gasto hacia variedades más baratas, así como hacia el bien numerario. 

Por otro lado, bajo esta parametrización, el tamaño medio de las empresas (tanto en producción como en ventas) y los beneficios son mayores en mercados más grandes: \textbf{el efecto directo del tamaño de mercado domina el efecto indirecto derivado de menores precios y markups}. Del mismo modo, los markups medios son menores, ya que el efecto directo del aumento de la competencia sobre los markups individuales domina el efecto de selección hacia empresas más eficientes (que, en términos relativos, fijan mayores markups). Asimismo, los beneficios medios y las ventas aumentan en la misma proporción cuando crece el tamaño del mercado, de modo que la rentabilidad media de la industria $\bar{\Pi}/\bar{r}$ no varía con el tamaño del mercado. 

Finalmente, esta parametrización permite determinar de forma inequívoca los efectos del tamaño del mercado sobre la dispersión de las variables empresariales: la varianza de los costes, precios y markups es menor en mercados grandes (el efecto de selección reduce el soporte de estas distribuciones), mientras que la varianza del tamaño de las empresas (en términos de producción o ingresos) es mayor debido al efecto amplificador directo del tamaño del mercado.\footnote{Estos resultados de estática comparada son consistentes con la evidencia empírica para establecimientos en Estados Unidos, como documentan CAMPBELL y HOPENHAYN (2005) y SYVERSON (2004). Estos estudios muestran que en mercados más grandes las empresas presentan mayores niveles medios de ventas y productividad, menor dispersión en productividad y precios, y una selección más estricta que eleva el umbral inferior de productividad.
}


$$\text{\textbf{3. Economía abierta}}$$
Ahora que conocemos el modelo de economía cerrada para analizar los efectos del tamaño del mercado sobre distintas variables de desempeño a nivel sectorial.\footnote{Este modelo de economía cerrada puede aplicarse directamente a un conjunto de economías abiertas perfectamente integradas mediante el comercio. En ese caso, la transición desde la autarquía al libre comercio es equivalente a un aumento del tamaño del mercado que induce: (i) aumentos en la \textbf{productividad} media; (ii) aumentos en la \textbf{variedad} de productos; y, (iii) reducciones en los \textbf{markups} medios.
}

El escenario de economía cerrada no puede extenderse de manera inmediata al caso en el que los bienes no comercian libremente (barreras y costes de transporte). Por tanto, para comprender las interrelaciones entre mercados, se extiende el modelo a un contexto de dos países.\footnote{Este marco puede generalizarse a un número arbitrario de países y a patrones de costes comerciales también arbitrarios (incluyendo costes asimétricos).
}

Supongamos que existen: $A$ y $B$, con $L^a$ y $L^b$ consumidores en cada uno de ellos. Los consumidores de ambos países comparten las mismas preferencias, lo que conduce a la misma función de demanda inversa. Los mercados están segmentados, aunque las empresas pueden producir en un país y vender en el otro, incurriendo en un coste de comercio por unidad:\footnote{Se mostrará posteriormente que las condiciones de equilibrio excluyen cualquier oportunidad de arbitraje rentable. Por simplicidad, no se introducen costes fijos de exportación, ya que reducirían considerablemente la tractabilidad del modelo sin aportar nuevas intuiciones. En este marco, con utilidad marginal acotada, los costes por unidad son suficientes para generar selección en los mercados de exportación.
} el coste entregado de una unidad producida con coste $c$ en el país $h$ ($h = A, B$) es $\tau^h c$, donde $\tau^h > 1$. 

De esta forma, los países pueden diferir en dos dimensiones: el tamaño del mercado $L^h$ y las barreras a la importación $\tau^h$.

Sea $p_{\max}^h$ el precio máximo compatible con demanda positiva en el mercado $h$. Entonces:

$p^h = \frac{1}{\eta N^h + \gamma} (\gamma \alpha + \eta N^h \bar{p}^h)\leq p_{\max}^h, \quad h = A, B$,

Donde $N^h$ es el número total de empresas que venden en el país $h$ (tanto empresas domésticas como exportadores extranjeros) y $\bar{p}^h$ es el precio medio en ese país (considerando ambas):\footnote{Sea $p_D^h(c)$ y $x_D^h(c)$ el precio y la cantidad que maximizan beneficios en el mercado doméstico para una empresa localizada en $h$ con coste $c$. Esta empresa también puede decidir producir una cantidad $x_X^h(c)$ para exportar, que venderá en el otro país a un precio entregado $p_X^h(c)$.
}

$$\begin{aligned}
    &\text{Para una empresa del país $h$, las funciones objetivo serán}:\\[1em]
    &\text{Óptimo \textbf{doméstico}}:\left\{\begin{aligned}
        &p_D^h(c)\equiv \text{Precio $\max \Pi$ en mercado doméstico}\\[0.5em]
        &x_D^h(c)\equiv \text{Cantidad $\max \Pi$ en mercado doméstico}
    \end{aligned}\right.\\[1em]
    &\text{Óptimo \textbf{exportador}}:\left\{\begin{aligned}
        &p_X^h(c)\equiv \text{Precio $\max \Pi$ en mercado extranjero}\\[0.5em]
        &x_X^h(c)\equiv \text{Cantidad $\max \Pi$ en mercado extranjero}
    \end{aligned}\right.\\[1em]
\end{aligned}$$

Dado que los mercados están segmentados y las empresas operan con rendimientos constantes a escala, \textbf{maximizan de manera independiente} los beneficios obtenidos en ventas domésticas y en exportaciones. Sea: (i) el beneficio en el mercado doméstico, $\Pi_D^h(c) = (p_D^h(c)-c) x_D^h(c)$; y, (ii) el beneficio por exportaciones hacia el país $h$, $\Pi_X^h(c) = (p_X^h(c) - \tau^b c) x_X^h(c)$, con $h \neq b$.\footnote{En todo este análisis, cualquier derivación que involucre $h$ y $b$ es válida para $h = A, B$ y $h \neq b$.
} De forma análoga, los precios y cantidades que maximizan beneficios deben satisfacer:

$$x_D^h(c) = \frac{L^h}{\gamma} (p_D^h(c) - c); \quad x_X^h(c) = \frac{L^b}{\gamma} (p_X^h(c) - \tau^b c)$$.

\textbf{Al igual que en la economía cerrada, solo las empresas que obtienen beneficios no negativos en un mercado (doméstico o de exportación) decidirán operar en dicho mercado}. Esto da lugar a \textbf{reglas de corte en costes similares para ambos mercados}. Sea $c_D^h$ el umbral de coste para las empresas que venden en su mercado doméstico, y $c_X^h$ el umbral para las empresas exportadoras desde $h$ hacia $b$. Estos umbrales satisfacen:

$$\begin{aligned}
    &\underset{\text{\scriptsize corte empresa doméstica}}{c_D^h = \sup \{ c : \Pi_D^h(c) > 0 \} = p_{\max}^h,}\\[1em]
    &\underset{\text{\scriptsize corte empresa exportadora}}{c_X^h = \sup \{ c : \Pi_X^h(c) > 0 \} = \frac{p_{\max}^b}{\tau^b}}
\end{aligned}$$

Esto tiene una implicación clave: $c_X^b = \frac{c_D^h}{\tau^h}$. Es decir, las \textbf{barreras comerciales dificultan que las empresas exportadoras alcancen el punto de beneficio nulo en comparación con las empresas domésticas}. 

Como en la economía cerrada, estos umbrales de coste resumen todos los efectos de las condiciones de mercado relevantes para el desempeño empresarial. En particular, los precios y cantidades óptimos pueden expresarse como funciones de dichos umbrales:

$$\begin{aligned}
    &\text{Por las condiciones de corte}:\\[1em]
    &\left\{\begin{aligned}
        &p_D^h(c) = \frac{1}{2}(c_D^h + c),\\[1em]
        &p_X^h(c) = \frac{\tau^b}{2}(c_X^h + c),
    \end{aligned}\right.\\[3em]
    &\left\{\begin{aligned}
        &x_D^h(c) = \frac{L^h}{2\gamma}(c_D^h - c),\\[1em]
        &x_X^h(c) = \frac{L^b}{2\gamma}\tau^b (c_X^h - c),
    \end{aligned}\right.\\[3em]
    &\text{Lo que conduce a los siguientes beneficios máximos:}\\[1em]
    &\left\{\begin{aligned}
        &\Pi_D^o(c) = \frac{L^o}{4\gamma}(c_D^o - c)^2,\\[1em]
        &\Pi_X^o(c) = \frac{L^h}{4\gamma}\left(\tau^h\right)^2 (c_X^o - c)^2.
    \end{aligned}\right.
\end{aligned}$$


$$\text{\textbf{3.1.} Condición de Libre Entrada}$$
La entrada es libre en ambos países. Las empresas \textbf{eligen su localización de producción antes de entrar y de incurrir en el coste hundido de entrada}. Con el fin de centrar el análisis en los efectos del tamaño de mercado y de las diferencias en costes comerciales, se supone que ambos países comparten la misma tecnología, caracterizada por el coste de entrada $F$ y la distribución de costes $C(c)$.\footnote{Se relaja este supuesto en el apéndice, donde se analizan las implicaciones de la ventaja comparativa ricardiana.
} 

La libre entrada de empresas domésticas en el país $h$ implica que el beneficio esperado es nulo en equilibrio:\footnote{Se mantiene la misma parametrización de PARETO para la distribución de costes $C(c)$ en ambos países.
}

$$\int_0^{c_D^h} \Pi_D^h(c)\, \mathrm dC(c) + \int_0^{c_X^h} \Pi_X^h(c)\, \mathrm dC(c) = F$$.

Dada la condición de maximización de beneficios, la condición de libre entrada puede reescribirse como:

$$L^h \left(c_D^h\right)^{k+2} + L^b \left(\tau^b\right)^2 \left(c_X^h\right)^{k+2} = \gamma\Gamma$$,

Donde $\Gamma = 2(k + 1)(k + 2)(c_M)^k F$ es un \textbf{índice tecnológico}: recoge conjuntamente los efectos de una mejor distribución de costes (menor cota superior, $c_M$) y de menores costes de entrada $F$. Esta condición de libre entrada se cumple siempre que exista una masa positiva de entrantes domésticos $N_E^h > 0$ en el país $h$.\footnote{En caso contrario, $\int_0^{c_D^l} \Pi_D^o(c)\, \mathrm dC(\varphi) + \int_0^{c_X^o} \Pi_X^o(c)\, \mathrm dC(\varphi) < F$, lo que implica $I_E^o = 0$ y que el país $o$ se especializa en el bien numerario. Para simplificar, se descarta este caso suponiendo que $\Gamma$ es suficientemente grande.
} En este análisis se considera el caso en el que ambos países producen el bien diferenciado y $N_E^h > 0$ para $h = A, B$. Entonces, dado que $c_X^b = c_D^h/\tau^h$, la condición de libre entrada puede reescribirse como:

$$L^h \left(c_D^h\right)^{k+2} + L^b \rho^b \left(c_D^b\right)^{k+2} = \gamma\Gamma,\qquad\text{donde},\,\,\underset{\text{\scriptsize "grado de libertad" del comercio en $h$}}{\rho^h \equiv \left(\tau^h\right)^{-k} \in (0,1)}$$,

Vemos como $\rho^o$ es una medida inversa de los costes comerciales (el \textbf{grado de libertad del comercio}). Este sistema, para $h = A,B$, puede resolverse para obtener los umbrales de costes en ambos países:

$$c_D^h = \left[ \frac{\gamma\Gamma}{L^h} \cdot \frac{1 - \rho^b}{1 - \rho^h \rho^b} \right]^{\frac{1}{k+2}}\qquad \forall h,b\in\text{Paises}:\,h\neq b$$.

$$\text{\textbf{3.2.} Precios, variedad y Bienestar}$$
En cuanto a precios, variedad y bienestar, los precios en el país $h$ reflejan tanto los precios domésticos de las empresas locales, $p_D^h(c)$, como los precios de los exportadores procedentes de $b$, $p_X^b(c)$. 

Utilizando las relaciones de barreras de costes y relaciones en equilibrio, estos precios pueden escribirse como:

$$\begin{aligned}
    &p_D^h(c) = \frac{1}{2}(p_{\max}^h + c), \quad c \in [0, c_D^h],\\[1em]
    &p_X^b(c) = \frac{1}{2}(p_{\max}^h + \tau^b c), \quad c \in \left[0, \frac{c_D^h}{\tau^h}\right]
\end{aligned}$$,

Donde $p_{\max}^h$ es el precio umbral para obtener una demanda positiva en el mercado $h$. 

Además, los costes de las empresas domésticas $c \in [0, c_D^h]$ y los costes entregados de los exportadores $\tau^h c \in [0, c_D^h]$ tienen distribuciones idénticas en este intervalo, dadas por $C^h(c) = (c/c_D^h)^k$. Por tanto, las distribuciones de precios de empresas domésticas y exportadoras en el país $h$ también coinciden. Y el precio medio en el país $h$ viene así dado por:

$$\bar{p}^h = \frac{2k + 1}{2k + 2} c_D^h$$.

Combinando este resultado con el precio umbral para demanda positiva, se obtiene el número de empresas que venden en el país $h$:

$$\left.\begin{aligned}
    &\underset{\text{\scriptsize aislamos $N^h$ y sustituimos $\bar p^h$ por:}}{p^h = \frac{1}{\eta N^h + \gamma} (\gamma \alpha + \eta N^h \bar{p}^h)}\\[0.5em]
    &\bar{p}^h = \frac{2k + 1}{2k + 2} c_D^h
\end{aligned}\right\}\Longrightarrow N^h = \frac{2(k + 1)\gamma}{\eta} \cdot \frac{\alpha - c_D^h}{c_D^h}$$.

Nótese que los resultados para número de variedades de producto y precios medios son idénticos a los del caso de economía cerrada, lo cual se debe a la coincidencia de las distribuciones de precios de empresas domésticas y exportadoras en cada mercado. 

Por último, y por la lo anterior, el bienestar en el país $h$ puede escribirse de forma análoga al caso de una economía cerrada:

$$U^h = 1 + \frac{1}{2\eta} (\alpha - c_D^h)\left[\alpha - \frac{k + 1}{k + 2} c_D^h\right]$$.

Una vez más, el bienestar varía de forma monótona con el umbral de costes doméstico $c_D^h$, que resume los efectos dominantes de la variedad y de los precios medios.\footnote{Las condiciones previamente establecidas sobre los parámetros de demanda $\alpha, \eta$ y $k$ garantizan que $x_0^o > 0$, tal como se había supuesto.
}

$$\text{\textbf{3.3.} Número de entrantes, Productores y exportadores}$$
En cuanto al número de entrantes, productores y exportadores, el número total de empresas que venden en el país $h$ (y, por tanto, la variedad de productos) está compuesto por productores domésticos y exportadores procedentes de $b$. Dada una masa positiva de (potenciales) entrantes $N_E^h$ en ambos países, existen $C(c_D^h) N_E^h$ productores domésticos y $C(c_X^b) N_E^b$ exportadores que finalmente venden en $h$, cumpliéndose:

$$C(c_D^h) N_E^h + C(c_X^b) N_E^b = N^h$$.

Esta condición (para cada país) puede resolverse para obtener el número de entrantes en cada país:

$$N_E^h = \frac{(c_M)^k}{1 - \rho^h \rho^b} \left[ \frac{N^h}{\left(c_D^h\right)^k}-\rho^h\frac{N^b}{\left(c_D^b\right)^k}\right]=\frac{2\left(c_M\right)^k(k+1)\gamma}{\eta(1-\rho^h\rho^b)}\left[\frac{\alpha-c_D^h}{\left(c_D^h\right)^{k+1}}-\rho^h\frac{\alpha-c_D^b}{\left(c_D^b\right)^{k+1}}\right]$$,

Se puede demostrar que la ecuación anterior implica que: $c_X^h < c_D^h$, de modo que solo un subconjunto de empresas relativamente más productivas exporta. Las empresas con costes más elevados (entre $c\in \left[c_D^h,c_X^h\right)$) únicamente sirven al mercado doméstico. Por tanto, $C(c_D^h) N_E^h$ representa también el número total de empresas que producen en el país $h$, ya que ninguna empresa produce sin abastecer al menos su mercado doméstico.

$$\text{\textbf{3.4.} \textit{Dumping} recíproco y oportunidades de arbitraje}$$
En esta sección se analiza el fenómeno del dumping recíproco y su relación con las condiciones de arbitraje. BRANDER y KRUGMAN (1983) mostraron que el dumping recíproco surge necesariamente en equilibrios de comercio intraindustrial bajo competencia à la Cournot, cuando empresas representativas producen un único bien homogéneo en dos países. 

Por su parte, OTTAVIANO, TABUCHI y THISSE (2002) demuestran que el dumping también puede aparecer con bienes diferenciados bajo competencia monopolística con empresas representativas. El modelo extiende este resultado a un entorno con heterogeneidad en costes, donde las empresas producen variedades diferenciadas y enfrentan elasticidades residuales de demanda distintas.\footnote{En trabajos relacionados, HOLMES y STEVENS (2004) muestran que la hipótesis de menores markups en mercados no locales, junto con diferencias en costes de transporte entre sectores, puede explicar diferencias entre mercados en la distribución del tamaño de las empresas.
}

Dadas las funciones de precios óptimos:

$$\left.\begin{aligned}
    &\underset{\text{\scriptsize precio doméstico empresas en $h$}}{p_D^h(c) = \frac{1}{2}(c_D^h + c)}\\[0.5em]
    &\underset{\text{\scriptsize precio de producto exportado de empresas en $h$}}{p_X^h(c) = \frac{\tau^b}{2}(c_X^h + c)}\\[0.5em]
    &\underset{\text{\scriptsize diferencias de productividad}}{\text{y, que}: c_X^h < c_D^h}
\end{aligned}\right\}\qquad\Longrightarrow \quad \frac{p_X^h(c)}{\tau^b} < p_D^h(c),\quad \forall c\ge c_X^h$$  

Por tanto, todas las empresas exportadoras fijan precios netos de costes de transporte (F.O.B) estrictamente inferiores a sus precios en el mercado doméstico.\footnote{Como ya señalaba WEINSTEIN (1992), esto implica que el dumping no debe interpretarse como una práctica de precios predatorios. Además, como muestran OTTAVIANO, TABUCHI y THISSE (2002), el dumping no requiere competencia oligopolística ni interacciones estratégicas entre empresas, elementos ausentes en este modelo.
}

Entonces el comportamiento de dumping está estrechamente vinculado a las condiciones de arbitraje en la reventa de bienes entre mercados:\footnote{Tal como explican FEENSTRA et al. (2001)
} el hecho de que $p_X^h(c)/\tau^b < p_D^h(c)$ equivale a una condición de no arbitraje que impide la reventa rentable, por parte de un tercero, de un bien producido y vendido en el país $h$. Asimismo, esta condición también impide la reventa rentable en sentido inverso, es decir, la reexportación hacia el país de origen $b$ de un bien previamente exportado al país $h$.\footnote{Dado que:
$$p_X^b(c) = p_D^h(c\tau^h) > p_X^h(c\tau^h)/\tau_h = p_D^b(c\tau^h\tau^b)/\tau^b > p_D^b(c)/\tau^b$$
}

$$\text{\textbf{3.5.} Impacto del Comercio}$$
En cuanto al impacto del comercio, se ha señalado anteriormente que la distribución del coste entregado de los exportadores $\tau^h c$ hacia el país $h$ coincide con la distribución de costes de las empresas domésticas en ese país. De ello se deduce que también coinciden las distribuciones de precios, markups, producción, ingresos y beneficios entre empresas domésticas y exportadoras en cada mercado. 

En consecuencia, la distribución de todas estas variables en equilibrio de economía abierta es idéntica a la de una economía cerrada con el mismo umbral de costes $c_D$. Por tanto, el análisis del comercio puede centrarse en la determinación de dicho umbral, gobernado por:

$$c_D^h = \left[ \frac{\gamma\Gamma}{L^h} \cdot \frac{1 - \rho^b}{1 - \rho^h \rho^b} \right]^{\frac{1}{k+2}}\qquad \forall h,b\in\text{Paises}:\,h\neq b$$

Comparando esta condición con la correspondiente a la economía cerrada, se observa inmediatamente que el umbral de costes es menor en economía abierta: el comercio aumenta la productividad agregada al forzar la salida de las empresas menos productivas. 

Este resultado es análogo al de MELITZ (2003), pero surge a través de un \textbf{canal distinto}:
\begin{la}
    \item \textbf{MELITZ (2003) -- Factores}. En MELITZ, el comercio intensifica la competencia en el mercado de factores: las empresas más productivas expanden su producción para exportar, elevando los salarios reales y expulsando a las menos eficientes. En ese modelo, la competencia de importaciones no afecta a la reasignación debido a la estructura CES (la elasticidad de demanda es constante y exógena); y,
    \item \textbf{MELITZ y OTTAVIANO -- Mercado de bienes}. En cambio, en el presente modelo, el único canal relevante es la \textbf{competencia en el mercado de bienes}. La oferta de trabajo es perfectamente elástica, por lo que no hay efectos vía salarios. La competencia de importaciones incrementa la intensidad competitiva, elevando las elasticidades residuales de demanda y obligando a salir a las empresas menos productivas. 
\end{la}

Se trata de un efecto equivalente al de un aumento del tamaño de mercado en economía cerrada: desplaza hacia abajo la distribución de markups. Por tanto, aunque sobreviven empresas más productivas: (i) el markup medio disminuye; (ii) los precios caen debido al efecto conjunto de selección y competencia; (iii) aumentan el tamaño medio de las empresas, los beneficios y la variedad de productos (como en el caso de mercados más grandes); y, (iv) las ganancias de bienestar derivadas del comercio proceden de tres fuentes: aumento de la productividad (vía selección), reducción de markups (efecto procompetitivo) y mayor variedad.

$$\text{\textbf{3.6.} Efecto tamaño de mercado}$$
En equilibrio de economía abierta, las diferencias entre países en términos de desempeño empresarial vienen determinadas por las diferencias en los umbrales de costes $c_D^h$. En esencia, se trata de incluir el hecho de que, incluso con comercio, los mercados no están completamente integrados: el tamaño del país sigue siendo determinante. 

Si los costes comerciales son simétricos, el país más grande presenta un umbral de costes más bajo, mayor productividad media, mayor variedad, y menores precios y markups. En consecuencia, el bienestar es mayor en el país grande, que además atrae más empresas y productores locales. En este sentido, las diferencias asociadas al tamaño del mercado en autarquía persisten bajo comercio, aunque atenuadas.

De forma sorprendente, el tamaño del país socio no afecta al umbral de costes propio. Esto refleja la existencia de efectos compensatorios:
\begin{la}
    \item \textbf{Compensación vía exportaciones}. Por el lado de las exportaciones, un socio mayor implica mayores oportunidades de mercado, pero también mayor competencia (más empresas productivas reduciendo markups); y,
    \item \textbf{Compensación vía importaciones}. Por el lado de las importaciones, un socio mayor implica más competencia externa, pero en equilibrio esto se compensa con una menor entrada doméstica y, por tanto, menor competencia interna.\footnote{Como muestra la ecuación $N_E^h$, las diferencias en tamaño inducen una mayor proporción (potencial) de entrada en el país grande. Además, estas diferencias deben estar acotadas para evitar la especialización completa: si el tamaño relativo es extremo, el número de entrantes en el país pequeño tiende a cero.
    }
\end{la}


$$\text{\textbf{4. Liberalización comercial}}$$
En esta sección se analizan los efectos de la liberalización comercial. Se ha mostrado previamente que la decisión de localización de las empresas (determinada por la libre entrada en el largo plazo) desempeña un papel fundamental en la intensidad de la competencia entre mercados en una economía abierta. Esta decisión de localización también resulta crucial para entender las consecuencias de la liberalización comercial en el largo plazo, especialmente cuando las reducciones en las barreras comerciales son asimétricas.

$$\text{\textbf{4.1.} Liberalización bilateral (simétrica)}$$
En primer lugar, se considera el caso de liberalización bilateral simétrica. Se supone que los costes comerciales son iguales en ambos países, $\tau^A = \tau^B = \tau$, y se analizan los efectos de una reducción en $\tau$ (equivalentemente, un aumento en la libertad de comercio: $\rho^A=\rho^B=\Delta\rho$). 

En este caso, la condición de equilibrio para el umbral de costes puede escribirse como:

$$c_D^h = \left[\frac{\gamma\Gamma}{L^h (1 +\rho)}\right]^{\frac{1}{k+2}},\qquad
\text{donde $\rho$ mide la \textit{freeness of trade}}.$$

La liberalización bilateral aumenta la competencia en ambos mercados, lo que reduce proporcionalmente los umbrales de costes y, por tanto, incrementa proporcionalmente la productividad agregada en ambos países.\footnote{El número de entrantes en la economía pequeña puede reducirse a cero si los costes comerciales caen por debajo de un cierto umbral, de modo que el país pequeño deja de producir el bien diferenciado. Se asume que los costes comerciales permanecen por encima de ese nivel.
} 

Los efectos cualitativos son los mismos que en la transición de autarquía a comercio: aumenta la variedad, disminuyen los markups y los precios, y el bienestar crece como resultado de mayor productividad, menor poder de mercado y mayor variedad.\footnote{En el corto plazo (extensión), la liberalización simétrica genera resultados cualitativamente idénticos. Aunque estos efectos no dependen del tamaño relativo de los países (siempre que ambos produzcan el bien diferenciado), las diferencias de tamaño sí afectan al patrón de entrada en el largo plazo. Si $L^h > L^b$, la diferencia positiva en el número de entrantes $N_E^h - N_E^b$ se amplía con la liberalización, ya que el mercado grande se vuelve relativamente más atractivo. Esto también incrementa la diferencia en el número de productores domésticos entre países.
}

$$\text{\textbf{4.2.} Liberalización unilateral (asimétrica)}$$
A continuación, se considera la liberalización unilateral por parte del país $h$ (un aumento en su grado de apertura, manteniendo constante el del otro país). 

A partir de la condición de umbral de costes $c_D^h$, esto implica un aumento del umbral $c_D^h$, es decir, una reducción de la competencia en el país que liberaliza, mientras que el umbral $c_D^b$ en el país socio disminuye, indicando mayor competencia allí. 

En consecuencia, el país que liberaliza sufre una pérdida de bienestar, mientras que su socio obtiene una ganancia.\footnote{Una vez más, la respuesta de los umbrales determina la respuesta de todas las demás variables agregadas.
} Este resultado se debe al cambio en la localización de las empresas: el número de entrantes en el país liberalizador disminuye ($N_E^h$), mientras que aumenta en el otro país ($N_E^b$).

\begin{specialblock}{Diferencias entre el corto y largo plazo -- Relocalización empresarial}
    \textcolor{red}{\textbf{OJO}!} Hay grandes diferencias entre el largo plazo y el corto plazo (extensión).
    
    En el corto plazo, la condición de equilibrio muestra que el umbral $c_D^h$ disminuye tras la liberalización (en $h$), mientras que $c_D^b$ permanece constante ($b$ no liberaliza). Esto pone de manifiesto el efecto procompetitivo inmediato: \textbf{el aumento de la competencia de importaciones obliga a salir a las empresas menos productivas del país liberalizador} ($h$). Sin embargo, la variedad aumenta, ya que el incremento de exportadores hacia ese país domina la reducción de productores domésticos. Por tanto, el bienestar aumenta en el corto plazo gracias a mayor productividad, menores markups y mayor variedad, mientras que en el país socio no cambia. Esto evidencia que la pérdida de bienestar en el largo plazo se debe exclusivamente al desplazamiento de la entrada hacia el país no liberalizador.

    Estos efectos de relocalización empresarial han sido ampliamente estudiados en la literatura (por ejemplo, Ingo Horstmann y James Markusen, 1986; Anthony Venables, 1985 y 1987; y la síntesis en Elhanan Helpman y Paul Krugman, 1989). La contribución del modelo es mostrar cómo la liberalización también afecta simultáneamente a la selección de empresas, la productividad agregada, la variedad y los markups dentro de un único marco.\footnote{La mayor parte de la literatura previa asume empresas representativas, por lo que no captura el vínculo entre política comercial y productividad agregada a través de la selección.
    }
\end{specialblock}

$$\text{\textbf{4.3.} Liberalización preferencial}$$
Finalmente, se analiza la liberalización preferencial en un entorno de tres países. Hasta ahora, el análisis con dos países ya proporciona resultados ricos, pero ignora la posición de cada país en la red global de comercio. 

Cuando las barreras son simétricas, el modelo multi-país es una extensión directa. Sin embargo, aparecen nuevos efectos cuando las barreras bilaterales difieren.

Se introduce un tercer país $T$, con el mismo tamaño que los otros ($L^T= L^A = L^B=L$). Los países difieren únicamente en las barreras comerciales bilaterales, que son simétricas por pares. El grado de libertad de comercio ($\rho$) entre dos países $h$ y $b$ se mide mediante $\rho^{h,b}=\left(\tau^{h,b}\right)^{-k}=\rho^{b,h}=\left(\tau^{b,h}\right)^{-k}$. De forma similar, $\rho^{h,t}$ y $\rho^{b,t}$ miden la libertad de comercio entre esos países y el tercero. Como en el caso de la liberalización unilateral, laliberalización preferencial (no proporcional en las tres barreras comerciales) induce un importante cambio en el patrón de entrada de los países en el largo plazo. 

Con tres países de igual tamaño, la condición de libre entrada (22) en el país $h$ se convierte en:

$$(c_D^h)^{k+2} + \rho^{h,b}(c_D^b)^{k+2} + \rho^{b,t}(c_D^t)^{k+2} = \frac{\gamma\Gamma}{L}, \quad h = \{A, B, T\}, \; t \neq l \neq h$$.

Esto proporciona un sistema de tres ecuaciones lineales en los tres umbrales domésticos. Cuando las barreras comerciales son simétricas por pares, los umbrales de largo plazo vienen dados por:
$$c_D^h = \left[\frac{\gamma\Gamma}{L}\cdot\frac{(1 - \rho^{b,t})[1+\rho^{b,t}-(\rho^{h,b}+\rho^{h,t})]}{1 + 2\rho^{h,b}\rho^{h,t}\rho^{b,t} - (\rho^{h,b})^2 - (\rho^{h,t})^2 - (\rho^{b,t})^2}\right]^{\frac{1}{k+2}}$$

El número correspondiente de empresas que venden en el país $h$ sigue estando dado por la ecuación que define $N^h$, mientras que el número de entrantes satisface:

$$N_E^h + \rho^{h,b} N_E^b + \rho^{h,t} N_E^t = N^h \left(\frac{c_M}{c_D^h}\right)^k$$.

Por tanto, las diferencias internacionales en los umbrales se derivan de la medida relativa de libertad comercial:

$$\left(1 - \rho^{b,t}\right)\left[1 + \rho^{b,t} - (\rho^{h,b} + \rho^{h,t})\right]$$,

Lo que implica que el umbral es más bajo en el país $h$ que presenta la menor suma de barreras bilaterales (es decir, mayor $\rho^{h,b} + \rho^{h,t}$). En efecto, este país actúa como la mejor base de exportación o \textit{hub}. Además, dado que $\rho^{b,t}$ entra en la expresión del umbral del país $h$, cualquier cambio en las barreras bilaterales afecta a los tres países: efecto tercer país.\footnote{Esto tiene implicaciones importantes para los acuerdos comerciales preferenciales. Se puede demostrar cómo este \textbf{efecto de tercer país} puede incorporarse en una ecuación de gravedad.
}

Para verlo con claridad, considérese el caso en el que inicialmente las barreras comerciales son simétricas $(\rho^{h,t} = \rho)$. Los umbrales iniciales son entonces idénticos y vienen dados por:

$$c_D = \left[\frac{\gamma\Gamma}{L} \cdot \frac{1}{1 + 2\rho}\right]^{\frac{1}{k+2}}$$.

Se introduce ahora un acuerdo comercial preferencial entre $A$ y $B$, de modo que $\rho^{A,B} = \rho' > \rho = \rho^{A,T} = \rho^{B,T}$. El nuevo régimen comercial afecta a los umbrales en todos los países. 

A partir de la ecuación de costes anterior:

$$\begin{aligned}
    &\text{Liberalizan}:\qquad c_D^A = c_D^B = \left[\frac{\gamma\Gamma}{L}\cdot\frac{1 - \rho}{1 - 2\rho^2 + \rho'}\right]^{\frac{1}{k+2}},\quad \Longrightarrow\quad N_E^A = N_E^B =\frac{2(c_M)^k (k + 1)\gamma}{\eta (1 - 2\rho^2 + \rho')}\left[ \frac{\alpha - c_D^A}{(c_D^B)^{k+1}} - \rho \frac{\alpha - c_D^T}{(c_D^T)^{k+1}}\right],\\[1em]
    &\text{No liberaliza}:\qquad c_D^T = \left[\frac{\gamma\Gamma}{L} \cdot \frac{(1 - \rho) + (\rho' - \rho)}{1 - 2\rho^2 + \rho'} \right]^{\frac{1}{k+2}},\quad \Longrightarrow\quad N_E^T =\frac{2(c_M)^k (k + 1)\gamma}{\eta (1 - 2\rho^2 + \rho')}\left[\frac{(1 + \rho')(\alpha - c_D^T)}{(c_D^T)^{k+1}} - 2\rho \frac{\alpha - c_D^H}{(c_D^H)^{k+1}} \right].
\end{aligned}$$

Comparando las ecuaciones se verifica fácilmente que la liberalización preferencial reduce los umbrales en los países que liberalizan y los aumenta en el tercer país. 

En consecuencia, los costes medios, precios y markups disminuyen en los países que liberalizan, mientras que aumentan en el país excluido. Los \textbf{países que participan en el acuerdo se convierten en mejores bases de exportación}: obtienen mejor acceso mutuo sin perder acceso al mercado del tercer país. Por tanto, la liberalización preferencial genera ganancias de bienestar en los países participantes y pérdidas de bienestar en el país excluido.



$$\text{\textbf{Extensión: Largo y corto plazo (distinción)}}$$
En las secciones siguientes se introduce una versión de economía abierta del modelo y se analizan las consecuencias de distintos escenarios de liberalización comercial. Los escenarios de liberalización asimétrica generan una relocalización bien conocida de empresas (en este modelo, de entrantes) entre países. Posteriormente, interesa distinguir estos \textbf{efectos de largo plazo} de los \textbf{efectos directos de corto plazo} de la liberalización sobre la competencia y la selección en los distintos mercados. Con este objetivo, se introduce una versión de corto plazo del modelo. Por ahora, se describen sus principales características y su equilibrio en una economía cerrada.

Hasta este punto, se ha considerado un escenario de largo plazo en el que las decisiones de entrada y salida estaban determinadas endógenamente. En contraste, el corto plazo se caracteriza por un número y una distribución de empresas incumbentes fijos. En este horizonte temporal, dichas empresas deciden si operar y producir o, por el contrario, cerrar temporalmente. En este último caso, pueden reanudar la producción posteriormente sin incurrir de nuevo en el coste de entrada. No se permite la entrada en el corto plazo: sea $\bar{I}$ el número fijo de empresas incumbentes y $\bar{C}(\varphi)$ su distribución de costes, con soporte en $[0, \bar{c}_M]$. Se mantiene la parametrización de PARETO para la productividad $1/c$, lo que implica que $\bar{C}(\varphi) = (c/\bar{c}_M)^k$. 

Al igual que en el largo plazo, únicamente producen las empresas que obtienen beneficios no negativos. Esto conduce a la misma determinación del umbral de costes $c_D$. Las empresas con $c > c_D$ cesan su actividad, mientras que las restantes $I = \bar{I}\,\bar{C}(\varphi_D) = \bar{I}(c_D/\bar{c}_M)^k$ producen. El umbral $c_D$ se determina mediante la condición de beneficio nulo en el corte:\footnote{Si la empresa menos productiva, con coste $\bar{c}_M$, obtiene beneficios no negativos, entonces $c_D = \bar{c}_M$ y todas las empresas producen en el corto plazo. 
}

$$I = \frac{2\gamma}{\eta}\cdot \frac{\alpha - c_D}{c_D - \bar{c}} = \frac{2(k+1)\gamma}{\eta}\cdot \frac{\alpha - c_D}{c_D}, \quad \text{siempre que } c_D < \bar{c}_M,$$

ya que el coste medio de las empresas que producen sigue siendo $\bar{c} = \frac{k}{k+1}c_D$, como en el equilibrio de largo plazo. Utilizando la nueva condición para el número de empresas, $I = \bar{I}(c_D/\bar{c}_M)^k$, la condición de beneficio nulo en el umbral implica:

$$\frac{(c_D)^{k+1}}{\alpha - c_D} = \frac{2(k+1)\gamma}{\eta}\cdot \frac{(\bar{c}_M)^k}{\bar{I}}, \quad \text{siempre que } c_D < \bar{c}_M$$,

lo que determina de manera única el umbral de costes de corto plazo $c_D$ y el número de empresas que producen $I$.

En este equilibrio de corto plazo, los cambios en el tamaño del mercado no alteran la distribución de empresas productoras (es decir, $c_D$ permanece constante), ni la distribución de precios y markups. Todas las empresas ajustan sus niveles de producción proporcionalmente al cambio en el tamaño del mercado.\footnote{Cuando el tamaño del mercado cambia, la curva de demanda inversa residual rota en torno al mismo precio máximo $p_{\max}$. Con curvas de demanda lineales, la curva de ingreso marginal rota de tal forma que el precio que maximiza beneficios para un coste marginal dado permanece inalterado.
} Únicamente la entrada en el largo plazo genera reasignaciones entre empresas (y el correspondiente cambio en el umbral $c_D$).



\clearpage
\anexo{Modelo ARA: Equilibrio General}\label{anx:modelo_ara}
Todo lo que se va a exponer aquí es una traducción literal del artículo: \href{https://tomohiroara.com/Hat.pdf}{\textit{Optimal Tariffs in the Melitz Model: A Sufficient Statistics Approach for Trade Policy}. ARA (2025)}

\begin{specialblock}{\textbf{Aclaración}.- Diferencias entre la perspectiva de MELITZ y OTTAVIANO (2008) y ARA (2025)}
    Se puede aceptar la siguiente reducción general de las diferencias como una aproximación pedagógica, siempre que aclaremos bien qué hay detrás de cada bloque.
    
    En primer lugar, con la inclusión de salarios y el paso a un marco de equilibrio general no pretende describir la microestructura del mercado, sino los resultados agregados suficientes para la evaluación del bienestar. En el enfoque de ARA los salarios pasan a ser endógenos, lo que introduce una restricción de equilibrio general: el mercado de trabajo debe vaciarse.\footnote{En Melitz and Ottaviano (2005), los salarios suelen tomarse como dados (o normalizados), lo que implica que el ajuste ante cambios en costes comerciales ocurre únicamente vía selección de firmas, precios y márgenes.
    } Esto implica que cualquier cambio en política comercial afecta simultáneamente a la estructura productiva y al nivel salarial, y por tanto al ingreso real. La consecuencia clave es que el bienestar ya no depende solo de los términos de intercambio (ToT) o de la selección de firmas, sino también del canal de renta: el arancel afecta a los precios relativos, pero también a los salarios, y por tanto al poder adquisitivo. Por tanto, ARA modifica la forma en que se interpreta la política comercial óptima, porque introduce trade-offs adicionales que no existen en el modelo original.
    
    En segundo lugar, tenemos la descomposición de la elasticidad comercial. ARA reinterpreta la fórmula del arancel óptimo en términos de elasticidades \textit{suficientes}. Pero lo verdaderamente novedoso no es solo descomponer la elasticidad, sino mostrar que el objeto relevante para la PC son sus componentes estructurales (margen extensivo, intensivo, entrada/salida, etc.), que a su vez dependen del entorno competitivo y del equilibrio general.\footnote{En MELITZ y OTTAVIANO (2005), muchas de las propiedades cuantitativas (por ejemplo, elasticidades constantes) dependen críticamente de asumir una distribución concreta (como PARETO). ARA evita esto al trabajar con elasticidades observables o estimables directamente, lo que hace que su marco sea más general y aplicable empíricamente.
    } 
    
    Aunque no son dos bloques independientes (uno de equilibrio general y otro de elasticidades), se pueden entender como dos caras de la misma contribución. La endogeneización de salarios y el equilibrio general es lo que hace que la elasticidad comercial relevante cambie de naturaleza; y la descomposición de esa elasticidad es lo que permite capturar de forma tractable los efectos de ese nuevo entorno.

    .En Melitz and Ottaviano (2005), el modelo está construido precisamente para seguir de forma explícita toda la cadena microeconómica: distribución de productividad, número de empresas, selección, markups, tamaño de las firmas, etc. Esto implica que, bajo el marco de MELITZ y OTTAVIANO se puede realizar un análisis sectorial muy fino: qué pasa con las firmas marginales, cómo cambian los markups, quién entra, quién sale, etc.  En cambio, ARA se colapsa toda esa riqueza micro en un conjunto muy reducido de objetos (las \textit{sufficient statistics}), en particular la elasticidad comercial y la cuota doméstica. El punto de partida del paper es que las ganancias del comercio pueden expresarse solo con esas dos estadísticas  
    
    Por ello el marco de ARA es menos útil para analizar la estructura del mercado porque ha eliminado (endógenamente) el nivel de detalle donde vive esa estructura. No sigue la distribución completa de firmas ni los markups individuales; todo queda resumido en elasticidades agregadas. Es decir, sabemos cuánto responde el comercio, pero no quién ajusta dentro del mercado. Imponer estas restricciones agregadas (o falta de detalle) persigue un objetivo claro:
    \begin{la}
        \item Bajo el marco de MELITZ y OTTAVIANO, si introducimos simultáneamente heterogeneidad de firmas, salarios endógenos, ajustes múltiples y Política Comercial Óptima, el modelo completo tipo MELITZ–OTTAVIANO deja de ser manejable analíticamente. No se pueden obtener fórmulas limpias de política comercial porque todo depende de la distribución completa de productividad y de su interacción con el equilibrio.
        \item ARA resuelve esto imponiendo dos tipos de simplificaciones estructurales: (i) sustituye la distribución completa de productividad por elasticidades agregadas, que capturan los márgenes intensivo y extensivo ($\gamma$ margen intensivo, $\sigma-1$ margen extensivo), sin necesidad de especificar toda la distribución (la propia elasticidad comercial pasa a ser una combinación de ambos márgenes); y, (ii) reorganiza el modelo para que el equilibrio pueda describirse con un sistema reducido de ecuaciones en cambios, en lugar de tener que resolver toda la distribución de equilibrio. De esta forma podemos derivar una fórmula de Política Comercial óptima en términos de objetos observables o estimables (elasticidades), sin depender de supuestos fuertes sobre la distribución de productividad. 
    \end{la}
\end{specialblock}

$$\text{Introducción}$$
Las ganancias del comercio pueden calcularse mediante únicamente dos estadísticas suficientes en una amplia clase de modelos de comercio (Costas Arkolakis et al., 2012): (i) la respuesta de los flujos comerciales a cambios en los costes comerciales; y, (ii) la cuota de gasto doméstico. 

Impulsada por estos resultados teóricos, la investigación empírica reciente ha estimado estas dos estadísticas suficientes y ha puesto de manifiesto una heterogeneidad significativa en dichas medidas, que varía sistemáticamente con las características de los países:
\begin{lnum}
    \item El primer elemento, la \textbf{elasticidad comercial}, tiende a diferir según si los pares de países son próximos o lejanos, y grandes o pequeños;\footnote{Por ejemplo, BAS et al. (2017) encuentran que la elasticidad comercial es menor para pares de países próximos, donde el volumen de comercio ya es elevado.
    } 
    \item El segundo elemento, la \textbf{cuota de comercio doméstico}, también se ve considerablemente afectado por el entorno comercial, en el sentido de que dicha cuota es mayor cuanto más grandes y menos abiertos son los países.\footnote{For the trend in the domestic trade share, see Eaton and Kortum (2002, 2011) with aggregate data, and Bernard et al. (2007) and Mayer and Ottaviano (2008) with firm-level data.
    } 
\end{lnum}

Estas evidencias empíricas implican que la liberalización comercial y el tamaño de mercado pueden tener efectos diferentes sobre las dos estadísticas suficientes que determinan las ganancias del comercio. Por tanto, es necesario tener en cuenta las diferencias entre pares de países en estas estadísticas para evaluar correctamente el impacto del comercio, como la reducción de los costes comerciales y la expansión del tamaño de mercado mediante acuerdos comerciales.


¿Cómo afecta el hecho de que estas dos estadísticas suficientes varíen entre pares de países a la política comercial óptima? Para abordar esta cuestión clave, desarrollamos una versión con países asimétricos del modelo de MELITZ (2003) con preferencias CES y competencia monopolística. Uno de los inconvenientes de este marco es que, cuando la productividad sigue una distribución de Pareto —una de las distribuciones más utilizadas en la literatura—, la elasticidad comercial es única para cualquier par de países, independientemente de sus características. Además, los markups de las empresas son constantes, lo que implica que, bajo heterogeneidad de empresas, el tamaño de mercado no tiene efecto sobre la cuota de comercio doméstico a través de la selección. 

Para superar estas limitaciones y proporcionar implicaciones de política más realistas, introducimos tres modificaciones principales respecto a la literatura existente: (i) en primer lugar, trabajamos con una distribución general de productividad que genera una elasticidad comercial específica para cada par de países; (ii) en segundo lugar, consideramos salarios endógenos que restablecen el papel del tamaño de mercado en la determinación de la cuota de comercio doméstico a través de la selección; y, por último, (iii) analizamos no solo los costes comerciales tipo iceberg, sino también los aranceles de importación que el gobierno elige para maximizar el bienestar. 

Estas diferencias, consideradas conjuntamente, nos permiten comprender los distintos efectos de las presiones competitivas sobre las estadísticas suficientes y analizar sus consecuencias para la política comercial óptima en un único marco unificado.

Nuestro punto de partida es observar que no solo los costes comerciales, sino también el tamaño de mercado, tienen un efecto crucial sobre los salarios.\footnote{Los costes comerciales han ido disminuyendo de forma sostenida a lo largo del tiempo tanto por mejoras tecnológicas como por negociaciones comerciales. Por ejemplo, HUMMELS (2007) encuentra que el precio del transporte aéreo internacional por tonelada se ha reducido más de diez veces en todo el mundo entre 1955 y 2004, debido en gran medida a la adopción de motores a reacción; de manera similar, el esfuerzo continuo de la Organización Mundial del Comercio (World Trade Organization) ha reducido los aranceles medios mundiales del 8,6\% al 3,2\% entre 1960 y 1995, incrementando significativamente los salarios de los países que comercian.
} La Figura 1 muestra la evolución de la población y del PIB per cápita como medidas, respectivamente, del tamaño de mercado y de los salarios. El panel A presenta el caso de Estados Unidos, mostrando una clara relación monótona entre población y PIB per cápita. En contraste, el panel B muestra el caso de Japón, donde la población está disminuyendo gradualmente debido principalmente a la baja tasa de natalidad.\footnote{Según la Oficina del Gabinete de Japón, se espera que la población disminuya de 124 millones en 2020 a 97 millones en 2050 y a 86 millones en 2060. Se suele argumentar que la contracción progresiva del tamaño de su mercado interno, junto con una fuerte dependencia de la demanda exterior, podría provocar una fuerte caída del PIB per cápita en Japón (Nikkei Asia, 2019).
} Esto muestra que los cambios tanto en los costes comerciales como en el tamaño de mercado han sido significativos en las últimas décadas, afectando de manera crítica a los salarios y, por tanto, al bienestar nacional.

\imagenfit{ARA_Model.png}{Población y PIB per cápita durante 1960-2020 (\textbf{Nota}: El eje izquierdo mide población en unidades de millones, PIB per cápita en millares a la derecha)}{WBG y reconstrucción del autor.}

Mostramos que, si los salarios son endógenos, un país grande implica una selección doméstica débil, lo que permite que empresas ineficientes sobrevivan en el mercado interno. Para entender esta razón, considérese una reducción unilateral de los costes comerciales, que tiene dos efectos sobre los beneficios esperados de las empresas. En primer lugar, dicha liberalización reduce directamente los beneficios esperados al disminuir los markups sobre los bienes extranjeros. En segundo lugar, incrementa indirectamente los beneficios esperados al reducir los salarios, es decir, los costes de producción de las empresas, lo que ocurre para restablecer el equilibrio comercial (Demidova y Rodríguez-Clare, 2013). En cambio, una expansión unilateral del tamaño de mercado no tiene un efecto directo sobre los beneficios esperados a través de los markups bajo preferencias CES y competencia monopolística, mientras que sí reduce indirectamente los beneficios esperados al elevar los salarios debido al efecto mercado interno (Krugman, 1980). Dado que el aumento de los salarios eleva los costes de producción y reduce la rentabilidad, la entrada de empresas se relocaliza desde el país en expansión, con salarios más altos, hacia el país no expansivo, con salarios más bajos. Esta relocalización permite que empresas ineficientes sobrevivan en el mercado doméstico.

La débil selección doméstica asociada a salarios más altos puede explicar el cambio en los patrones de comercio que experimentó Japón a finales de los años noventa. Por ejemplo, utilizando datos a nivel de empresa del sector manufacturero en Japón, Fukao et al. (2008) analizan cómo las diferencias de productividad afectaron a la rotación de empresas entre 1990 y 2003. Encuentran que la tasa de rotación era significativamente mayor para las empresas menos productivas, pero que casi la mitad del 10\% superior de las empresas más productivas también abandonó el mercado. Este hecho aparentemente paradójico puede explicarse por el aumento de los salarios en Japón, que indujo a estas empresas altamente productivas a buscar mano de obra más barata en mercados extranjeros como China, al tiempo que permitió que empresas menos productivas sobrevivieran en el mercado doméstico japonés.

Nuestro efecto de selección asociado al tamaño de mercado contradice el resultado de MELITZ y OTTAVIANO (2008). La razón clave radica en la presencia de un bien exterior que hace que los salarios estén fijados exógenamente en su modelo. En ese caso, las diferencias en tamaño de mercado generan patrones de comercio tales que un país grande (respectivamente pequeño) se especializa en el sector de bienes diferenciados (respectivamente en el bien exterior). Como resultado, cuanto mayor es el tamaño de mercado, más intensa es la competencia en el sector de bienes diferenciados, lo que obliga a las empresas menos eficientes a salir del mercado. Sin embargo, si este bien exterior está ausente, los salarios se determinan endógenamente por el equilibrio comercial y las diferencias en tamaño de mercado no generan estos patrones de comercio a través de cambios en salarios. Por tanto, no resulta sorprendente que el efecto del tamaño de mercado sobre la selección difiera del de Melitz y Ottaviano (2008).\footnote{También difiere en la especificación de las preferencias. Sin embargo, la ausencia del bien exterior es suficiente incluso con mark-ups variables (Demidova, 2017).
} Aunque un país grande sufre una selección doméstica débil, puede, no obstante, disfrutar de ganancias de bienestar derivadas de su tamaño de mercado, ya que el efecto negativo sobre la selección puede verse compensado por el efecto positivo sobre la variedad de productos.

Dado que los salarios endógenos generan efectos distintos de los costes comerciales y del tamaño de mercado sobre la selección, ¿qué implicaciones tiene esto para la política económica? Al analizar la política comercial óptima, mostramos que esta diferencia es crucial para la caracterización de los \textbf{aranceles óptimos, es decir, aquellos que maximizan el bienestar que cada país impondría sin temer represalias}. En el presente modelo, los aranceles óptimos están inversamente relacionados con la elasticidad de la oferta de exportaciones del socio comercial, que a su vez se compone tanto de la cuota de comercio doméstico como de la elasticidad comercial, como en modelos previos. Sin embargo, encontramos que la reducción de los costes comerciales y la expansión del tamaño de mercado no conducen necesariamente a aranceles óptimos más elevados en nuestro modelo. Desde el punto de vista de la política económica, el efecto del tamaño de mercado sobre los aranceles óptimos es particularmente relevante: un país grande no siempre se beneficia de aranceles elevados: un país grande puede obtener una ganancia en términos de intercambio al fijar aranceles, como en la teoría convencional del arancel óptimo, pero también sufre una pérdida de bienestar derivada de una débil selección doméstica, ya que los aranceles amplifican esta pérdida al proteger a las empresas ineficientes de la competencia extranjera. 

A partir de este trade-off asociado a la selección, el hecho de que el beneficio derivado de los aranceles supere o no su coste depende de si la elasticidad comercial es constante o variable. Si la elasticidad comercial es variable y difiere entre mercados los aranceles óptimos pueden incluso disminuir con el tamaño de mercado a través de la respuesta endógena de dicha elasticidad.\footnote{See, for example, Helpman et al. (2008), Novy (2013), Spearot (2013) and Bas et al. (2017). Our theoretical approach is close to Helpman et al. (2008) and Bas et al. (2017) who rest on CES preferences and monopolistic competition to provide evidence.
}

Para comprender mejor este resultado de política económica se descompone la elasticidad comercial en:
\begin{la}
    \item \textbf{Margen intensivo}. Elasticidad del margen intensivo, es la elasticidad de las ventas de cada empresa incumbente. Es decir, como varía la intensidad en ventas de las empresas que sobreviven o perduran en el mercado doméstico; y,
    \item \textbf{Margen extensivo}. Elasticidad del margen extensivo, es la elasticidad de las ventas de nuevas empresas entrantes. Es decir, como disminuye o se amplía las variedades en el mercado doméstico por la incorporación de más empresas exportadoras (más variedades importadas).
\end{la} 

Dado que la elasticidad del margen intensivo es constante bajo preferencias CES y competencia monopolística, la naturaleza variable de la elasticidad comercial debe provenir del margen extensivo, que a su vez depende de la microestructura del modelo.\footnote{En el modelo con empresas homogéneas donde todas exportan, no existe margen de ajuste vía entrada de nuevas empresas (es decir, la elasticidad extensiva es cero) y, por tanto, la elasticidad comercial coincide con la elasticidad intensiva. En el modelo con empresas heterogéneas donde la productividad sigue una distribución de PARETO, la elasticidad extensiva es constante y, por tanto, también lo es la elasticidad comercial (Chaney, 2008). En estos casos especiales, el tamaño de mercado no tiene efecto sobre la elasticidad comercial y los aranceles óptimos siempre aumentan con el tamaño de mercado únicamente a través de la cuota de comercio doméstico (Gros, 1987; Felbermayr et al., 2013).
} Aquí se produce una moficiación sustancial: si la distribución de productividad no está restringida (no PARETO), el resultado teórico de una elasticidad comercial constante deja de cumplirse. Es importante destacar que la evidencia empírica ha encontrado variaciones sustanciales en la elasticidad comercial entre pares de países, algo que no puede ser explicado por la mayoría de modelos previos. Por tanto, debido a este canal adicional, los aranceles óptimos no necesariamente aumentan con el tamaño de mercado.

Aunque este trabajo analiza principalmente el aspecto cualitativo de los aranceles óptimos con elasticidad comercial variable, nuestros resultados analíticos también permiten medir cuantitativamente la discrepancia en los aranceles óptimos que surge cuando se supone erróneamente que la elasticidad comercial es constante cuando en realidad es variable. Nuestro modelo calibrado con datos de Estados Unidos muestra que los aranceles óptimos con elasticidad variable son sustancialmente más bajos que aquellos obtenidos bajo elasticidad constante. En nuestro ejercicio numérico, \textbf{los niveles de aranceles óptimos con elasticidad variable son aproximadamente dos tercios (o incluso menos de la mitad) de los correspondientes a una elasticidad constante en modelos con empresas heterogéneas (u homogéneas)}. Esta diferencia se debe al supuesto implícito de elasticidad constante en gran parte de la literatura. Utilizando soluciones analíticas de estática comparada, también encontramos que el efecto del tamaño de mercado sobre los aranceles óptimos es cuantitativamente mucho menor que el efecto de los costes comerciales variables. Este resultado cuantitativo es coherente con nuestra predicción teórica de que un país grande no siempre se beneficia de fijar aranceles elevados. Asimismo, discutimos cómo la calibración de los aranceles óptimos con datos de Estados Unidos tiene relevancia práctica para las políticas comerciales reales, como los recientes aumentos arancelarios bajo la administración Trump, proporcionando así información útil para los responsables de política económica.


Un conjunto de trabajos ha analizado las implicaciones en términos de bienestar y política económica tanto en modelos con empresas homogéneas como heterogéneas. En lo que respecta al bienestar, Costas Arkolakis et al. (2012) derivan una fórmula sencilla que permite capturar las ganancias del comercio únicamente a partir de la cuota de comercio doméstico y de la elasticidad comercial. Dado que este resultado se aplica a una sorprendentemente amplia clase de modelos de comercio, trabajos posteriores han estudiado extensiones y la robustez de dichos resultados. Por ejemplo, Arkolakis et al. (2019) analizan funciones de demanda generales que generan markups variables, Felbermayr et al. (2015) introducen aranceles que generan ingresos públicos, y Head et al. (2014) y Melitz y Redding (2015) consideran una distribución no Pareto que hace que la elasticidad comercial sea variable. Mostramos que la fórmula de bienestar de Arkolakis et al. (2012) también es aplicable al resultado convencional sobre aranceles óptimos. En particular, condicionando a sus estadísticas suficientes para el bienestar, el nivel óptimo de los aranceles de importación es el mismo en distintos modelos de comercio con elasticidad comercial constante, pero, en general, dicho nivel depende de la microestructura que hace que la elasticidad comercial sea variable.[Footnote: 4Our sufficient statistics approach to optimal trade policy is related to that by Lashkaripour (2021) in that such policy can be calculated by the welfare formula by Arkolakis et al. (2012). One of the critical differences is that his baseline analysis builds on the Ricardian setup of Eaton and Kortum (2002), which means that the trade elasticity is constant at the supply side parameter.] También encontramos que la heterogeneidad de empresas derivada de una distribución distinta de Pareto puede afectar a la medición del bienestar, como en Head et al. (2014) y Melitz y Redding (2015); sin embargo, el alcance de este trabajo difiere del suyo, ya que mostramos analíticamente una nueva fórmula de arancel óptimo con elasticidad comercial variable e investigamos cuantitativamente el efecto de las dos principales fuentes de presión competitiva sobre la política comercial óptima. Además, al igual que en el trabajo original de Marc Melitz (2003), estos trabajos consideran principalmente comercio entre países simétricos y, por tanto, no pueden distinguir entre los efectos bilaterales y unilaterales de perturbaciones exógenas sobre la política comercial óptima.

En cuanto a las implicaciones de política económica, existe una amplia literatura sobre aranceles óptimos. Gros (1987) deriva aranceles óptimos en el modelo con empresas homogéneas, los cuales están inversamente relacionados con la elasticidad comercial y con la cuota de comercio doméstico del socio comercial. Utilizando el marco de Ossa (2011), que incorpora efectos de relocalización de la producción inducidos por aranceles, Ossa (2014) proporciona un análisis exhaustivo de los aranceles óptimos en un modelo multisectorial de equilibrio general que integra los motivos tradicionales (términos de intercambio), de nueva teoría del comercio (desplazamiento de beneficios) y de economía política en el modelo de empresas homogéneas. Estos análisis de aranceles óptimos se extienden al modelo con empresas heterogéneas por Demidova y Rodríguez-Clare (2009) para una economía pequeña y por Felbermayr et al. (2013) para una economía grande. En este último caso, Felbermayr et al. (2013) muestran que los aranceles óptimos son menores en el modelo con empresas heterogéneas que en el modelo con empresas homogéneas, manteniendo constante la cuota de comercio doméstico. Aunque estos trabajos han contribuido a mejorar nuestra comprensión de la política comercial óptima, una de sus limitaciones es que la elasticidad comercial es constante tanto en modelos con empresas homogéneas como heterogéneas. Sin embargo, la existencia de una elasticidad comercial constante es altamente sensible a las restricciones paramétricas, y los cambios en bienestar pueden estimarse incorrectamente cuando la elasticidad “real” es variable (Melitz y Redding, 2015). En el contexto de la política comercial, el nivel óptimo de los aranceles de importación también puede estimarse erróneamente cuando se imponen esas mismas restricciones paramétricas. Destacamos esta limitación clave no solo caracterizando analíticamente los aranceles óptimos con una elasticidad comercial variable, sino también midiendo cuantitativamente estas magnitudes a partir de nuestro modelo y de modelos existentes calibrados con datos de Estados Unidos.

Recientemente, Arnaud Costinot et al. (2020) ofrecen una generalización estricta de Gros (1987) en el modelo con empresas homogéneas, y de Demidova y Rodríguez-Clare (2009) y Felbermayr et al. (2013) en el modelo con empresas heterogéneas con distribución Pareto. Encuentran que, cuando los aranceles son uniformes para todas las empresas, los aranceles óptimos pueden ser menores en el modelo con empresas heterogéneas que en el modelo con empresas homogéneas debido a la no convexidad de los bienes agregados entre mercados domésticos y extranjeros. En contraste, nosotros mostramos que el mismo resultado puede surgir debido a la variabilidad de la elasticidad comercial entre mercados domésticos y extranjeros. Aunque este nuevo elemento en la determinación de los aranceles óptimos está estrechamente relacionado con el nuestro, en el sentido de que ambos surgen en presencia de selección, nuestro enfoque es relativamente fácil de medir a partir de datos a nivel de empresa, lo que lo hace directamente aplicable a la cuantificación de los aranceles óptimos. En este sentido, investigamos un canal distinto pero complementario en el análisis de la política comercial óptima.


$$\text{\textbf{2.} Modelo}$$
En esta sección, describimos en primer lugar el comportamiento de los consumidores, las empresas y los gobiernos que desempeñan un papel crucial en el modelo. A continuación, definimos las condiciones de equilibrio que determinan las principales variables endógenas del modelo.

$$\text{\textbf{2.1} Fundamentos}$$
Existen dos países, indexados por $i$, $j$, que utilizan únicamente trabajo para producir bienes diferenciados en un único sector. El país $i$ está poblado por una masa $L_i$ de consumidores idénticos cuyas preferencias vienen dadas por:

$$U_i =\left(\sum_{n=i,j}\int_{\omega \in \Omega_n}q_{ni}(\omega)^{\rho} \, d\omega\right)^{1/\rho},\qquad 0 < \rho < 1,$$

Donde la elasticidad de sustitución entre variedades es $\sigma = \frac{1}{1-\rho} > 1$. A lo largo del trabajo, denotamos el país exportador (importador) mediante el primer (segundo) subíndice; así, $q_{ji}(\omega)$ es la cantidad enviada desde el país $j$ al país $i$ de la variedad $\omega$. Sea $p_{ji}(\omega)$ el precio de la variedad correspondiente. Dado que los consumidores obtienen utilidad del consumo tanto de variedades domésticas como extranjeras, la restricción presupuestaria se expresa como $\sum_n \int_\omega p_{ni}(\omega) q_{ni}(\omega)\, \mathrm d\omega$.

La maximización de la utilidad del consumidor sujeta a esta restricción da lugar a la siguiente demanda:

$$q_{ji}(\omega) = R_i P_i^{-1} \, p_{ji}(\omega)^{-\sigma},\qquad \text{donde},\,\, P_i = \left( \sum_n \int_\omega p_{ni}(\omega)^{1-\sigma} d\omega \right)^{1/(1-\sigma)}$$

$P_i$ es el índice de precios asociado a un bien agregado definido como $Q_i \equiv U_i$, y $R_i = P_i Q_i$ es el gasto agregado.

El comportamiento de las empresas es similar al modelizado por MELITZ (2003). Tras pagar un coste fijo de entrada $f_i^e$, una masa $M_i^e$ de entrantes extrae aleatoriamente una productividad $\varphi$ de una distribución $G_i(\varphi)$. La distribución tiene soporte $(\varphi_{\min}, \varphi_{\max})$, donde el límite superior puede ser finito $(\varphi_{\max} < \infty)$ o infinito $(\varphi_{\max} = \infty)$. Tras observar su productividad $\varphi$, cada empresa decide si salir o producir. 

Si una empresa del país $i$ decide producir para el mercado doméstico, incurre en un coste fijo $f_{ii}$ y en costes marginales constantes inversamente proporcionales a la productividad:

$$l_{ii}(\varphi) = f_{ii} + \frac{q_{ii}(\varphi)}{\varphi}$$.

Por otro lado, si una empresa del país $j$ decide producir para el mercado del país $i$, incurre en costes comerciales variables $\theta_{ji}$ y en costes fijos de exportación $f_{ji}$, con la misma forma funcional:

$$l_{ji}(\varphi) = f_{ji} + \frac{\theta_{ji} q_{ji}(\varphi)}{\varphi}$$.

Suponemos que $\tau_{ji} > 1$ (con $\theta_{jj} = 1$) y que todos los costes de producción se miden en unidades de trabajo del país de origen. Por ejemplo, una empresa del país $j$ que vende en el país $i$ incurre en un coste $w_j l_{ji}(\varphi)$, donde $w_j$ es el salario en el país $j$.

El gobierno impone aranceles ad valorem $\tau_{ji} = 1 + t_{ji} > 1$ (con $\tau_{jj} = 1$ y $t_{jj} = 0$) sobre los bienes extranjeros. En este trabajo, los aranceles se modelizan como “demand shifters” que afectan directamente a la demanda de los consumidores por bienes extranjeros.\footnote{Aunque los aranceles modifican el precio de estos bienes, pueden aplicarse antes o después de que las empresas fijen sus markups, lo que se conoce como aranceles como “cost shifters” o “demand shifters”.

Los aranceles que son cost shiffters afectan directamente al coste marginal de las empresas extranjeras. FELBERMAYR et al. (2015) demuestra que la distinción entre ambos tipos de aranceles (aunque la evidencia no es concluyente) puede ser cuantitativamente importante para el nivel de bienestar.
} Por tanto, una empresa del país $j$ que vende en el país $i$ recibe únicamente el precio neto de arancel $p_{ji}(\varphi)/\tau_{ji}$ (es decir, $p_{ji}(\varphi)$ es el precio con arancel incluido en el país $i$), mientras que el gobierno del país $i$ recauda ingresos arancelarios iguales a $(\tau_{ji}-1)\,p_{ji}(\varphi)/\tau_{ji}$ de dicha empresa.

Con ello se completa la descripción de las preferencias de los consumidores, la tecnología de las empresas y la política del gobierno en el modelo. A partir de esta estructura, el beneficio de la empresa viene dado por:

$$\pi_{ji}(\varphi) =\left(\frac{p_{ji}(\varphi)}{\tau_{ji}} - \frac{\theta_{ji} w_j}{\varphi}\right) q_{ji}(\varphi) - w_j f_{ji}$$


$$\text{\textbf{2.2} Equilibrio en niveles}$$
Como es habitual en competencia monopolística con un gran número de empresas, la decisión de precios de una empresa individual tiene un efecto despreciable sobre el comportamiento del resto. Por tanto, cada empresa fija su precio maximizando beneficios, tomando como dados el índice de precios y el gasto agregado en la demanda del consumidor. 

Dado que una empresa del país $j$ que vende en el país $i$, con productividad $\varphi$, incurre en costes marginales $\tau_{ji} w_j / \varphi$ así como en aranceles $\tau_{ji}$, la maximización de beneficios implica que el precio óptimo consiste en aplicar un markup constante $\sigma = (\sigma - 1)^{-1} = 1/\rho$ sobre dichos costes:

$$p_{ji}(\varphi) = \frac{\tau_{ji}\,\theta_{ji} w_j}{\rho \,\varphi}$$.

Utilizando este precio óptimo para $\pi_{ji}(\varphi)$ y definiendo los ingresos de la empresa netos de aranceles como $r_{ji}(\varphi) \equiv p_{ji}(\varphi) q_{ji}(\varphi)/\tau_{ji}$, los beneficios variables de la empresa se expresan como $r_{ji}(\varphi)/\sigma$, los cuales son estrictamente crecientes en la productividad $\varphi$. Dado esto, existe un único umbral de productividad a partir del cual una empresa exportadora obtiene beneficios nulos, es decir, $r_{ji}(\varphi_{ji}^*)/\sigma = w_j f_{ji}$. 

Esto se denomina \textbf{condición de beneficio nulo en el umbral}:

$$B_i \, \tau_{ji}^{-\sigma} (\tau_{ji} w_j)^{1-\sigma} (\varphi_{ji}^*)^{\sigma-1} = \sigma w_j f_{ji}, \tag{1}$$

Donde $B_i \equiv R_i (\rho P_i)^{\sigma-1}$ es el índice de demanda del mercado. Nótese que (1) también determina el umbral de productividad doméstico $\varphi_{jj}^*$ cuando $i = j$. 

Nos centramos en el caso en el que existe selección en la exportación, es decir, $\varphi_{ji}^* > \varphi_{jj}^*$.

Utilizando (1), puede demostrarse fácilmente que la condición de selección se cumple cuando los costes comerciales son suficientemente elevados y el tamaño de los mercados no difiere en exceso; esto último implica que la demanda relativa $B_i/B_j$, proporcional al tamaño relativo del mercado medido por $R_i/R_j$, no es ni demasiado grande ni demasiado pequeña.

La libre entrada exige que los beneficios esperados obtenidos en todos los mercados de operación igualen los costes fijos de entrada:

$$J_i(\varphi^*) \equiv \int_{\varphi^*}^{\varphi_{\max}} \left[ \left(\frac{\varphi}{\varphi^*}\right)^{\sigma-1} - 1 \right] \mathrm dG_i(\varphi)$$

A partir de la definición de $\varphi_{ij}^*$ en (1), cualquier empresa del país $i$ que vende en el país $j$ obtiene beneficios esperados $w_i f_{ij} J_i(\varphi_{ij}^*)$. Sin embargo, como se ha descrito anteriormente, cualquier empresa que entra en el país $i$ debe incurrir en el coste fijo de entrada $w_i f_i^e$ antes de observar su productividad. En conjunto, la condición de libre entrada (FE) se expresa como:

$$\sum_{n=i,j} f_{in} J_i(\varphi_{in}^*) = f_i^e. \tag{2}$$

La condición de libre entrada determina las demandas de mercado $B_i$, $B_j$ ajustando los índices de precios $P_i$, $P_j$ en equilibrio, de modo que los entrantes potenciales obtengan beneficios esperados nulos.

El trabajo se utiliza tanto en la entrada como en la producción, y debe igualar la oferta agregada de trabajo en la economía. A partir de (1) y (2), la condición de \textbf{equilibrio en el mercado de trabajo} (LMC) se expresa como:

$$\frac{R_i - T_i}{w_i} = L_i$$,

Donde $T_i \equiv (\tau_{ji} - 1) R_{ji}$ son los ingresos arancelarios agregados y $R_{ji}$ es el gasto agregado en bienes del país $j$ en el país $i$.\footnote{Como se ha definido $r_{ji}(\varphi)$ libre de aranceles, $R_{ji}$ también queda definido libre de aranceles, agregando $r_{ji}(\varphi)$ sobre la productividad $\varphi\in(\varphi_{ji}^*,\varphi_{\max})$ sobre una masa $M_i^e$ de entrantes. Por tanto, $R_i$ es el consumo agregado libre de aranceles.
}

La condición LMC determina los salarios $w_i$, $w_j$, lo cual se entiende mejor reescribiéndola como:

$$R_i = w_i L_i + T_i$$

Es decir, el salario en el país $i$ se determina por la igualdad entre el gasto agregado $R_i$ y la suma de la renta laboral $w_i L_i$ más los ingresos arancelarios $T_i$. No existe ningún excedente neto adicional más allá de $w_i L_i$ y $T_i$, ya que la libre entrada hace que los beneficios esperados netos sean nulos en equilibrio.

Para trabajar en equilibrio general, relacionamos la condición de equilibrio en el mercado de trabajo (LMC) con la condición de equilibrio de la balanza comercial (TB). Mientras que la condición TB requiere que $R_{ij} = R_{ji}$, es equivalente a la condición LMC, en el sentido de que ambas implican la misma igualdad, $R_i = w_i L_i + T_i$. 

Obsérvese que el gasto agregado en el país $i$ consiste en el gasto en bienes domésticos en $i$ y en bienes extranjeros procedentes del país $j$ incluyendo aranceles $\left(R_i = \sum_n \tau_{ni} R_{ni}\right)$, ya que los consumidores acceden a los bienes extranjeros con aranceles incluidos. Por otro lado, la renta laboral agregada en el país $i$ consiste en los ingresos de las empresas domésticas y de las empresas exportadoras en $i$, netos de aranceles $\left(w_i L_i = \sum_n R_{in}\right)$, ya que las empresas perciben los ingresos exteriores netos de aranceles.

Sea $\lambda_{ji} \equiv \tau_{ji} R_{ji} / \sum_n \tau_{ni} R_{ni}$ la cuota de comercio exterior en el país $i$ incluyendo aranceles. De forma similar, sea $\tilde{\lambda}_{ji} \equiv R_{ji} / \sum_n R_{ni}$ la cuota de comercio correspondiente neta de aranceles. Entonces, podemos expresar el gasto agregado en bienes domésticos y extranjeros mediante las cuotas de gasto netas de aranceles: $R_{ii} = \tilde{\lambda}_{ii} w_i L_i, \quad R_{ji} = \tilde{\lambda}_{ji} w_i L_i$.\footnote{El resultado sigue de observar: $w_iL_i=\sum_nR_{in}=\sum_nR_{ni}$ (de $R_{ij}=R_{ji}$) en el denominados de $\tilde\lambda_{ji}$
}

Sustituyendo estas expresiones en la condición LMC, obtenemos una forma familiar de la condición de equilibrio de la balanza comercial en presencia de aranceles:

$$w_i L_i = \sum_{n=i,j} \tilde{\lambda}_{in} w_n L_n. \tag{3}$$

Estamos ahora en condiciones de caracterizar las variables de equilibrio cuando los países tienen la posibilidad de fijar aranceles. Dados los valores de las variables exógenas, el equilibrio en niveles se define como el conjunto del vector $\{\varphi_{ii}^*, \varphi_{ij}^*, B_i, w_i\}$, que viene determinado conjuntamente por el sistema de ocho ecuaciones en (1), (2) y (3) para $i$, $j$. 

Por la ley de Walras, el salario en el país $j$ puede normalizarse a la unidad, es decir, $w_j = 1$. Una vez determinados los niveles de estas variables clave, el resto de variables endógenas puede expresarse como función de ellas. En particular, utilizando la definición de $B_i$ en (1), el bienestar por trabajador, equivalente al salario real, se expresa como sigue (véase el Apéndice A.3):

$$W_i =\left(\frac{L_i}{\sigma f_{ii}}\right)^{\frac{1}{\sigma-1}}(\mu_i)^{\frac{1}{\rho}} \,\varphi_{ii}^*, \tag{4}$$

Donde $\mu_i \equiv R_i/(w_i L_i)$ se denomina \textbf{multiplicador arancelario} (FELBERMAYR et al., 2015) en el análisis posterior. Esta variable aparece en la expresión de bienestar en (4), ya que los ingresos arancelarios se redistribuyen a los consumidores. Esto implica que el multiplicador arancelario $\mu_i$, así como la cuota de comercio doméstico $\tilde{\lambda}_{ji}$, están estrechamente relacionados con los aranceles $\tau_{ji}$.

A partir de las definiciones de estas variables y de la relación $\mu_i = 1 + (\tau_{ji} - 1)\tilde{\lambda}_{ji}$, derivada de $R_i = w_i L_i + T_i$, obtenemos:

$$\tilde{\lambda}_{ji} = \frac{\lambda_{ji}}{\tau_{ji}(1 - \lambda_{ji}) + \lambda_{ji}}, \qquad \mu_i = \frac{\tau_{ji}}{\tau_{ji}(1 - \lambda_{ji}) + \lambda_{ji}}$$.

Es evidente que $\tilde{\lambda}_{ji} = \lambda_{ji}$ y $\mu_i = 1$ cuando el país $i$ no impone aranceles sobre los bienes extranjeros del país $j$ ($\tau_{ji} = 1$).

$$\text{\textbf{3} Impacto de la competencia}$$

La sección anterior ha definido las variables de equilibrio en niveles. Esta sección define las variables de equilibrio en cambios, con el fin de analizar los efectos de la liberalización comercial y del tamaño de mercado. Aunque nos centraremos principalmente en los costes comerciales variables en esta sección, los efectos de los costes comerciales fijos y de los aranceles son muy similares. Sin embargo, a diferencia de los costes comerciales variables y fijos, que consumen recursos reales, los aranceles generan ingresos públicos y se utilizan para manipular los términos de intercambio. En la Sección 4 caracterizaremos los aranceles óptimos que maximizan el bienestar teniendo en cuenta este motivo.

$$\text{\textbf{3.1} Equilibrio en variaciones}$$

Supongamos que el país $i$ reduce unilateralmente los costes comerciales variables $\theta_{ji}$ y expande el tamaño de mercado $L_i$, manteniendo constantes el resto de variables exógenas. Nos interesa el impacto sobre la competencia inducido por estos cambios, que potencialmente puede diferir entre ambos. Denotemos con un \textit{sombrero} los cambios proporcionales de una variable $(\hat{x} \equiv dx/x)$.

Tomando logaritmos y diferenciando totalmente la condición de beneficio nulo en el umbral (ZCP) (1), obtenemos:
$$\hat{B}_i + (\sigma - 1)\,\hat{\varphi}_{ji}^* = \sigma \,\hat{w}_j + (\sigma - 1)\,\hat{\tau}_{ji}. \tag{5}$$

Aunque permitimos cambios tanto en $\theta_{ji}$ como en $L_i$, la ecuación (5) solo incluye cambios en $\theta_{ji}$, ya que (1) no depende directamente de $L_i$. Por tanto, (5) muestra que las variaciones en el umbral de productividad de exportación $\hat{\varphi}_{ji}^*$ están vinculadas a cambios en el resto de variables endógenas y en los costes comerciales variables. De forma análoga a (1), las variaciones en el umbral de productividad doméstico $\hat{\varphi}_{jj}^*$ también vienen dadas por (5) cuando $i = j$, teniendo en cuenta que $\hat{\theta}_{jj} = 0$.

De manera similar, diferenciando totalmente la condición de libre entrada (2) y reordenando, se obtiene:

$$\hat{\varphi}_{ij}^* = -\alpha_i \,\hat{\varphi}_{ii}^*, \tag{6}\qquad \text{donde},\,\alpha_i \equiv\frac{f_{ii} J_i'(\varphi_{ii}^*) \varphi_{ii}^*}{f_{ij} J_i'(\varphi_{ij}^*) \varphi_{ij}^*}.$$

Dado que (2) no depende directamente de $\theta_{ji}$ ni de $L_i$, los cambios en estas variables no aparecen en (6). Observando $\alpha_i$, tenemos que $\alpha_i > 0$, ya que $J_i(\varphi^*)$ es estrictamente decreciente en $\varphi^*$. Así (por (6)), los cambios en el umbral de productividad de exportación $\hat{\varphi}_{ij}^*$ se producen en dirección opuesta a los cambios en el umbral de productividad doméstico $\hat{\varphi}_{ii}^*$. Además, utilizando el hecho de que $J_i(\varphi^*)$ está relacionado con los beneficios esperados y que la condición de equilibrio de la balanza comercial exige $R_{ij} = R_{ji}$, puede demostrarse que $\alpha_i$ es igual a $R_{ii}/R_{ji}$, es decir, la razón entre el gasto agregado en bienes domésticos y extranjeros en el país $i$.

Esto permite expresar las cuotas de comercio exterior y el multiplicador arancelario en función de $\alpha_i$:

$$\lambda_{ji} = \frac{\tau_{ji}}{\alpha_i + \tau_{ji}},\qquad\tilde{\lambda}_{ji} = \frac{1}{\alpha_i + 1}, \qquad \mu_i = \frac{\alpha_i + \tau_{ji}}{\alpha_i + 1}.$$

Finalmente, utilizando $\tilde{\lambda}_{ji}$ definida anteriormente, podemos reescribir la condición de equilibrio de la balanza comercial (3) como

$$\frac{w_i L_i}{\alpha_i + 1} = \frac{w_j L_j}{\alpha_j + 1}$$,

Donde $\alpha_i$ es función de $\varphi_{ii}^*$, $\varphi_{ij}^*$. Tomando logaritmos y diferenciando totalmente esta igualdad, y utilizando (6), obtenemos:

$$\begin{aligned}
    &\hat{w}_i - \hat{w}_j = -\beta_i \,\hat{\varphi}_{ii}^* + \beta_j \,\hat{\varphi}_{jj}^* - \hat{L}_i,\quad (7)\quad \text{donde},\,\beta_i \equiv\frac{\alpha_i}{\alpha_i + 1}\left[(\sigma - 1) + \gamma_{ii} + (\sigma - 1 + \gamma_{ij}) \alpha_i\right],\\[2em]
    &\text{y},\qquad \gamma_{in} \equiv - \frac{\mathrm d \ln \int_{\varphi_{in}^*}^{\varphi_{\max}} \varphi^{\sigma-1} \mathrm dC_i(\varphi)}{\mathrm d \ln \varphi_{in}^*}
\end{aligned}$$

Siendo $\gamma_{in}$ la elasticidad del margen extensivo que surge en presencia de heterogeneidad de empresas. Se defina el margen extensivo con signo negativo para que $\gamma_{in}>0$ (como la elasticidad del margen intensivo $\sigma-1>0$).

La ecuación (7) muestra que los cambios en los salarios $\hat{w}_i$, $\hat{w}_j$ están vinculados a cambios en los umbrales de productividad domésticos $\hat{\varphi}_{ii}^*$, $\hat{\varphi}_{jj}^*$, donde únicamente los cambios en el tamaño de mercado intervienen directamente. Observando $\beta_i (> 0)$, se identifican dos canales a través de los cuales operan estos cambios:
\begin{la}
    \item El primero es el margen intensivo: los shocks exógenos inducen a las empresas incumbentes a ajustar su producción mediante cambios en la demanda de los consumidores, con una elasticidad $\sigma - 1$ bajo preferencias CES; y,
    \item El segundo es el margen extensivo: los shocks exógenos inducen la entrada de nuevas empresas o la salida de empresas existentes mediante cambios en la competitividad de cada mercado, con una elasticidad $\gamma_{in}$ en presencia de heterogeneidad de empresas.
\end{la}

Ahora estamos en condiciones de caracterizar las variables de equilibrio en variaciones. Del mismo modo que (1), (2) y (3) se utilizan para resolver el equilibrio en niveles, (5), (6) y (7) se utilizan para resolver el equilibrio en cambios. En concreto, el equilibrio en cambios se define como el conjunto del vector $\{\hat{\varphi}_{ii}^*, \hat{\varphi}_{ij}^*, \hat{B}_i, \hat{w}_i\}$, caracterizado conjuntamente por el sistema de ocho ecuaciones en (5), (6) y (7) para $i$, $j$. Por la ley de Walras, los cambios en el salario del país $j$ pueden normalizarse a cero, es decir, $\hat{w}_j = 0$. Una vez determinados los cambios en estas variables clave, los cambios en el resto de variables endógenas pueden escribirse como función de ellas. En particular, tomando logaritmos y diferenciando totalmente (4), los cambios en el bienestar por trabajador se expresan como sigue:

$$\hat{W}_i =\left(\frac{(\tau_{ji}-1)\lambda_{ii}}{\rho}\frac{\beta_i}{\alpha_i}+1\right)\hat{\varphi}_{ii}^*+ \frac{\hat{L}_i}{\sigma-1}.\qquad (8)$$

La ecuación (8) implica que \textbf{el umbral de productividad doméstico $\varphi_{ii}^*$ es una estadística suficiente para el bienestar incluso en presencia de ingresos arancelarios}: cualquier shock exógeno genera cambios en el bienestar a través de cambios en el umbral $\varphi_{ii}^*$, que están directamente relacionados con reasignaciones de recursos al estilo de MELITZ (2003). En presencia de aranceles, dichos shocks también inducen cambios en el bienestar a través de cambios en el multiplicador arancelario $\mu_i$, que recoge la devolución de los ingresos arancelarios a los consumidores. 

Sin embargo, el primer término de (8) muestra que los cambios en bienestar asociados a estos efectos quedan resumidos por los cambios en $\varphi_{ii}^*$. Este impacto se aplica no solo a los costes comerciales variables, sino también al tamaño de mercado, ya que el umbral se ve afectado endógenamente a través de cambios en los salarios (véase (1)). Además, los cambios en el tamaño de mercado tienen un canal adicional sobre el bienestar a través de cambios en la variedad de productos, recogido en el segundo término de (8).

$$\text{\textbf{3.2} Costes comerciales}$$

En la Sección 3.1 hemos caracterizado el equilibrio en cambios teniendo en cuenta simultáneamente los costes comerciales y el tamaño de mercado. Aunque esto permite estudiar ambos efectos utilizando un único sistema de ecuaciones, la mayoría de modelos los analizan por separado. Demidova y Rodríguez-Clare (2013), por ejemplo, consideran el efecto sobre el bienestar de una liberalización comercial asimétrica, manteniendo constante el tamaño de mercado. Muestran que, cuando los salarios se determinan endógenamente por la condición de equilibrio de la balanza comercial, una liberalización comercial unilateral aumenta el bienestar en el país que liberaliza. Este resultado implica que el efecto de la liberalización comercial se invierte, ya que se sabe que tal liberalización reduce el bienestar en el país liberalizador cuando los salarios están fijados exógenamente por la existencia de un bien exterior (Demidova, 2008; Melitz y Ottaviano, 2008). Aquí proporcionamos soluciones analíticas de los resultados de Demidova y Rodríguez-Clare (2013).

Podemos aislar el impacto de una liberalización comercial unilateral sobre el equilibrio manteniendo constante el tamaño de mercado. Recordemos que el modelo consta de un sistema de ocho ecuaciones ((5), (6), (7)) para ocho incógnitas $(\hat{\varphi}_{ii}^*, \hat{\varphi}_{ij}^*, \hat{B}_i, \hat{w}_i \text{ para } i,j)$, donde el trabajo en el país $j$ se toma como numerario. Resolviendo este sistema imponiendo $\hat{L}_i = 0$, obtenemos:

$$\begin{aligned}
    &\hat{\varphi}_{ii}^* =-\frac{\rho(\beta_j+\rho)}{\Xi}\,\hat{\theta}_{ji},\qquad\hat{\varphi}_{jj}^* =-\frac{\rho(\beta_i-\rho \alpha_i)}{\Xi}\,\hat{\theta}_{ji},\qquad\hat{w}_i =\frac{\rho^2(\beta_i+\alpha_i\beta_j)}{\Xi}\,\hat{\theta}_{ji}, \quad (9)\\[2em]
    &\text{donde}\qquad \beta_i - \rho \alpha_i > 0\qquad  (\text{por definiciones de }\,\,\,  \alpha_i\,\,\,  \text{y}\,\,\,  \beta_i)\\[1em]
    &\text{y}\qquad  \Xi \equiv \prod_n (\beta_n+\rho) - \prod_n (\beta_n-\rho\alpha_n) > 0.
\end{aligned}$$

La ecuación (9) muestra que una reducción en $\tau_{ji}$ eleva $\varphi_{ii}^*$ y $\varphi_{jj}^*$, y reduce $w_i$. A partir de (8), se sigue entonces que el bienestar aumenta no solo en el país $j$, sino también en el país $i$, porque la caída en $w_i$ es menor que la caída en $P_i$ (es decir, los salarios reales $w_i/P_i$ aumentan). Además, los ingresos arancelarios redistribuidos a los consumidores aumentan a través del incremento en $\mu_i$, lo que contribuye adicionalmente al bienestar.

La intuición de estos resultados se aprecia con claridad resolviendo primero (5) y (6), sin utilizar todavía (7):

$$\begin{aligned}
    \hat{\varphi}_{ii}^*=\frac{1}{\alpha_i\alpha_j-1}\,\hat{\theta}_{ji}-\frac{\alpha_j+1}{\rho(\alpha_i\alpha_j-1)}\,\hat{w}_i,\\[2em]
    \hat{\varphi}_{jj}^*=-\frac{\alpha_j}{\alpha_i\alpha_j-1}\,\hat{\theta}_{ji}+\frac{\alpha_i+1}{\rho(\alpha_i\alpha_j-1)}\,\hat{w}_i, 
\end{aligned}\qquad (10)\quad \text{donde},\quad  \alpha_i\alpha_j-1>0
$$

En (10), el primer término es el efecto directo de los costes comerciales variables y el segundo término es el efecto indirecto a través de los cambios en salarios. Para ser preciso, los cambios en salarios son cambios en salarios relativos a $i$, ya que los salarios en $j$ están normalizados a $1$.

El efecto directo reduce los beneficios esperados en el país que liberaliza al disminuir los markups sobre las importaciones, pero incrementa los beneficios esperados en el país que no liberaliza al permitir que sus empresas exporten con mayor facilidad. Como resultado, una reducción de $\theta_{ji}$ desincentiva la entrada en el país $i$ e incentiva la entrada en el país $j$, reduciendo $\varphi_{ii}^*$ y aumentando $\varphi_{jj}^*$. Nótese que este efecto existe incluso cuando los salarios están fijados exógenamente por un bien exterior. A partir de (8), estos cambios implican que una liberalización comercial unilateral reduce el bienestar en el país que liberaliza, pero lo incrementa en el país que no liberaliza (Demidova, 2008; Melitz y Ottaviano, 2008).

Si los salarios son endógenos, una reducción de $\theta_{ji}$ (costes de importación) provoca un aumento de las importaciones en el país $i$, lo que es contrarrestado por una caída de los salarios:
\begin{la}
    \item El efecto indirecto aumenta los beneficios esperados en el país que liberaliza ($i$) al reducir los costes de producción de las empresas ($\uparrow \varphi_{ii}^*$), pero reduce los beneficios esperados en el país que no liberaliza al elevar allí los costes de producción en términos relativos frente al país liberalizador ($\varphi_{jj}^*$). En consecuencia, una reducción de $w_i$ (causada por una reducción de $\theta_{ji}$) incentiva la entrada en el país $i$, pero desincentiva la entrada en el país $j$, elevando $\varphi_{ii}^*$ y reduciendo $\varphi_{jj}^*$;
    \item El efecto indirecto opera en direcciones opuestas al efecto directo, pero (9) muestra que tanto $\varphi_{ii}^*$ como $\varphi_{jj}^*$ aumentan como resultado de una reducción de $\tau_{ji}$.
\end{la}

Este resultado de equilibrio solo se cumple cuando el efecto indirecto supera al efecto directo en el caso de $\varphi_{ii}^*$, mientras que ocurre lo contrario en el caso de $\varphi_{jj}^*$.

Dado que el umbral de productividad doméstico aumenta, la ecuación (8) sugiere que una liberalización comercial unilateral incrementa el bienestar en ambos países. Además, puesto que el umbral de productividad de exportación disminuye (véase (6)), dicha liberalización permite que un mayor número de empresas exporte y reduce la cuota de comercio doméstico.

Aunque nos hemos centrado en el impacto de los costes comerciales variables de importación $\theta_{ji}$, el efecto de cualquier coste comercial $(\theta_{ij}, \theta_{ji}, f_{ij}, f_{ji}, \tau_{ij}, \tau_{ji})$ sobre los umbrales de productividad es cualitativamente similar. En el caso de los costes comerciales variables de exportación $\theta_{ij}$, por ejemplo, obtenemos:

$$\begin{aligned}
    &\hat{\varphi}_{ii}^*=-\frac{\rho(\beta_j-\rho\alpha_j)}{\Xi}\,\hat{\theta}_{ij},\\[1em]
    &\hat{\varphi}_{jj}^*=-\frac{\rho(\beta_i+\rho)}{\Xi}\,\hat{\theta}_{ij},\\[1em]
    &\hat{w}_i=-\frac{\rho^2(\beta_j+\alpha_j\beta_i)}{\Xi}\,\hat{\theta}_{ij}
\end{aligned}$$.

Por tanto, una reducción de los costes de exportación $\theta_{ij}$ también incrementa el umbral de productividad doméstico en ambos países, como antes. La única diferencia es que una reducción de los costes de importación $\theta_{ji}$ reduce $w_i$, mientras que una reducción de los costes de exportación $\theta_{ij}$ eleva $w_i$. La diferencia en los cambios salariales refleja el hecho de que una reducción de $\theta_{ij}$ conduce a un aumento de las exportaciones en el país $i$, lo que debe ser contrarrestado por una subida de los salarios en ese país para restablecer el equilibrio comercial.

Por último, partiendo de una situación simétrica, las ganancias de bienestar derivadas de una liberalización comercial unilateral son siempre mayores en el país que liberaliza que en el país que no lo hace. Consideremos el efecto de los costes comerciales variables de importación $\theta_{ji}$. Evaluar (9) en $\alpha_i=\alpha_j$ y $\beta_i=\beta_j$ revela que $|\hat{\varphi}_{ii}^*| > |\hat{\varphi}_{jj}^*|$, lo que implica que $\hat{W}_i > \hat{W}_j$ a partir de (8) imponiendo $\hat{L}_i=0$. Así, una reducción de $\theta_{ji}$ genera mayores ganancias de bienestar en el país $i$ que en el país $j$. Esto también se cumple para los costes comerciales variables de exportación $\theta_{ij}$, en el sentido de que, partiendo de una situación simétrica, una reducción de $\theta_{ij}$ genera mayores ganancias de bienestar en el país $j$ que en el país $i$.

\textbf{Lema 1.} La liberalización comercial unilateral tiene los siguientes efectos sobre las variables de equilibrio:
\begin{lnum}
    \item El salario disminuye en el país que liberaliza;
    \item El umbral de productividad doméstico aumenta tanto en el país que liberaliza como en el que no lo hace. Como resultado, la cuota de comercio doméstico disminuye en ambos países;
    \item La liberalización comercial unilateral aumenta siempre el bienestar en ambos países. Partiendo de una situación simétrica, el efecto sobre el bienestar es siempre mayor en el país que liberaliza que en el que no lo hace.
\end{lnum}

El Lema 1 es esencialmente equivalente al resultado principal de DEMIDOVA y RODRÍGUEZ-CLARE (2013). Estos autores muestran, sin recurrir a una distribución específica de productividad ni a la existencia de un bien exterior, que \textbf{los salarios endógenos invierten el efecto de la liberalización comercial asimétrica sobre el bienestar en el país que liberaliza}, debido a la \textbf{ausencia del efecto mercado interno} en los patrones de comercio. Una de las diferencias clave es que ellos presentan su resultado gráficamente mediante una figura sencilla, mientras que nosotros lo demostramos analíticamente mediante álgebra en tasas de variación (“hat algebra”). Más importante aún es la tractabilidad de nuestro enfoque para analizar el impacto de otra medida de competencia, el tamaño de mercado, que se estudia de forma paralela a la liberalización comercial sin recurrir a una distribución específica de productividad ni a la existencia de un bien exterior (Sección 3.3). Además, nuestras soluciones analíticas de estática comparada permiten abordar tanto los efectos cualitativos como cuantitativos de cambios unilaterales en estas presiones competitivas sobre la política comercial óptima (Secciones 4 y 5).

$$\text{\textbf{3.3} Tamaño de mercado}$$

A continuación, analizamos el impacto del tamaño de mercado, manteniendo constantes el resto de variables exógenas, cuestión que también ha sido ampliamente estudiada en la literatura. MELITZ y OTTAVIANO (2008) fueron los primeros en mostrar que un país de mayor tamaño presenta mayor productividad y mayor bienestar debido a una competencia más intensa en el mercado doméstico. Debido a la presencia de un bien exterior que implica salarios fijos, encuentran que \textbf{la liberalización comercial y el tamaño de mercado tienen efectos opuestos sobre el bienestar: un país que liberaliza unilateralmente se ve perjudicado debido a la relocalización de la entrada desde el país liberalizador hacia el país no liberalizador}, como en la Sección 3.2. En esta sección mostramos que, \textbf{en ausencia de un bien exterior, el efecto del tamaño de mercado se invierte: un país que se expande unilateralmente se ve perjudicado debido a la relocalización de la entrada desde el país en expansión hacia el país que no se expande}. Este resultado implica que las \textbf{empresas ineficientes sobreviven en el mercado doméstico de gran tamaño, lo que denominamos selección doméstica débil}. 

Aunque este canal a través de la selección afecta negativamente a la productividad y al bienestar, un país grande puede, no obstante, beneficiarse de su tamaño de mercado, ya que el efecto negativo de la selección débil puede verse compensado por el efecto positivo derivado del aumento de la variedad de productos.

Al igual que en la Sección 3.2, podemos aislar el impacto de una expansión unilateral del mercado sobre el equilibrio manteniendo constantes los costes comerciales variables. Resolviendo (5), (6) y (7) imponiendo $\hat{\theta}_{ji} = 0$ para las ocho incógnitas $(\hat{\varphi}_{ii}^*, \hat{\varphi}_{ij}^*, \hat{B}_i, \hat{w}_i$ para $i,j)$, obtenemos las siguientes relaciones de equilibrio:

$$\begin{aligned}
    \hat{\varphi}_{ii}^* =-\frac{\rho(\alpha_j+1)}{\Xi}\,\hat{L}_i,\qquad\hat{\varphi}_{jj}^* =\frac{\rho(\alpha_i+1)}{\Xi}\,\hat{L}_i,\qquad\hat{w}_i =\frac{\rho^2(\alpha_i\alpha_j-1)}{\Xi}\,\hat{L}_i
\end{aligned}$$.

Por tanto, una expansión en $L_i$ reduce $\varphi_{ii}^*$, pero aumenta $\varphi_{jj}^*$ y $w_i$. El efecto del tamaño de mercado es distinto del de la liberalización comercial en (9), aunque ambos shocks inducen una mayor presión competitiva en el mercado del país $i$. A partir de (8), se deduce entonces que el bienestar siempre aumenta en el país $j$, pero puede aumentar o disminuir en el país $i$ dependiendo de hasta qué punto la expansión del tamaño de mercado reduce su umbral de productividad doméstico.

La intuición se entiende claramente resolviendo primero (5) y (6) sin (7):

$$\begin{aligned}
    \hat{\varphi}_{ii}^* =-\frac{\alpha_j + 1}{\rho(\alpha_i\alpha_j - 1)}\,\hat{w}_i\\[2em]
    \hat{\varphi}_{jj}^* =\frac{\alpha_i + 1}{\sigma(\alpha_i\alpha_j - 1)}\,\hat{w}_i.
\end{aligned}\qquad (12)$$

La comparación entre (10) y (12) revela inmediatamente que \textbf{el efecto directo del tamaño de mercado está ausente} en una expansión unilateral de mercado, debido a la propiedad particular y restrictiva de las preferencias CES y la competencia monopolística, y que \textbf{solo existe el efecto indirecto a través de cambios en los salarios}. Así, si los salarios son exógenos (determinados por un bien exterior), el tamaño de mercado no tiene ningún efecto sobre el umbral de productividad doméstico. Dado que la expansión del tamaño de mercado no induce directamente entrada ni salida, (8) implica que \textbf{una expansión unilateral del mercado aumenta el bienestar en el país que se expande únicamente a través del aumento de la variedad de productos}, lo cual es el resultado estándar tanto en modelos con heterogeneidad de empresas (Marc J. Melitz, 2003) como en modelos con empresas homogéneas (Paul Krugman, 1980).

Si los salarios son endógenos, un aumento en $L_i$ conduce a mayores salarios en el país $i$ cuando las empresas incurren en costes comerciales. El efecto indirecto reduce los beneficios esperados en el país que se expande al aumentar los costes de producción, mientras que ocurre lo contrario en el país que no se expande. Por tanto, un aumento en $w_i$ (debido al aumento en $L_i$) desincentiva la entrada en el país $i$ pero la incentiva en el país $j$. Como resultado, la cuota de comercio doméstico aumenta en el país $i$ pero disminuye en el país $j$. Obsérvese que el efecto negativo sobre $\varphi_{ii}^*$ proviene del efecto de mercado interno sobre los salarios de KRUGMAN (1980), aunque dicho efecto negativo no aparece en su modelo donde la productividad es exógena. En nuestro modelo, sin embargo, mayores salarios implican mayores costes marginales y menor rentabilidad en el país que se expande. \textbf{Esto relocaliza la entrada desde el país en expansión (con salarios altos) hacia el país que no se expande (con salarios más bajos), lo que permite que empresas ineficientes sobrevivan en el mercado doméstico.} Esta intuición explica por qué la expansión del tamaño de mercado en un país afecta simultáneamente a la productividad en otro país, algo que no ocurre cuando los salarios son exógenos (Gianmarco Ottaviano y Melitz, 2008).\footnote{Al igual que la liberalización comercial unilateral, la expansión unilateral del mercado no genera el efecto de mercado interno sobre los patrones de comercio cuando los salarios son endógenos. En la versión working paper de este artículo (Ara, 2021), mostramos que la expansión del mercado en el país $i$ modifica los patrones de comercio a favor del país $j$ tanto a través del margen intensivo como del margen extensivo. Esto es similar al resultado de Bertoletti y Etro (2017), quienes muestran que la expansión del mercado conduce a un cambio en los patrones de comercio que puede interpretarse como destrucción empresarial: un país en expansión con salarios altos se caracteriza por una mayor concentración de empresas exportadoras.
}

¿Qué ocurre con el efecto de $L_i$ sobre el bienestar en el país que se expande? La inspección de (8) muestra que depende de la magnitud de la reducción en $\varphi_{ii}^*$ y del aumento en $L_i$, que capturan respectivamente los efectos de la selección doméstica débil y del aumento de la variedad de productos, los cuales afectan al bienestar en direcciones opuestas. Para ver cuál domina en equilibrio, expresamos los cambios en el bienestar únicamente en función de los cambios en $\varphi_{ii}^*$:

$$\hat{W}_i =\frac{1}{\sigma - 1}\left[(\sigma - 1)(\beta_i + \rho) -\frac{\sigma\beta_i}{\mu_i}-(\beta_j-\rho\alpha_j)\left(\frac{\alpha_i+1}{\alpha_j+1}\right)\right]\hat{\varphi}_{ii}^*.(13)$$

Dado que un aumento en $L_i$ reduce $\varphi_{ii}^*$, (13) implica que el país $i$ se beneficia de la expansión del mercado si el término entre corchetes es negativo. Desafortunadamente, esto no siempre se cumple, por lo que no puede afirmarse de forma general que la expansión unilateral del mercado genere ganancias de bienestar en este modelo. Sin embargo, puede demostrarse que, partiendo de una situación simétrica $(\alpha_i = \alpha_j, \beta_i = \beta_j)$ y de libre comercio ($\mu_i = 1$), dicha expansión mejora inequívocamente el bienestar tanto en el país que se expande como en el que no.

\textbf{Lema 2.} La expansión unilateral del tamaño de mercado tiene los siguientes efectos sobre las variables de equilibrio:
\begin{lnum}
    \item El salario aumenta en el país que se expande;
    \item El umbral de productividad doméstico disminuye en el país que se expande, pero aumenta en el país que no se expande. Como resultado, la cuota de comercio doméstico aumenta en el país que se expande y disminuye en el país que no se expande;
    \item Partiendo de una situación simétrica y de libre comercio, la expansión unilateral del mercado aumenta el bienestar tanto en el país que se expande como en el que no lo hace.
\end{lnum}

El resultado del Lema 2 presenta una diferencia notable con respecto a estudios previos: la expansión del mercado no siempre conduce a mejoras de bienestar debido a las distorsiones derivadas de esa selección débil.\footnote{En su influyente trabajo sobre eficiencia asignativa bajo preferencias generales, Swati Dhingra y John Morrow (2019) encuentran que la expansión del mercado genera ganancias de bienestar cuando las preferencias de los consumidores están alineadas, es decir, cuando los desplazamientos de la demanda afectan a los márgenes privados y sociales en la misma dirección. Su resultado implica que la expansión del mercado mejora el bienestar bajo preferencias CES (ya que los márgenes no dependen del tamaño de mercado), pero esto no se cumple en el presente modelo. Como muestran Dhingra y Morrow (2019), una condición suficiente para obtener ganancias de bienestar es que la expansión del mercado no tenga un impacto negativo sobre la productividad. Esta condición no se cumple aquí, porque la expansión del mercado genera selección doméstica débil, lo que reduce la productividad en el país que se expande.

Gabriel Felbermayr y Jung (2018) también muestran que un país más grande tiende a presentar una selección doméstica más débil; sin embargo, su análisis del tamaño de mercado se limita a una distribución de Pareto.
}


$$\text{\textbf{4} Política Comercial}$$

Hasta ahora hemos analizado el impacto de cambios exógenos en las medidas de competencia sobre las principales variables endógenas sin especificar una distribución concreta de productividad ni recurrir a un bien exterior. En esta sección mostramos que este grado de generalidad es clave para caracterizar la política comercial óptima de un país.

$$\text{\textbf{4.1} Aranceles óptimos}$$

Supongamos que el gobierno fija el arancel con el objetivo de maximizar el bienestar. A continuación nos centramos en caracterizar el arancel óptimo del país $i$, es decir, el arancel que este país impondría sin temor a represalias por parte del país $j$.

Consideremos el efecto del arancel $\tau_{ji}$ fijado por el país $i$, manteniendo constantes el resto de variables exógenas (incluidos los aranceles $\tau_{ij}$ del país $j$). En el país $j$, el efecto de $\tau_{ji}$ es esencialmente el mismo que el de los costes comerciales variables $\theta_{ji}$, en el sentido de que los cambios en el bienestar vienen determinados únicamente por los cambios en el umbral de productividad doméstico $\varphi_{jj}^*$. A partir del Lema 1, se deduce que un aumento en $\tau_{ji}$ reduce dicho umbral y, por tanto, empeora el bienestar del país $j$.

En el país $i$, en cambio, aparece un efecto adicional del arancel $\tau_{ji}$ sobre el bienestar: un aumento en $\tau_{ji}$ mejora los términos de intercambio del país $i$, lo cual opera a través de cambios en el multiplicador arancelario $\mu_i$. Debido a este nuevo canal, los cambios en el bienestar del país $i$, en analogía con (8), se expresan como:

$$\hat{W}_i =\left(\frac{(\tau_{ji}-1)\lambda_{ii}}{\rho}\cdot\frac{\beta_i}{\alpha_i}+1\right)\hat{\varphi}_{ii}^*+\frac{\lambda{ji}}{\rho}\hat{\tau}_{ji}.$$

El primer término representa una pérdida de bienestar derivada de los aranceles, al proteger a las empresas ineficientes de la competencia extranjera, mientras que el segundo término representa una ganancia de bienestar debida a la mejora en los términos de intercambio. Tras reorganizar, los cambios en el bienestar asociados a estos dos efectos pueden capturarse únicamente mediante los cambios en $\varphi_{ii}^*$:

$$\hat{W}_i =\frac{\lambda_{ji}(\beta_i - \rho\alpha_i)}{\rho}\left[\frac{\beta_j -\rho\alpha_j}{\beta_j+\rho}-\frac{1}{\tau_{ji}}\right]\hat{\varphi}_{ii}^*.(14)$$

Consideremos primero el efecto de los aranceles sobre el bienestar del país $i$ en el entorno del libre comercio, es decir, en $\tau_{ji} = 1$. Recordando el Lema 1, un aumento en $\tau_{ji}$ también reduce $\varphi_{ii}^*$. Evaluando (14) en $\tau_{ji} = 1$, se observa inmediatamente que un pequeño incremento en los aranceles a partir del libre comercio mejora de forma inequívoca el bienestar del país $i$ (a costa del país $j$), y que los aranceles óptimos que maximizan el bienestar son estrictamente positivos para el país $i$. También puede demostrarse que, partiendo de una situación simétrica, la ganancia del país $i$ no compensa la pérdida del país $j$, de modo que el efecto de $\tau_{ji}$ sobre el bienestar mundial es siempre negativo.

Antes de caracterizar los aranceles óptimos del país $i$, resulta útil relacionar la expresión (14) con la de los modelos cuantitativos existentes en la literatura. Utilizando $\lambda_{ii}$ y $\mu_i$ en función de $\alpha_i$ definidos en la Sección 3.1, podemos reescribir (14) como:

$$\begin{aligned}
    \hat{W}_i =- \frac{\alpha_i}{\beta_i}\,\hat{\lambda}_{ii}+\frac{\beta_i-\rho\alpha_i}{\rho\,\beta_i}\,\hat{\mu}_i.\qquad (15)
\end{aligned}$$

Los cambios en el bienestar en (15) engloban los resultados de ARKOLAKIS et al. (2012) sin ingresos arancelarios y los de FELBERMAYR et al. (2015) con ingresos arancelarios para el modelo de MELITZ (2003) bajo una distribución de PARETO, que es una de las más utilizadas en la literatura. Aunque esta hipótesis distribucional proporciona una buena aproximación a la distribución del tamaño de las empresas, conlleva ciertas limitaciones específicas.

En particular, cuando la productividad sigue una distribución de PARETO con parámetro de forma $k$, la elasticidad del margen extensivo es constante e igual a $\gamma_{ii} = \gamma_{ij} = \gamma = k - (\sigma - 1)$, lo que implica que el efecto de los costes comerciales sobre la entrada y salida de empresas es el mismo en los mercados domésticos y extranjeros. Además, sustituyendo $\gamma_{ii}$ y $\gamma_{ij}$ en $\gamma_{ji}$ definido en la Sección 3.1, se obtiene que $\beta_i/\alpha_i = \sigma - 1 + \gamma_{ij}$, es decir, la elasticidad del comercio derivada originalmente por CHANEY (2008) bajo una distribución de PARETO. 

Dado que la elasticidad del margen extensivo es constante, la elasticidad del comercio también lo es en todos los mercados. Denotando esta elasticidad única como $\varepsilon = \sigma - 1 + \gamma$, la expresión (15) puede escribirse como:
$$\begin{aligned}
    \hat{W}_i =- \frac{1}{\varepsilon}\,\hat{\lambda}_{ii}+\left(1+\frac{\eta}{\varepsilon}\right)\,\hat{\mu}_i,\qquad \text{donde},\,\,\, \eta = \frac{k}{\sigma - 1}\left(1 + \frac{1 - \sigma}{k}\right) > 0
\end{aligned}$$.

Esta expresión muestra que los cambios en el bienestar pueden resumirse únicamente mediante dos estadísticas suficientes, $\lambda_{ii}$ y $\varepsilon$, en ausencia de aranceles (primer término), como señalan ARKOLAKIS et al. (2012). Sin embargo, su fórmula de bienestar debe modificarse en presencia de aranceles que generan ingresos públicos (segundo término), como destacan FELBERMAYR et al. (2015).

No obstante, estos resultados dependen críticamente de la hipótesis de que la elasticidad del comercio es única en todos los mercados, como subrayan  MELITZ y REDDING (2015). A partir de la definición de $\gamma_{ij} = - d\ln \int_{\varphi_{ij}^*}^{\varphi_{\max}} \varphi^{\sigma-1} \mathrm dC_i(\varphi) / \mathrm d\ln \varphi_{ij}^*$, la elasticidad del margen extensivo difiere en general entre mercados cuando la distribución de productividad es no restringida, en cuyo caso la elasticidad del comercio pasa a ser específica para cada par de países $i,j$. 

Denotando esta elasticidad variable como $\varepsilon_{ij} = \sigma - 1 + \eta_{ij}$, la expresión (15) puede reescribirse finalmente como:

$$\begin{aligned}
    \hat{W}_i =\frac{1}{\varepsilon_{ij} + \gamma_{ii} - \gamma_{ij}}\left(\hat{M}_i^e - \hat{\lambda}_{ii}\right)+\left(\frac{1}{\rho} -\frac{1}{\varepsilon_{ij} + \gamma_{ii} - \gamma_{ij}}\right)\hat{\mu}_i
\end{aligned}$$.

Esta expresión es análoga a la de MELITZ y REDDING (2015, ecuación (33)), aunque aquí derivamos los cambios en el bienestar en presencia de aranceles que generan ingresos públicos. 

Obsérvese que, además de la cuota de comercio doméstico $\lambda_{ii}$ y de la elasticidad del comercio $\varepsilon_{ij}$, los cambios en el bienestar dependen también del diferencial de elasticidad del margen extensivo entre el mercado doméstico y el de exportación, $\gamma_{ii} - \gamma_{ij}$, que aparece cuando la elasticidad del comercio difiere entre mercados. Estos autores argumentan que, cuando existe dicho diferencial, la cuota de comercio doméstico y la elasticidad del comercio dejan de ser estadísticas suficientes para el bienestar, y que los cambios en bienestar pueden estimarse de forma sustancialmente errónea si se supone que la elasticidad del comercio es constante cuando en realidad es variable. La expresión (16) muestra que esta crítica también se aplica a los cambios en bienestar asociados a los aranceles. Por ejemplo, cuando el margen extensivo es más elástico en el mercado de exportación que en el doméstico $(\gamma_{ii} - \gamma_{ij} < 0)$, los cambios en el bienestar son menores que en ausencia de dicho diferencial $(\gamma_{ii} - \gamma_{ij} = 0)$.

La evidencia empírica reciente documenta que la elasticidad del comercio difiere de forma significativa entre pares de países, respaldando este resultado sobre el bienestar.\footnote{Manteniendo preferencias CES y competencia monopolística —de modo que el margen intensivo es constante—, Elhanan Helpman et al. (2008) encuentran una variación sustancial en la elasticidad del comercio entre pares de países debida al margen extensivo. De manera similar, Bas et al. (2017) muestran que la variación del margen extensivo entre pares de países desempeña un papel clave en la cuantificación de la elasticidad del comercio. Estas evidencias sugieren la existencia de un diferencial en la elasticidad del margen extensivo.
}


Pasamos ahora a caracterizar los aranceles óptimos del país $i$. Como en gran parte de la literatura sobre política comercial, utilizamos la condición de primer orden de maximización del bienestar, suponiendo que la condición de suficiencia se cumple. Igualando $\hat{W}_i = 0$ en (14) y resolviendo para $\tau_{ji}$, se obtiene la siguiente expresión para los aranceles óptimos:

$$\begin{aligned}
    \tau_{ji}^* = 1 +\underbrace{\frac{\rho}{\frac{\alpha_j}{\alpha_j+1}\left(\frac{\beta_j}{\alpha_j}-\rho\right)}}_{t_{ji}^*}=\frac{\beta_j+\rho}{\beta_j+\rho\,\alpha_j}>1
\end{aligned}$$.

Además, teniendo en cuenta que $\tilde{\lambda}_{jj} = \alpha_j / (\alpha_j + 1)$ y reescribiendo la definición de $\beta_j$ en función de la elasticidad del comercio $\varepsilon_{ji}$, se obtiene que los aranceles óptimos pueden expresarse implícitamente como función de momentos observables clave:

$$\begin{aligned}
    t_{ji}^* =\frac{{\rho}}{\tilde{\lambda}_{jj}(\varepsilon_{ji} + (\gamma_{jj} - \gamma_{ji})(1 - \tilde{\lambda}_{jj}) - \rho)}.\qquad (17)
\end{aligned}$$

La ecuación (17) muestra que los aranceles óptimos del país $i$ son \textbf{inversamente proporcionales a la elasticidad de oferta de exportaciones del país $j$}, la cual está compuesta tanto por la cuota de comercio doméstico $\tilde{\lambda}_{jj}$ como por la elasticidad del comercio $\varepsilon_{ji}$, tal como ocurre en modelos tradicionales de comercio. Sin embargo, una diferencia clave en este modelo es que el diferencial de elasticidad del margen extensivo $\gamma_{jj} - \gamma_{ji}$ aparece explícitamente en la expresión de los aranceles óptimos, reflejando el hecho de que la elasticidad del comercio no tiene por qué ser única entre distintos mercados.

La fórmula de arancel óptimo (17) puede interpretarse como una generalización de varios resultados clásicos en la literatura de política comercial:
\begin{la}
    \item \textbf{Empresas heterogenas} (PARETO). Si la productividad sigue una distribución de Pareto con parámetro de forma $k$, entonces la elasticidad del margen extensivo es constante $\gamma_{jj} = \gamma_{ji} = \gamma = k - (\sigma - 1)$, y la elasticidad del comercio también es constante $\varepsilon_{ji} = \varepsilon = \sigma - 1 + \gamma$, como se ha mostrado anteriormente. Dado que en este caso $\varepsilon = k$, la ecuación (17) se reduce a:
    $$\begin{aligned}
    t_{ji}^* =\frac{\tilde{\lambda}_{jj}}{k - \sigma}.\qquad(18)
    \end{aligned}$$
    Esta expresión coincide exactamente con la obtenida por Gabriel Felbermayr et al. (2013) en un modelo de empresas heterogéneas à la Melitz (2003) en el que la productividad sigue una distribución de Pareto. Además, puede considerarse el modelo de empresas homogéneas como un caso particular del modelo de empresas heterogéneas en el que la productividad toma un valor constante o nulo (REDDING y MELITZ, 2015).
    \item \textbf{Empresas homogeneas}. Si los costes comerciales son suficientemente bajos como para que todas las empresas homogéneas exporten en esta clase de modelos, puede demostrarse fácilmente que la elasticidad del margen extensivo es constante e igual a $\gamma_{jj} = \gamma_{ji} = \gamma_ = 0$, y que la elasticidad del comercio también es constante e igual a $\varepsilon_{ji} = \varepsilon = \sigma - 1$. En ese caso, la expresión (17) se reduce a:
    $$\begin{aligned}
    t_{ji}^* = \frac{1}{\tilde{\lambda}_{jj}(\sigma - 1)}\qquad (19)
    \end{aligned}$$.
    Esta expresión coincide exactamente con la obtenida por François Gros (1987) en un modelo de empresas homogéneas à la KRUGMAN (1980).
\end{la}

Desde este punto de vista, la fórmula del arancel óptimo (17) plantea dos advertencias importantes:
\begin{lnum}
    \item En primer lugar, no puede afirmarse en general que los aranceles óptimos sean menores en el modelo de empresas heterogéneas que en el modelo de empresas homogéneas. Del mismo modo que los distintos modelos generan diferentes cuotas de comercio doméstico $\tilde{\lambda}_{jj}$, también generan diferentes elasticidades del comercio $\varepsilon_{ji}$. Esto implica que los aranceles óptimos en distintos modelos no son directamente comparables sin controlar por las diferencias en la elasticidad del comercio. Nuestra fórmula permite clarificar este punto. Sustituyendo en (17) la condición $\gamma_{jj} - \gamma_{ji} = 0$, que se cumple tanto en el modelo heterogéneo con distribución de PARETO como en el modelo homogéneo con distribución degenerada, se obtiene que, condicionando por los dos momentos observables empíricamente (cuota de comercio doméstico y elasticidad del comercio), los aranceles óptimos son idénticos entre modelos;\footnote{Este resultado se obtiene aplicando la fórmula de bienestar de ARKOLAKIS et al. (2012) a nuestra expresión de arancel óptimo: condicionando por las dos estadísticas suficientes para el bienestar, los cambios en bienestar asociados a los aranceles son iguales y, en consecuencia, también lo son los niveles de los aranceles óptimos-
    }
    \item En segundo lugar, la equivalencia de los aranceles óptimos entre modelos solo se mantiene si el diferencial de elasticidad del margen extensivo es nulo $(\gamma_{jj} - \gamma_{ji} = 0)$. Si esta condición no se cumple, los aranceles óptimos difieren incluso controlando por las dos estadísticas suficientes. Consideremos, por ejemplo, el caso en el que el margen extensivo es más elástico en el mercado de exportación que en el doméstico $(\gamma_{jj} - \gamma_{ji} < 0)$. Como se observa en (16), los cambios en bienestar asociados a los aranceles son menores que en ausencia de dicho diferencial. Dado que los aranceles óptimos que maximizan el bienestar son estrictamente positivos, esto implica que el gobierno enfrenta una menor pérdida de bienestar por imponer aranceles y, por tanto, tiene un mayor incentivo a fijar aranceles más altos. En efecto, (17) muestra que los niveles de arancel óptimo son mayores cuando $\gamma_{jj} - \gamma_{ji} < 0$ que cuando $\gamma_{jj} - \gamma_{ji} = 0$, manteniendo constantes las dos estadísticas suficientes. Esto se debe a que el gobierno no internaliza las diferencias en el impacto de los aranceles sobre la entrada y salida de empresas entre mercados. El resultado opuesto se obtiene cuando $\gamma_{jj} - \gamma_{ji} > 0$, en cuyo caso los aranceles óptimos son más bajos que en ausencia de este diferencial. En general, por tanto, la cuota de comercio doméstico y la elasticidad del comercio dejan de ser estadísticas suficientes no solo para el bienestar, como señalan MELITZ y REDDING (2015), sino también para la política comercial óptima.
\end{lnum}

\textbf{Proposición 1.} Condicionando por la cuota de comercio doméstico y la elasticidad del comercio, los niveles de los aranceles óptimos presentan las siguientes propiedades:
\begin{lnum}
    \item Cuando la elasticidad del margen extensivo es la misma en los mercados doméstico y de exportación, los aranceles óptimos son iguales entre distintos modelos de comercio;
    \item Cuando el margen extensivo es más (menos) elástico en el mercado de exportación que en el doméstico, los aranceles óptimos son mayores (menores) que en ausencia de dicho diferencial.
\end{lnum}

En la Proposición 1 se comparan los aranceles óptimos entre distintos modelos manteniendo constantes tanto la cuota de comercio doméstico como la elasticidad del comercio que surgen endógenamente en cada modelo. Si la comparación se realiza sin este condicionamiento, la proposición deja de cumplirse. Los aranceles óptimos en (17), (18) y (19) dependen de la cuota de comercio doméstico, que a su vez depende de los aranceles, por lo que no es constante. El hecho de que los aranceles óptimos estén caracterizados de forma implícita implica que no pueden resolverse en forma cerrada como en trabajos anteriores (Gros, 1987; Felbermayr et al., 2013). Para evitar esta dificultad, Gabriel Felbermayr et al. (2013) comparan los aranceles óptimos entre modelos homogéneos y heterogéneos manteniendo constante únicamente la cuota de comercio doméstico. Más recientemente, Arnaud Costinot et al. (2020) muestran que los aranceles óptimos pueden ser menores bajo una distribución no Pareto (en comparación con una Pareto). Aunque también subrayan la importancia de una distribución general de productividad, comparan los aranceles bajo el mismo tipo de condicionamiento que Felbermayr et al. (2013). Sin embargo, en nuestro caso no es posible adoptar ese enfoque, ya que no solo la cuota de comercio doméstico, sino también la elasticidad del comercio y el diferencial de elasticidad del margen extensivo dependen de los aranceles. Por ello, no es posible determinar qué nivel de arancel óptimo es menor entre (17), (18) y (19) sin fijar previamente las variables de equilibrio de los distintos modelos. Por esta razón, en la Sección 5 recurrimos a soluciones numéricas para evaluar si la variabilidad de la elasticidad del comercio genera diferencias cuantitativamente relevantes en los niveles de los aranceles óptimos.


























\end{document}